\documentclass[11pt]{article}

    \usepackage[breakable]{tcolorbox}
    \usepackage{parskip} % Stop auto-indenting (to mimic markdown behaviour)
    

    % Basic figure setup, for now with no caption control since it's done
    % automatically by Pandoc (which extracts ![](path) syntax from Markdown).
    \usepackage{graphicx}
    % Keep aspect ratio if custom image width or height is specified
    \setkeys{Gin}{keepaspectratio}
    % Maintain compatibility with old templates. Remove in nbconvert 6.0
    \let\Oldincludegraphics\includegraphics
    % Ensure that by default, figures have no caption (until we provide a
    % proper Figure object with a Caption API and a way to capture that
    % in the conversion process - todo).
    \usepackage{caption}
    \DeclareCaptionFormat{nocaption}{}
    \captionsetup{format=nocaption,aboveskip=0pt,belowskip=0pt}

    \usepackage{float}
    \floatplacement{figure}{H} % forces figures to be placed at the correct location
    \usepackage{xcolor} % Allow colors to be defined
    \usepackage{enumerate} % Needed for markdown enumerations to work
    \usepackage{geometry} % Used to adjust the document margins
    \usepackage{amsmath} % Equations
    \usepackage{amssymb} % Equations
    \usepackage{textcomp} % defines textquotesingle
    % Hack from http://tex.stackexchange.com/a/47451/13684:
    \AtBeginDocument{%
        \def\PYZsq{\textquotesingle}% Upright quotes in Pygmentized code
    }
    \usepackage{upquote} % Upright quotes for verbatim code
    \usepackage{eurosym} % defines \euro

    \usepackage{iftex}
    \ifPDFTeX
        \usepackage[T1]{fontenc}
        \IfFileExists{alphabeta.sty}{
              \usepackage{alphabeta}
          }{
              \usepackage[mathletters]{ucs}
              \usepackage[utf8x]{inputenc}
          }
    \else
        \usepackage{fontspec}
        \usepackage{unicode-math}
    \fi

    \usepackage{fancyvrb} % verbatim replacement that allows latex
    \usepackage{grffile} % extends the file name processing of package graphics
                         % to support a larger range
    \makeatletter % fix for old versions of grffile with XeLaTeX
    \@ifpackagelater{grffile}{2019/11/01}
    {
      % Do nothing on new versions
    }
    {
      \def\Gread@@xetex#1{%
        \IfFileExists{"\Gin@base".bb}%
        {\Gread@eps{\Gin@base.bb}}%
        {\Gread@@xetex@aux#1}%
      }
    }
    \makeatother
    \usepackage[Export]{adjustbox} % Used to constrain images to a maximum size
    \adjustboxset{max size={0.9\linewidth}{0.9\paperheight}}

    % The hyperref package gives us a pdf with properly built
    % internal navigation ('pdf bookmarks' for the table of contents,
    % internal cross-reference links, web links for URLs, etc.)
    \usepackage{hyperref}
    % The default LaTeX title has an obnoxious amount of whitespace. By default,
    % titling removes some of it. It also provides customization options.
    \usepackage{titling}
    \usepackage{longtable} % longtable support required by pandoc >1.10
    \usepackage{booktabs}  % table support for pandoc > 1.12.2
    \usepackage{array}     % table support for pandoc >= 2.11.3
    \usepackage{calc}      % table minipage width calculation for pandoc >= 2.11.1
    \usepackage[inline]{enumitem} % IRkernel/repr support (it uses the enumerate* environment)
    \usepackage[normalem]{ulem} % ulem is needed to support strikethroughs (\sout)
                                % normalem makes italics be italics, not underlines
    \usepackage{soul}      % strikethrough (\st) support for pandoc >= 3.0.0
    \usepackage{mathrsfs}
    

    
    % Colors for the hyperref package
    \definecolor{urlcolor}{rgb}{0,.145,.698}
    \definecolor{linkcolor}{rgb}{.71,0.21,0.01}
    \definecolor{citecolor}{rgb}{.12,.54,.11}

    % ANSI colors
    \definecolor{ansi-black}{HTML}{3E424D}
    \definecolor{ansi-black-intense}{HTML}{282C36}
    \definecolor{ansi-red}{HTML}{E75C58}
    \definecolor{ansi-red-intense}{HTML}{B22B31}
    \definecolor{ansi-green}{HTML}{00A250}
    \definecolor{ansi-green-intense}{HTML}{007427}
    \definecolor{ansi-yellow}{HTML}{DDB62B}
    \definecolor{ansi-yellow-intense}{HTML}{B27D12}
    \definecolor{ansi-blue}{HTML}{208FFB}
    \definecolor{ansi-blue-intense}{HTML}{0065CA}
    \definecolor{ansi-magenta}{HTML}{D160C4}
    \definecolor{ansi-magenta-intense}{HTML}{A03196}
    \definecolor{ansi-cyan}{HTML}{60C6C8}
    \definecolor{ansi-cyan-intense}{HTML}{258F8F}
    \definecolor{ansi-white}{HTML}{C5C1B4}
    \definecolor{ansi-white-intense}{HTML}{A1A6B2}
    \definecolor{ansi-default-inverse-fg}{HTML}{FFFFFF}
    \definecolor{ansi-default-inverse-bg}{HTML}{000000}

    % common color for the border for error outputs.
    \definecolor{outerrorbackground}{HTML}{FFDFDF}

    % commands and environments needed by pandoc snippets
    % extracted from the output of `pandoc -s`
    \providecommand{\tightlist}{%
      \setlength{\itemsep}{0pt}\setlength{\parskip}{0pt}}
    \DefineVerbatimEnvironment{Highlighting}{Verbatim}{commandchars=\\\{\}}
    % Add ',fontsize=\small' for more characters per line
    \newenvironment{Shaded}{}{}
    \newcommand{\KeywordTok}[1]{\textcolor[rgb]{0.00,0.44,0.13}{\textbf{{#1}}}}
    \newcommand{\DataTypeTok}[1]{\textcolor[rgb]{0.56,0.13,0.00}{{#1}}}
    \newcommand{\DecValTok}[1]{\textcolor[rgb]{0.25,0.63,0.44}{{#1}}}
    \newcommand{\BaseNTok}[1]{\textcolor[rgb]{0.25,0.63,0.44}{{#1}}}
    \newcommand{\FloatTok}[1]{\textcolor[rgb]{0.25,0.63,0.44}{{#1}}}
    \newcommand{\CharTok}[1]{\textcolor[rgb]{0.25,0.44,0.63}{{#1}}}
    \newcommand{\StringTok}[1]{\textcolor[rgb]{0.25,0.44,0.63}{{#1}}}
    \newcommand{\CommentTok}[1]{\textcolor[rgb]{0.38,0.63,0.69}{\textit{{#1}}}}
    \newcommand{\OtherTok}[1]{\textcolor[rgb]{0.00,0.44,0.13}{{#1}}}
    \newcommand{\AlertTok}[1]{\textcolor[rgb]{1.00,0.00,0.00}{\textbf{{#1}}}}
    \newcommand{\FunctionTok}[1]{\textcolor[rgb]{0.02,0.16,0.49}{{#1}}}
    \newcommand{\RegionMarkerTok}[1]{{#1}}
    \newcommand{\ErrorTok}[1]{\textcolor[rgb]{1.00,0.00,0.00}{\textbf{{#1}}}}
    \newcommand{\NormalTok}[1]{{#1}}

    % Additional commands for more recent versions of Pandoc
    \newcommand{\ConstantTok}[1]{\textcolor[rgb]{0.53,0.00,0.00}{{#1}}}
    \newcommand{\SpecialCharTok}[1]{\textcolor[rgb]{0.25,0.44,0.63}{{#1}}}
    \newcommand{\VerbatimStringTok}[1]{\textcolor[rgb]{0.25,0.44,0.63}{{#1}}}
    \newcommand{\SpecialStringTok}[1]{\textcolor[rgb]{0.73,0.40,0.53}{{#1}}}
    \newcommand{\ImportTok}[1]{{#1}}
    \newcommand{\DocumentationTok}[1]{\textcolor[rgb]{0.73,0.13,0.13}{\textit{{#1}}}}
    \newcommand{\AnnotationTok}[1]{\textcolor[rgb]{0.38,0.63,0.69}{\textbf{\textit{{#1}}}}}
    \newcommand{\CommentVarTok}[1]{\textcolor[rgb]{0.38,0.63,0.69}{\textbf{\textit{{#1}}}}}
    \newcommand{\VariableTok}[1]{\textcolor[rgb]{0.10,0.09,0.49}{{#1}}}
    \newcommand{\ControlFlowTok}[1]{\textcolor[rgb]{0.00,0.44,0.13}{\textbf{{#1}}}}
    \newcommand{\OperatorTok}[1]{\textcolor[rgb]{0.40,0.40,0.40}{{#1}}}
    \newcommand{\BuiltInTok}[1]{{#1}}
    \newcommand{\ExtensionTok}[1]{{#1}}
    \newcommand{\PreprocessorTok}[1]{\textcolor[rgb]{0.74,0.48,0.00}{{#1}}}
    \newcommand{\AttributeTok}[1]{\textcolor[rgb]{0.49,0.56,0.16}{{#1}}}
    \newcommand{\InformationTok}[1]{\textcolor[rgb]{0.38,0.63,0.69}{\textbf{\textit{{#1}}}}}
    \newcommand{\WarningTok}[1]{\textcolor[rgb]{0.38,0.63,0.69}{\textbf{\textit{{#1}}}}}
    \makeatletter
    \newsavebox\pandoc@box
    \newcommand*\pandocbounded[1]{%
      \sbox\pandoc@box{#1}%
      % scaling factors for width and height
      \Gscale@div\@tempa\textheight{\dimexpr\ht\pandoc@box+\dp\pandoc@box\relax}%
      \Gscale@div\@tempb\linewidth{\wd\pandoc@box}%
      % select the smaller of both
      \ifdim\@tempb\p@<\@tempa\p@
        \let\@tempa\@tempb
      \fi
      % scaling accordingly (\@tempa < 1)
      \ifdim\@tempa\p@<\p@
        \scalebox{\@tempa}{\usebox\pandoc@box}%
      % scaling not needed, use as it is
      \else
        \usebox{\pandoc@box}%
      \fi
    }
    \makeatother

    % Define a nice break command that doesn't care if a line doesn't already
    % exist.
    \def\br{\hspace*{\fill} \\* }
    % Math Jax compatibility definitions
    \def\gt{>}
    \def\lt{<}
    \let\Oldtex\TeX
    \let\Oldlatex\LaTeX
    \renewcommand{\TeX}{\textrm{\Oldtex}}
    \renewcommand{\LaTeX}{\textrm{\Oldlatex}}
    % Document parameters
    % Document title
    \title{Ch10-deeplearning-lab}
    
    
    
    
    
    
    
% Pygments definitions
\makeatletter
\def\PY@reset{\let\PY@it=\relax \let\PY@bf=\relax%
    \let\PY@ul=\relax \let\PY@tc=\relax%
    \let\PY@bc=\relax \let\PY@ff=\relax}
\def\PY@tok#1{\csname PY@tok@#1\endcsname}
\def\PY@toks#1+{\ifx\relax#1\empty\else%
    \PY@tok{#1}\expandafter\PY@toks\fi}
\def\PY@do#1{\PY@bc{\PY@tc{\PY@ul{%
    \PY@it{\PY@bf{\PY@ff{#1}}}}}}}
\def\PY#1#2{\PY@reset\PY@toks#1+\relax+\PY@do{#2}}

\@namedef{PY@tok@w}{\def\PY@tc##1{\textcolor[rgb]{0.73,0.73,0.73}{##1}}}
\@namedef{PY@tok@c}{\let\PY@it=\textit\def\PY@tc##1{\textcolor[rgb]{0.24,0.48,0.48}{##1}}}
\@namedef{PY@tok@cp}{\def\PY@tc##1{\textcolor[rgb]{0.61,0.40,0.00}{##1}}}
\@namedef{PY@tok@k}{\let\PY@bf=\textbf\def\PY@tc##1{\textcolor[rgb]{0.00,0.50,0.00}{##1}}}
\@namedef{PY@tok@kp}{\def\PY@tc##1{\textcolor[rgb]{0.00,0.50,0.00}{##1}}}
\@namedef{PY@tok@kt}{\def\PY@tc##1{\textcolor[rgb]{0.69,0.00,0.25}{##1}}}
\@namedef{PY@tok@o}{\def\PY@tc##1{\textcolor[rgb]{0.40,0.40,0.40}{##1}}}
\@namedef{PY@tok@ow}{\let\PY@bf=\textbf\def\PY@tc##1{\textcolor[rgb]{0.67,0.13,1.00}{##1}}}
\@namedef{PY@tok@nb}{\def\PY@tc##1{\textcolor[rgb]{0.00,0.50,0.00}{##1}}}
\@namedef{PY@tok@nf}{\def\PY@tc##1{\textcolor[rgb]{0.00,0.00,1.00}{##1}}}
\@namedef{PY@tok@nc}{\let\PY@bf=\textbf\def\PY@tc##1{\textcolor[rgb]{0.00,0.00,1.00}{##1}}}
\@namedef{PY@tok@nn}{\let\PY@bf=\textbf\def\PY@tc##1{\textcolor[rgb]{0.00,0.00,1.00}{##1}}}
\@namedef{PY@tok@ne}{\let\PY@bf=\textbf\def\PY@tc##1{\textcolor[rgb]{0.80,0.25,0.22}{##1}}}
\@namedef{PY@tok@nv}{\def\PY@tc##1{\textcolor[rgb]{0.10,0.09,0.49}{##1}}}
\@namedef{PY@tok@no}{\def\PY@tc##1{\textcolor[rgb]{0.53,0.00,0.00}{##1}}}
\@namedef{PY@tok@nl}{\def\PY@tc##1{\textcolor[rgb]{0.46,0.46,0.00}{##1}}}
\@namedef{PY@tok@ni}{\let\PY@bf=\textbf\def\PY@tc##1{\textcolor[rgb]{0.44,0.44,0.44}{##1}}}
\@namedef{PY@tok@na}{\def\PY@tc##1{\textcolor[rgb]{0.41,0.47,0.13}{##1}}}
\@namedef{PY@tok@nt}{\let\PY@bf=\textbf\def\PY@tc##1{\textcolor[rgb]{0.00,0.50,0.00}{##1}}}
\@namedef{PY@tok@nd}{\def\PY@tc##1{\textcolor[rgb]{0.67,0.13,1.00}{##1}}}
\@namedef{PY@tok@s}{\def\PY@tc##1{\textcolor[rgb]{0.73,0.13,0.13}{##1}}}
\@namedef{PY@tok@sd}{\let\PY@it=\textit\def\PY@tc##1{\textcolor[rgb]{0.73,0.13,0.13}{##1}}}
\@namedef{PY@tok@si}{\let\PY@bf=\textbf\def\PY@tc##1{\textcolor[rgb]{0.64,0.35,0.47}{##1}}}
\@namedef{PY@tok@se}{\let\PY@bf=\textbf\def\PY@tc##1{\textcolor[rgb]{0.67,0.36,0.12}{##1}}}
\@namedef{PY@tok@sr}{\def\PY@tc##1{\textcolor[rgb]{0.64,0.35,0.47}{##1}}}
\@namedef{PY@tok@ss}{\def\PY@tc##1{\textcolor[rgb]{0.10,0.09,0.49}{##1}}}
\@namedef{PY@tok@sx}{\def\PY@tc##1{\textcolor[rgb]{0.00,0.50,0.00}{##1}}}
\@namedef{PY@tok@m}{\def\PY@tc##1{\textcolor[rgb]{0.40,0.40,0.40}{##1}}}
\@namedef{PY@tok@gh}{\let\PY@bf=\textbf\def\PY@tc##1{\textcolor[rgb]{0.00,0.00,0.50}{##1}}}
\@namedef{PY@tok@gu}{\let\PY@bf=\textbf\def\PY@tc##1{\textcolor[rgb]{0.50,0.00,0.50}{##1}}}
\@namedef{PY@tok@gd}{\def\PY@tc##1{\textcolor[rgb]{0.63,0.00,0.00}{##1}}}
\@namedef{PY@tok@gi}{\def\PY@tc##1{\textcolor[rgb]{0.00,0.52,0.00}{##1}}}
\@namedef{PY@tok@gr}{\def\PY@tc##1{\textcolor[rgb]{0.89,0.00,0.00}{##1}}}
\@namedef{PY@tok@ge}{\let\PY@it=\textit}
\@namedef{PY@tok@gs}{\let\PY@bf=\textbf}
\@namedef{PY@tok@ges}{\let\PY@bf=\textbf\let\PY@it=\textit}
\@namedef{PY@tok@gp}{\let\PY@bf=\textbf\def\PY@tc##1{\textcolor[rgb]{0.00,0.00,0.50}{##1}}}
\@namedef{PY@tok@go}{\def\PY@tc##1{\textcolor[rgb]{0.44,0.44,0.44}{##1}}}
\@namedef{PY@tok@gt}{\def\PY@tc##1{\textcolor[rgb]{0.00,0.27,0.87}{##1}}}
\@namedef{PY@tok@err}{\def\PY@bc##1{{\setlength{\fboxsep}{\string -\fboxrule}\fcolorbox[rgb]{1.00,0.00,0.00}{1,1,1}{\strut ##1}}}}
\@namedef{PY@tok@kc}{\let\PY@bf=\textbf\def\PY@tc##1{\textcolor[rgb]{0.00,0.50,0.00}{##1}}}
\@namedef{PY@tok@kd}{\let\PY@bf=\textbf\def\PY@tc##1{\textcolor[rgb]{0.00,0.50,0.00}{##1}}}
\@namedef{PY@tok@kn}{\let\PY@bf=\textbf\def\PY@tc##1{\textcolor[rgb]{0.00,0.50,0.00}{##1}}}
\@namedef{PY@tok@kr}{\let\PY@bf=\textbf\def\PY@tc##1{\textcolor[rgb]{0.00,0.50,0.00}{##1}}}
\@namedef{PY@tok@bp}{\def\PY@tc##1{\textcolor[rgb]{0.00,0.50,0.00}{##1}}}
\@namedef{PY@tok@fm}{\def\PY@tc##1{\textcolor[rgb]{0.00,0.00,1.00}{##1}}}
\@namedef{PY@tok@vc}{\def\PY@tc##1{\textcolor[rgb]{0.10,0.09,0.49}{##1}}}
\@namedef{PY@tok@vg}{\def\PY@tc##1{\textcolor[rgb]{0.10,0.09,0.49}{##1}}}
\@namedef{PY@tok@vi}{\def\PY@tc##1{\textcolor[rgb]{0.10,0.09,0.49}{##1}}}
\@namedef{PY@tok@vm}{\def\PY@tc##1{\textcolor[rgb]{0.10,0.09,0.49}{##1}}}
\@namedef{PY@tok@sa}{\def\PY@tc##1{\textcolor[rgb]{0.73,0.13,0.13}{##1}}}
\@namedef{PY@tok@sb}{\def\PY@tc##1{\textcolor[rgb]{0.73,0.13,0.13}{##1}}}
\@namedef{PY@tok@sc}{\def\PY@tc##1{\textcolor[rgb]{0.73,0.13,0.13}{##1}}}
\@namedef{PY@tok@dl}{\def\PY@tc##1{\textcolor[rgb]{0.73,0.13,0.13}{##1}}}
\@namedef{PY@tok@s2}{\def\PY@tc##1{\textcolor[rgb]{0.73,0.13,0.13}{##1}}}
\@namedef{PY@tok@sh}{\def\PY@tc##1{\textcolor[rgb]{0.73,0.13,0.13}{##1}}}
\@namedef{PY@tok@s1}{\def\PY@tc##1{\textcolor[rgb]{0.73,0.13,0.13}{##1}}}
\@namedef{PY@tok@mb}{\def\PY@tc##1{\textcolor[rgb]{0.40,0.40,0.40}{##1}}}
\@namedef{PY@tok@mf}{\def\PY@tc##1{\textcolor[rgb]{0.40,0.40,0.40}{##1}}}
\@namedef{PY@tok@mh}{\def\PY@tc##1{\textcolor[rgb]{0.40,0.40,0.40}{##1}}}
\@namedef{PY@tok@mi}{\def\PY@tc##1{\textcolor[rgb]{0.40,0.40,0.40}{##1}}}
\@namedef{PY@tok@il}{\def\PY@tc##1{\textcolor[rgb]{0.40,0.40,0.40}{##1}}}
\@namedef{PY@tok@mo}{\def\PY@tc##1{\textcolor[rgb]{0.40,0.40,0.40}{##1}}}
\@namedef{PY@tok@ch}{\let\PY@it=\textit\def\PY@tc##1{\textcolor[rgb]{0.24,0.48,0.48}{##1}}}
\@namedef{PY@tok@cm}{\let\PY@it=\textit\def\PY@tc##1{\textcolor[rgb]{0.24,0.48,0.48}{##1}}}
\@namedef{PY@tok@cpf}{\let\PY@it=\textit\def\PY@tc##1{\textcolor[rgb]{0.24,0.48,0.48}{##1}}}
\@namedef{PY@tok@c1}{\let\PY@it=\textit\def\PY@tc##1{\textcolor[rgb]{0.24,0.48,0.48}{##1}}}
\@namedef{PY@tok@cs}{\let\PY@it=\textit\def\PY@tc##1{\textcolor[rgb]{0.24,0.48,0.48}{##1}}}

\def\PYZbs{\char`\\}
\def\PYZus{\char`\_}
\def\PYZob{\char`\{}
\def\PYZcb{\char`\}}
\def\PYZca{\char`\^}
\def\PYZam{\char`\&}
\def\PYZlt{\char`\<}
\def\PYZgt{\char`\>}
\def\PYZsh{\char`\#}
\def\PYZpc{\char`\%}
\def\PYZdl{\char`\$}
\def\PYZhy{\char`\-}
\def\PYZsq{\char`\'}
\def\PYZdq{\char`\"}
\def\PYZti{\char`\~}
% for compatibility with earlier versions
\def\PYZat{@}
\def\PYZlb{[}
\def\PYZrb{]}
\makeatother


    % For linebreaks inside Verbatim environment from package fancyvrb.
    \makeatletter
        \newbox\Wrappedcontinuationbox
        \newbox\Wrappedvisiblespacebox
        \newcommand*\Wrappedvisiblespace {\textcolor{red}{\textvisiblespace}}
        \newcommand*\Wrappedcontinuationsymbol {\textcolor{red}{\llap{\tiny$\m@th\hookrightarrow$}}}
        \newcommand*\Wrappedcontinuationindent {3ex }
        \newcommand*\Wrappedafterbreak {\kern\Wrappedcontinuationindent\copy\Wrappedcontinuationbox}
        % Take advantage of the already applied Pygments mark-up to insert
        % potential linebreaks for TeX processing.
        %        {, <, #, %, $, ' and ": go to next line.
        %        _, }, ^, &, >, - and ~: stay at end of broken line.
        % Use of \textquotesingle for straight quote.
        \newcommand*\Wrappedbreaksatspecials {%
            \def\PYGZus{\discretionary{\char`\_}{\Wrappedafterbreak}{\char`\_}}%
            \def\PYGZob{\discretionary{}{\Wrappedafterbreak\char`\{}{\char`\{}}%
            \def\PYGZcb{\discretionary{\char`\}}{\Wrappedafterbreak}{\char`\}}}%
            \def\PYGZca{\discretionary{\char`\^}{\Wrappedafterbreak}{\char`\^}}%
            \def\PYGZam{\discretionary{\char`\&}{\Wrappedafterbreak}{\char`\&}}%
            \def\PYGZlt{\discretionary{}{\Wrappedafterbreak\char`\<}{\char`\<}}%
            \def\PYGZgt{\discretionary{\char`\>}{\Wrappedafterbreak}{\char`\>}}%
            \def\PYGZsh{\discretionary{}{\Wrappedafterbreak\char`\#}{\char`\#}}%
            \def\PYGZpc{\discretionary{}{\Wrappedafterbreak\char`\%}{\char`\%}}%
            \def\PYGZdl{\discretionary{}{\Wrappedafterbreak\char`\$}{\char`\$}}%
            \def\PYGZhy{\discretionary{\char`\-}{\Wrappedafterbreak}{\char`\-}}%
            \def\PYGZsq{\discretionary{}{\Wrappedafterbreak\textquotesingle}{\textquotesingle}}%
            \def\PYGZdq{\discretionary{}{\Wrappedafterbreak\char`\"}{\char`\"}}%
            \def\PYGZti{\discretionary{\char`\~}{\Wrappedafterbreak}{\char`\~}}%
        }
        % Some characters . , ; ? ! / are not pygmentized.
        % This macro makes them "active" and they will insert potential linebreaks
        \newcommand*\Wrappedbreaksatpunct {%
            \lccode`\~`\.\lowercase{\def~}{\discretionary{\hbox{\char`\.}}{\Wrappedafterbreak}{\hbox{\char`\.}}}%
            \lccode`\~`\,\lowercase{\def~}{\discretionary{\hbox{\char`\,}}{\Wrappedafterbreak}{\hbox{\char`\,}}}%
            \lccode`\~`\;\lowercase{\def~}{\discretionary{\hbox{\char`\;}}{\Wrappedafterbreak}{\hbox{\char`\;}}}%
            \lccode`\~`\:\lowercase{\def~}{\discretionary{\hbox{\char`\:}}{\Wrappedafterbreak}{\hbox{\char`\:}}}%
            \lccode`\~`\?\lowercase{\def~}{\discretionary{\hbox{\char`\?}}{\Wrappedafterbreak}{\hbox{\char`\?}}}%
            \lccode`\~`\!\lowercase{\def~}{\discretionary{\hbox{\char`\!}}{\Wrappedafterbreak}{\hbox{\char`\!}}}%
            \lccode`\~`\/\lowercase{\def~}{\discretionary{\hbox{\char`\/}}{\Wrappedafterbreak}{\hbox{\char`\/}}}%
            \catcode`\.\active
            \catcode`\,\active
            \catcode`\;\active
            \catcode`\:\active
            \catcode`\?\active
            \catcode`\!\active
            \catcode`\/\active
            \lccode`\~`\~
        }
    \makeatother

    \let\OriginalVerbatim=\Verbatim
    \makeatletter
    \renewcommand{\Verbatim}[1][1]{%
        %\parskip\z@skip
        \sbox\Wrappedcontinuationbox {\Wrappedcontinuationsymbol}%
        \sbox\Wrappedvisiblespacebox {\FV@SetupFont\Wrappedvisiblespace}%
        \def\FancyVerbFormatLine ##1{\hsize\linewidth
            \vtop{\raggedright\hyphenpenalty\z@\exhyphenpenalty\z@
                \doublehyphendemerits\z@\finalhyphendemerits\z@
                \strut ##1\strut}%
        }%
        % If the linebreak is at a space, the latter will be displayed as visible
        % space at end of first line, and a continuation symbol starts next line.
        % Stretch/shrink are however usually zero for typewriter font.
        \def\FV@Space {%
            \nobreak\hskip\z@ plus\fontdimen3\font minus\fontdimen4\font
            \discretionary{\copy\Wrappedvisiblespacebox}{\Wrappedafterbreak}
            {\kern\fontdimen2\font}%
        }%

        % Allow breaks at special characters using \PYG... macros.
        \Wrappedbreaksatspecials
        % Breaks at punctuation characters . , ; ? ! and / need catcode=\active
        \OriginalVerbatim[#1,codes*=\Wrappedbreaksatpunct]%
    }
    \makeatother

    % Exact colors from NB
    \definecolor{incolor}{HTML}{303F9F}
    \definecolor{outcolor}{HTML}{D84315}
    \definecolor{cellborder}{HTML}{CFCFCF}
    \definecolor{cellbackground}{HTML}{F7F7F7}

    % prompt
    \makeatletter
    \newcommand{\boxspacing}{\kern\kvtcb@left@rule\kern\kvtcb@boxsep}
    \makeatother
    \newcommand{\prompt}[4]{
        {\ttfamily\llap{{\color{#2}[#3]:\hspace{3pt}#4}}\vspace{-\baselineskip}}
    }
    

    
    % Prevent overflowing lines due to hard-to-break entities
    \sloppy
    % Setup hyperref package
    \hypersetup{
      breaklinks=true,  % so long urls are correctly broken across lines
      colorlinks=true,
      urlcolor=urlcolor,
      linkcolor=linkcolor,
      citecolor=citecolor,
      }
    % Slightly bigger margins than the latex defaults
    
    \geometry{verbose,tmargin=1in,bmargin=1in,lmargin=1in,rmargin=1in}
    
    

\begin{document}
    
    \maketitle
    
    

    
    \hypertarget{deep-learning}{%
\section{Deep Learning}\label{deep-learning}}

    In this section we demonstrate how to fit the examples discussed in the
text. We use the \texttt{Python} \texttt{torch} package, along with the
\texttt{pytorch\_lightning} package which provides utilities to simplify
fitting and evaluating models. This code can be impressively fast with
certain special processors, such as Apple's new M1 chip. The package is
well-structured, flexible, and will feel comfortable to \texttt{Python}
users. A good companion is the site
\href{https://pytorch.org/tutorials/beginner/basics/intro.html}{pytorch.org/tutorials}.
Much of our code is adapted from there, as well as the
\texttt{pytorch\_lightning} documentation. \{The precise URLs at the
time of writing are
\url{https://pytorch.org/tutorials/beginner/basics/intro.html} and
\url{https://pytorch-lightning.readthedocs.io/en/latest/}.\}

We start with several standard imports that we have seen before.

    \begin{tcolorbox}[breakable, size=fbox, boxrule=1pt, pad at break*=1mm,colback=cellbackground, colframe=cellborder]
\prompt{In}{incolor}{1}{\boxspacing}
\begin{Verbatim}[commandchars=\\\{\}]
\PY{k+kn}{import}\PY{+w}{ }\PY{n+nn}{numpy}\PY{+w}{ }\PY{k}{as}\PY{+w}{ }\PY{n+nn}{np}\PY{o}{,}\PY{+w}{ }\PY{n+nn}{pandas}\PY{+w}{ }\PY{k}{as}\PY{+w}{ }\PY{n+nn}{pd}
\PY{k+kn}{from}\PY{+w}{ }\PY{n+nn}{matplotlib}\PY{n+nn}{.}\PY{n+nn}{pyplot}\PY{+w}{ }\PY{k+kn}{import} \PY{n}{subplots}
\PY{k+kn}{from}\PY{+w}{ }\PY{n+nn}{sklearn}\PY{n+nn}{.}\PY{n+nn}{linear\PYZus{}model}\PY{+w}{ }\PY{k+kn}{import} \PYZbs{}
     \PY{p}{(}\PY{n}{LinearRegression}\PY{p}{,}
      \PY{n}{LogisticRegression}\PY{p}{,}
      \PY{n}{Lasso}\PY{p}{)}
\PY{k+kn}{from}\PY{+w}{ }\PY{n+nn}{sklearn}\PY{n+nn}{.}\PY{n+nn}{preprocessing}\PY{+w}{ }\PY{k+kn}{import} \PY{n}{StandardScaler}
\PY{k+kn}{from}\PY{+w}{ }\PY{n+nn}{sklearn}\PY{n+nn}{.}\PY{n+nn}{model\PYZus{}selection}\PY{+w}{ }\PY{k+kn}{import} \PY{n}{KFold}
\PY{k+kn}{from}\PY{+w}{ }\PY{n+nn}{sklearn}\PY{n+nn}{.}\PY{n+nn}{pipeline}\PY{+w}{ }\PY{k+kn}{import} \PY{n}{Pipeline}
\PY{k+kn}{from}\PY{+w}{ }\PY{n+nn}{ISLP}\PY{+w}{ }\PY{k+kn}{import} \PY{n}{load\PYZus{}data}
\PY{k+kn}{from}\PY{+w}{ }\PY{n+nn}{ISLP}\PY{n+nn}{.}\PY{n+nn}{models}\PY{+w}{ }\PY{k+kn}{import} \PY{n}{ModelSpec} \PY{k}{as} \PY{n}{MS}
\PY{k+kn}{from}\PY{+w}{ }\PY{n+nn}{sklearn}\PY{n+nn}{.}\PY{n+nn}{model\PYZus{}selection}\PY{+w}{ }\PY{k+kn}{import} \PYZbs{}
     \PY{p}{(}\PY{n}{train\PYZus{}test\PYZus{}split}\PY{p}{,}
      \PY{n}{GridSearchCV}\PY{p}{)}
\end{Verbatim}
\end{tcolorbox}

    \hypertarget{torch-specific-imports}{%
\subsubsection{Torch-Specific Imports}\label{torch-specific-imports}}

There are a number of imports for \texttt{torch}. (These are not
included with \texttt{ISLP}, so must be installed separately.) First we
import the main library and essential tools used to specify
sequentially-structured networks.

    \begin{tcolorbox}[breakable, size=fbox, boxrule=1pt, pad at break*=1mm,colback=cellbackground, colframe=cellborder]
\prompt{In}{incolor}{2}{\boxspacing}
\begin{Verbatim}[commandchars=\\\{\}]
\PY{k+kn}{import}\PY{+w}{ }\PY{n+nn}{torch}
\PY{k+kn}{from}\PY{+w}{ }\PY{n+nn}{torch}\PY{+w}{ }\PY{k+kn}{import} \PY{n}{nn}
\PY{k+kn}{from}\PY{+w}{ }\PY{n+nn}{torch}\PY{n+nn}{.}\PY{n+nn}{optim}\PY{+w}{ }\PY{k+kn}{import} \PY{n}{RMSprop}
\PY{k+kn}{from}\PY{+w}{ }\PY{n+nn}{torch}\PY{n+nn}{.}\PY{n+nn}{utils}\PY{n+nn}{.}\PY{n+nn}{data}\PY{+w}{ }\PY{k+kn}{import} \PY{n}{TensorDataset}
\end{Verbatim}
\end{tcolorbox}

    There are several other helper packages for \texttt{torch}. For
instance, the \texttt{torchmetrics} package has utilities to compute
various metrics to evaluate performance when fitting a model. The
\texttt{torchinfo} package provides a useful summary of the layers of a
model. We use the \texttt{read\_image()} function when loading test
images in Section 10.9.4.

If you have not already installed the packages \texttt{torchvision} and
\texttt{torchinfo} you can install them by running
\texttt{pip\ install\ torchinfo\ torchvision}. We can now import from
\texttt{torchinfo}.

    \begin{tcolorbox}[breakable, size=fbox, boxrule=1pt, pad at break*=1mm,colback=cellbackground, colframe=cellborder]
\prompt{In}{incolor}{3}{\boxspacing}
\begin{Verbatim}[commandchars=\\\{\}]
\PY{k+kn}{from}\PY{+w}{ }\PY{n+nn}{torchmetrics}\PY{+w}{ }\PY{k+kn}{import} \PY{p}{(}\PY{n}{MeanAbsoluteError}\PY{p}{,}
                          \PY{n}{R2Score}\PY{p}{)}
\PY{k+kn}{from}\PY{+w}{ }\PY{n+nn}{torchinfo}\PY{+w}{ }\PY{k+kn}{import} \PY{n}{summary}
\end{Verbatim}
\end{tcolorbox}

    The package \texttt{pytorch\_lightning} is a somewhat higher-level
interface to \texttt{torch} that simplifies the specification and
fitting of models by reducing the amount of boilerplate code needed
(compared to using \texttt{torch} alone).

    \begin{tcolorbox}[breakable, size=fbox, boxrule=1pt, pad at break*=1mm,colback=cellbackground, colframe=cellborder]
\prompt{In}{incolor}{4}{\boxspacing}
\begin{Verbatim}[commandchars=\\\{\}]
\PY{k+kn}{from}\PY{+w}{ }\PY{n+nn}{pytorch\PYZus{}lightning}\PY{+w}{ }\PY{k+kn}{import} \PY{n}{Trainer}
\PY{k+kn}{from}\PY{+w}{ }\PY{n+nn}{pytorch\PYZus{}lightning}\PY{n+nn}{.}\PY{n+nn}{loggers}\PY{+w}{ }\PY{k+kn}{import} \PY{n}{CSVLogger}
\end{Verbatim}
\end{tcolorbox}

    In order to reproduce results we use \texttt{seed\_everything()}. We
will also instruct \texttt{torch} to use deterministic algorithms where
possible.

    \begin{tcolorbox}[breakable, size=fbox, boxrule=1pt, pad at break*=1mm,colback=cellbackground, colframe=cellborder]
\prompt{In}{incolor}{5}{\boxspacing}
\begin{Verbatim}[commandchars=\\\{\}]
\PY{k+kn}{from}\PY{+w}{ }\PY{n+nn}{pytorch\PYZus{}lightning}\PY{+w}{ }\PY{k+kn}{import} \PY{n}{seed\PYZus{}everything}
\PY{n}{seed\PYZus{}everything}\PY{p}{(}\PY{l+m+mi}{0}\PY{p}{,} \PY{n}{workers}\PY{o}{=}\PY{k+kc}{True}\PY{p}{)}
\PY{n}{torch}\PY{o}{.}\PY{n}{use\PYZus{}deterministic\PYZus{}algorithms}\PY{p}{(}\PY{k+kc}{True}\PY{p}{,} \PY{n}{warn\PYZus{}only}\PY{o}{=}\PY{k+kc}{True}\PY{p}{)}
\end{Verbatim}
\end{tcolorbox}

    \begin{Verbatim}[commandchars=\\\{\}]
Seed set to 0
    \end{Verbatim}

    We will use several datasets shipped with \texttt{torchvision} for our
examples: a pretrained network for image classification, as well as some
transforms used for preprocessing.

    \begin{tcolorbox}[breakable, size=fbox, boxrule=1pt, pad at break*=1mm,colback=cellbackground, colframe=cellborder]
\prompt{In}{incolor}{6}{\boxspacing}
\begin{Verbatim}[commandchars=\\\{\}]
\PY{k+kn}{from}\PY{+w}{ }\PY{n+nn}{torchvision}\PY{n+nn}{.}\PY{n+nn}{io}\PY{+w}{ }\PY{k+kn}{import} \PY{n}{read\PYZus{}image}
\PY{k+kn}{from}\PY{+w}{ }\PY{n+nn}{torchvision}\PY{n+nn}{.}\PY{n+nn}{datasets}\PY{+w}{ }\PY{k+kn}{import} \PY{n}{MNIST}\PY{p}{,} \PY{n}{CIFAR100}
\PY{k+kn}{from}\PY{+w}{ }\PY{n+nn}{torchvision}\PY{n+nn}{.}\PY{n+nn}{models}\PY{+w}{ }\PY{k+kn}{import} \PY{p}{(}\PY{n}{resnet50}\PY{p}{,}
                                \PY{n}{ResNet50\PYZus{}Weights}\PY{p}{)}
\PY{k+kn}{from}\PY{+w}{ }\PY{n+nn}{torchvision}\PY{n+nn}{.}\PY{n+nn}{transforms}\PY{+w}{ }\PY{k+kn}{import} \PY{p}{(}\PY{n}{Resize}\PY{p}{,}
                                    \PY{n}{Normalize}\PY{p}{,}
                                    \PY{n}{CenterCrop}\PY{p}{,}
                                    \PY{n}{ToTensor}\PY{p}{)}
\end{Verbatim}
\end{tcolorbox}

    We have provided a few utilities in \texttt{ISLP} specifically for this
lab. The \texttt{SimpleDataModule} and \texttt{SimpleModule} are simple
versions of objects used in \texttt{pytorch\_lightning}, the high-level
module for fitting \texttt{torch} models. Although more advanced uses
such as computing on graphical processing units (GPUs) and parallel data
processing are possible in this module, we will not be focusing much on
these in this lab. The \texttt{ErrorTracker} handles collections of
targets and predictions over each mini-batch in the validation or test
stage, allowing computation of the metric over the entire validation or
test data set.

    \begin{tcolorbox}[breakable, size=fbox, boxrule=1pt, pad at break*=1mm,colback=cellbackground, colframe=cellborder]
\prompt{In}{incolor}{7}{\boxspacing}
\begin{Verbatim}[commandchars=\\\{\}]
\PY{k+kn}{from}\PY{+w}{ }\PY{n+nn}{ISLP}\PY{n+nn}{.}\PY{n+nn}{torch}\PY{+w}{ }\PY{k+kn}{import} \PY{p}{(}\PY{n}{SimpleDataModule}\PY{p}{,}
                        \PY{n}{SimpleModule}\PY{p}{,}
                        \PY{n}{ErrorTracker}\PY{p}{,}
                        \PY{n}{rec\PYZus{}num\PYZus{}workers}\PY{p}{)}
\end{Verbatim}
\end{tcolorbox}

    In addition we have included some helper functions to load the
\texttt{IMDb} database, as well as a lookup that maps integers to
particular keys in the database. We've included a slightly modified copy
of the preprocessed \texttt{IMDb} data from \texttt{keras}, a separate
package for fitting deep learning models. This saves us significant
preprocessing and allows us to focus on specifying and fitting the
models themselves.

    \begin{tcolorbox}[breakable, size=fbox, boxrule=1pt, pad at break*=1mm,colback=cellbackground, colframe=cellborder]
\prompt{In}{incolor}{8}{\boxspacing}
\begin{Verbatim}[commandchars=\\\{\}]
\PY{k+kn}{from}\PY{+w}{ }\PY{n+nn}{ISLP}\PY{n+nn}{.}\PY{n+nn}{torch}\PY{n+nn}{.}\PY{n+nn}{imdb}\PY{+w}{ }\PY{k+kn}{import} \PY{p}{(}\PY{n}{load\PYZus{}lookup}\PY{p}{,}
                             \PY{n}{load\PYZus{}tensor}\PY{p}{,}
                             \PY{n}{load\PYZus{}sparse}\PY{p}{,}
                             \PY{n}{load\PYZus{}sequential}\PY{p}{)}
\end{Verbatim}
\end{tcolorbox}

    Finally, we introduce some utility imports not directly related to
\texttt{torch}. The \texttt{glob()} function from the \texttt{glob}
module is used to find all files matching wildcard characters, which we
will use in our example applying the \texttt{ResNet50} model to some of
our own images. The \texttt{json} module will be used to load a JSON
file for looking up classes to identify the labels of the pictures in
the \texttt{ResNet50} example.

    \begin{tcolorbox}[breakable, size=fbox, boxrule=1pt, pad at break*=1mm,colback=cellbackground, colframe=cellborder]
\prompt{In}{incolor}{9}{\boxspacing}
\begin{Verbatim}[commandchars=\\\{\}]
\PY{k+kn}{from}\PY{+w}{ }\PY{n+nn}{glob}\PY{+w}{ }\PY{k+kn}{import} \PY{n}{glob}
\PY{k+kn}{import}\PY{+w}{ }\PY{n+nn}{json}
\end{Verbatim}
\end{tcolorbox}

    \hypertarget{single-layer-network-on-hitters-data}{%
\subsection{Single Layer Network on Hitters
Data}\label{single-layer-network-on-hitters-data}}

We start by fitting the models in Section 10.6 on the \texttt{Hitters}
data.

    \begin{tcolorbox}[breakable, size=fbox, boxrule=1pt, pad at break*=1mm,colback=cellbackground, colframe=cellborder]
\prompt{In}{incolor}{10}{\boxspacing}
\begin{Verbatim}[commandchars=\\\{\}]
\PY{n}{Hitters} \PY{o}{=} \PY{n}{load\PYZus{}data}\PY{p}{(}\PY{l+s+s1}{\PYZsq{}}\PY{l+s+s1}{Hitters}\PY{l+s+s1}{\PYZsq{}}\PY{p}{)}\PY{o}{.}\PY{n}{dropna}\PY{p}{(}\PY{p}{)}
\PY{n}{n} \PY{o}{=} \PY{n}{Hitters}\PY{o}{.}\PY{n}{shape}\PY{p}{[}\PY{l+m+mi}{0}\PY{p}{]}
\end{Verbatim}
\end{tcolorbox}

    We will fit two linear models (least squares and lasso) and compare
their performance to that of a neural network. For this comparison we
will use mean absolute error on a validation dataset. \begin{equation*}
\begin{split}
\mbox{MAE}(y,\hat{y}) = \frac{1}{n} \sum_{i=1}^n |y_i-\hat{y}_i|.
\end{split}
\end{equation*} We set up the model matrix and the response.

    \begin{tcolorbox}[breakable, size=fbox, boxrule=1pt, pad at break*=1mm,colback=cellbackground, colframe=cellborder]
\prompt{In}{incolor}{11}{\boxspacing}
\begin{Verbatim}[commandchars=\\\{\}]
\PY{n}{model} \PY{o}{=} \PY{n}{MS}\PY{p}{(}\PY{n}{Hitters}\PY{o}{.}\PY{n}{columns}\PY{o}{.}\PY{n}{drop}\PY{p}{(}\PY{l+s+s1}{\PYZsq{}}\PY{l+s+s1}{Salary}\PY{l+s+s1}{\PYZsq{}}\PY{p}{)}\PY{p}{,} \PY{n}{intercept}\PY{o}{=}\PY{k+kc}{False}\PY{p}{)}
\PY{n}{X} \PY{o}{=} \PY{n}{model}\PY{o}{.}\PY{n}{fit\PYZus{}transform}\PY{p}{(}\PY{n}{Hitters}\PY{p}{)}\PY{o}{.}\PY{n}{to\PYZus{}numpy}\PY{p}{(}\PY{p}{)}
\PY{n}{Y} \PY{o}{=} \PY{n}{Hitters}\PY{p}{[}\PY{l+s+s1}{\PYZsq{}}\PY{l+s+s1}{Salary}\PY{l+s+s1}{\PYZsq{}}\PY{p}{]}\PY{o}{.}\PY{n}{to\PYZus{}numpy}\PY{p}{(}\PY{p}{)}
\end{Verbatim}
\end{tcolorbox}

    The \texttt{to\_numpy()} method above converts \texttt{pandas} data
frames or series to \texttt{numpy} arrays. We do this because we will
need to use \texttt{sklearn} to fit the lasso model, and it requires
this conversion. We also use a linear regression method from
\texttt{sklearn}, rather than the method in Chapter\textasciitilde3 from
\texttt{statsmodels}, to facilitate the comparisons.

    We now split the data into test and training, fixing the random state
used by \texttt{sklearn} to do the split.

    \begin{tcolorbox}[breakable, size=fbox, boxrule=1pt, pad at break*=1mm,colback=cellbackground, colframe=cellborder]
\prompt{In}{incolor}{12}{\boxspacing}
\begin{Verbatim}[commandchars=\\\{\}]
\PY{p}{(}\PY{n}{X\PYZus{}train}\PY{p}{,} 
 \PY{n}{X\PYZus{}test}\PY{p}{,}
 \PY{n}{Y\PYZus{}train}\PY{p}{,}
 \PY{n}{Y\PYZus{}test}\PY{p}{)} \PY{o}{=} \PY{n}{train\PYZus{}test\PYZus{}split}\PY{p}{(}\PY{n}{X}\PY{p}{,}
                            \PY{n}{Y}\PY{p}{,}
                            \PY{n}{test\PYZus{}size}\PY{o}{=}\PY{l+m+mi}{1}\PY{o}{/}\PY{l+m+mi}{3}\PY{p}{,}
                            \PY{n}{random\PYZus{}state}\PY{o}{=}\PY{l+m+mi}{1}\PY{p}{)}
\end{Verbatim}
\end{tcolorbox}

    \hypertarget{linear-models}{%
\subsubsection{Linear Models}\label{linear-models}}

We fit the linear model and evaluate the test error directly.

    \begin{tcolorbox}[breakable, size=fbox, boxrule=1pt, pad at break*=1mm,colback=cellbackground, colframe=cellborder]
\prompt{In}{incolor}{13}{\boxspacing}
\begin{Verbatim}[commandchars=\\\{\}]
\PY{n}{hit\PYZus{}lm} \PY{o}{=} \PY{n}{LinearRegression}\PY{p}{(}\PY{p}{)}\PY{o}{.}\PY{n}{fit}\PY{p}{(}\PY{n}{X\PYZus{}train}\PY{p}{,} \PY{n}{Y\PYZus{}train}\PY{p}{)}
\PY{n}{Yhat\PYZus{}test} \PY{o}{=} \PY{n}{hit\PYZus{}lm}\PY{o}{.}\PY{n}{predict}\PY{p}{(}\PY{n}{X\PYZus{}test}\PY{p}{)}
\PY{n}{np}\PY{o}{.}\PY{n}{abs}\PY{p}{(}\PY{n}{Yhat\PYZus{}test} \PY{o}{\PYZhy{}} \PY{n}{Y\PYZus{}test}\PY{p}{)}\PY{o}{.}\PY{n}{mean}\PY{p}{(}\PY{p}{)}
\end{Verbatim}
\end{tcolorbox}

            \begin{tcolorbox}[breakable, size=fbox, boxrule=.5pt, pad at break*=1mm, opacityfill=0]
\prompt{Out}{outcolor}{13}{\boxspacing}
\begin{Verbatim}[commandchars=\\\{\}]
259.71528833146243
\end{Verbatim}
\end{tcolorbox}
        
    Next we fit the lasso using \texttt{sklearn}. We are using mean absolute
error to select and evaluate a model, rather than mean squared error.
The specialized solver we used in Section 6.5.2 uses only mean squared
error. So here, with a bit more work, we create a cross-validation grid
and perform the cross-validation directly.

We encode a pipeline with two steps: we first normalize the features
using a \texttt{StandardScaler()} transform, and then fit the lasso
without further normalization.

    \begin{tcolorbox}[breakable, size=fbox, boxrule=1pt, pad at break*=1mm,colback=cellbackground, colframe=cellborder]
\prompt{In}{incolor}{14}{\boxspacing}
\begin{Verbatim}[commandchars=\\\{\}]
\PY{n}{scaler} \PY{o}{=} \PY{n}{StandardScaler}\PY{p}{(}\PY{n}{with\PYZus{}mean}\PY{o}{=}\PY{k+kc}{True}\PY{p}{,} \PY{n}{with\PYZus{}std}\PY{o}{=}\PY{k+kc}{True}\PY{p}{)}
\PY{n}{lasso} \PY{o}{=} \PY{n}{Lasso}\PY{p}{(}\PY{n}{warm\PYZus{}start}\PY{o}{=}\PY{k+kc}{True}\PY{p}{,} \PY{n}{max\PYZus{}iter}\PY{o}{=}\PY{l+m+mi}{30000}\PY{p}{)}
\PY{n}{standard\PYZus{}lasso} \PY{o}{=} \PY{n}{Pipeline}\PY{p}{(}\PY{n}{steps}\PY{o}{=}\PY{p}{[}\PY{p}{(}\PY{l+s+s1}{\PYZsq{}}\PY{l+s+s1}{scaler}\PY{l+s+s1}{\PYZsq{}}\PY{p}{,} \PY{n}{scaler}\PY{p}{)}\PY{p}{,}
                                 \PY{p}{(}\PY{l+s+s1}{\PYZsq{}}\PY{l+s+s1}{lasso}\PY{l+s+s1}{\PYZsq{}}\PY{p}{,} \PY{n}{lasso}\PY{p}{)}\PY{p}{]}\PY{p}{)}
\end{Verbatim}
\end{tcolorbox}

    We need to create a grid of values for \(\lambda\). As is common
practice, we choose a grid of 100 values of \(\lambda\), uniform on the
log scale from \texttt{lam\_max} down to \texttt{0.01*lam\_max}. Here
\texttt{lam\_max} is the smallest value of \(\lambda\) with an all-zero
solution. This value equals the largest absolute inner-product between
any predictor and the (centered) response. \{The derivation of this
result is beyond the scope of this book.\}

    \begin{tcolorbox}[breakable, size=fbox, boxrule=1pt, pad at break*=1mm,colback=cellbackground, colframe=cellborder]
\prompt{In}{incolor}{15}{\boxspacing}
\begin{Verbatim}[commandchars=\\\{\}]
\PY{n}{X\PYZus{}s} \PY{o}{=} \PY{n}{scaler}\PY{o}{.}\PY{n}{fit\PYZus{}transform}\PY{p}{(}\PY{n}{X\PYZus{}train}\PY{p}{)}
\PY{n}{n} \PY{o}{=} \PY{n}{X\PYZus{}s}\PY{o}{.}\PY{n}{shape}\PY{p}{[}\PY{l+m+mi}{0}\PY{p}{]}
\PY{n}{lam\PYZus{}max} \PY{o}{=} \PY{n}{np}\PY{o}{.}\PY{n}{fabs}\PY{p}{(}\PY{n}{X\PYZus{}s}\PY{o}{.}\PY{n}{T}\PY{o}{.}\PY{n}{dot}\PY{p}{(}\PY{n}{Y\PYZus{}train} \PY{o}{\PYZhy{}} \PY{n}{Y\PYZus{}train}\PY{o}{.}\PY{n}{mean}\PY{p}{(}\PY{p}{)}\PY{p}{)}\PY{p}{)}\PY{o}{.}\PY{n}{max}\PY{p}{(}\PY{p}{)} \PY{o}{/} \PY{n}{n}
\PY{n}{param\PYZus{}grid} \PY{o}{=} \PY{p}{\PYZob{}}\PY{l+s+s1}{\PYZsq{}}\PY{l+s+s1}{lasso\PYZus{}\PYZus{}alpha}\PY{l+s+s1}{\PYZsq{}}\PY{p}{:} \PY{n}{np}\PY{o}{.}\PY{n}{exp}\PY{p}{(}\PY{n}{np}\PY{o}{.}\PY{n}{linspace}\PY{p}{(}\PY{l+m+mi}{0}\PY{p}{,} \PY{n}{np}\PY{o}{.}\PY{n}{log}\PY{p}{(}\PY{l+m+mf}{0.01}\PY{p}{)}\PY{p}{,} \PY{l+m+mi}{100}\PY{p}{)}\PY{p}{)}
             \PY{o}{*} \PY{n}{lam\PYZus{}max}\PY{p}{\PYZcb{}}
\end{Verbatim}
\end{tcolorbox}

    Note that we had to transform the data first, since the scale of the
variables impacts the choice of \(\lambda\). We now perform
cross-validation using this sequence of \(\lambda\) values.

    \begin{tcolorbox}[breakable, size=fbox, boxrule=1pt, pad at break*=1mm,colback=cellbackground, colframe=cellborder]
\prompt{In}{incolor}{16}{\boxspacing}
\begin{Verbatim}[commandchars=\\\{\}]
\PY{n}{cv} \PY{o}{=} \PY{n}{KFold}\PY{p}{(}\PY{l+m+mi}{10}\PY{p}{,}
           \PY{n}{shuffle}\PY{o}{=}\PY{k+kc}{True}\PY{p}{,}
           \PY{n}{random\PYZus{}state}\PY{o}{=}\PY{l+m+mi}{1}\PY{p}{)}
\PY{n}{grid} \PY{o}{=} \PY{n}{GridSearchCV}\PY{p}{(}\PY{n}{standard\PYZus{}lasso}\PY{p}{,}
                    \PY{n}{param\PYZus{}grid}\PY{p}{,}
                    \PY{n}{cv}\PY{o}{=}\PY{n}{cv}\PY{p}{,}
                    \PY{n}{scoring}\PY{o}{=}\PY{l+s+s1}{\PYZsq{}}\PY{l+s+s1}{neg\PYZus{}mean\PYZus{}absolute\PYZus{}error}\PY{l+s+s1}{\PYZsq{}}\PY{p}{)}
\PY{n}{grid}\PY{o}{.}\PY{n}{fit}\PY{p}{(}\PY{n}{X\PYZus{}train}\PY{p}{,} \PY{n}{Y\PYZus{}train}\PY{p}{)}\PY{p}{;}
\end{Verbatim}
\end{tcolorbox}

    We extract the lasso model with best cross-validated mean absolute
error, and evaluate its performance on \texttt{X\_test} and
\texttt{Y\_test}, which were not used in cross-validation.

    \begin{tcolorbox}[breakable, size=fbox, boxrule=1pt, pad at break*=1mm,colback=cellbackground, colframe=cellborder]
\prompt{In}{incolor}{17}{\boxspacing}
\begin{Verbatim}[commandchars=\\\{\}]
\PY{n}{trained\PYZus{}lasso} \PY{o}{=} \PY{n}{grid}\PY{o}{.}\PY{n}{best\PYZus{}estimator\PYZus{}}
\PY{n}{Yhat\PYZus{}test} \PY{o}{=} \PY{n}{trained\PYZus{}lasso}\PY{o}{.}\PY{n}{predict}\PY{p}{(}\PY{n}{X\PYZus{}test}\PY{p}{)}
\PY{n}{np}\PY{o}{.}\PY{n}{fabs}\PY{p}{(}\PY{n}{Yhat\PYZus{}test} \PY{o}{\PYZhy{}} \PY{n}{Y\PYZus{}test}\PY{p}{)}\PY{o}{.}\PY{n}{mean}\PY{p}{(}\PY{p}{)}
\end{Verbatim}
\end{tcolorbox}

            \begin{tcolorbox}[breakable, size=fbox, boxrule=.5pt, pad at break*=1mm, opacityfill=0]
\prompt{Out}{outcolor}{17}{\boxspacing}
\begin{Verbatim}[commandchars=\\\{\}]
235.67548374780287
\end{Verbatim}
\end{tcolorbox}
        
    This is similar to the results we got for the linear model fit by least
squares. However, these results can vary a lot for different train/test
splits; we encourage the reader to try a different seed in code block 12
and rerun the subsequent code up to this point.

\hypertarget{specifying-a-network-classes-and-inheritance}{%
\subsubsection{Specifying a Network: Classes and
Inheritance}\label{specifying-a-network-classes-and-inheritance}}

To fit the neural network, we first set up a model structure that
describes the network. Doing so requires us to define new classes
specific to the model we wish to fit. Typically this is done in
\texttt{pytorch} by sub-classing a generic representation of a network,
which is the approach we take here. Although this example is simple, we
will go through the steps in some detail, since it will serve us well
for the more complex examples to follow.

    \begin{tcolorbox}[breakable, size=fbox, boxrule=1pt, pad at break*=1mm,colback=cellbackground, colframe=cellborder]
\prompt{In}{incolor}{18}{\boxspacing}
\begin{Verbatim}[commandchars=\\\{\}]
\PY{k}{class}\PY{+w}{ }\PY{n+nc}{HittersModel}\PY{p}{(}\PY{n}{nn}\PY{o}{.}\PY{n}{Module}\PY{p}{)}\PY{p}{:}

    \PY{k}{def}\PY{+w}{ }\PY{n+nf+fm}{\PYZus{}\PYZus{}init\PYZus{}\PYZus{}}\PY{p}{(}\PY{n+nb+bp}{self}\PY{p}{,} \PY{n}{input\PYZus{}size}\PY{p}{)}\PY{p}{:}
        \PY{n+nb}{super}\PY{p}{(}\PY{n}{HittersModel}\PY{p}{,} \PY{n+nb+bp}{self}\PY{p}{)}\PY{o}{.}\PY{n+nf+fm}{\PYZus{}\PYZus{}init\PYZus{}\PYZus{}}\PY{p}{(}\PY{p}{)}
        \PY{n+nb+bp}{self}\PY{o}{.}\PY{n}{flatten} \PY{o}{=} \PY{n}{nn}\PY{o}{.}\PY{n}{Flatten}\PY{p}{(}\PY{p}{)}
        \PY{n+nb+bp}{self}\PY{o}{.}\PY{n}{sequential} \PY{o}{=} \PY{n}{nn}\PY{o}{.}\PY{n}{Sequential}\PY{p}{(}
            \PY{n}{nn}\PY{o}{.}\PY{n}{Linear}\PY{p}{(}\PY{n}{input\PYZus{}size}\PY{p}{,} \PY{l+m+mi}{50}\PY{p}{)}\PY{p}{,}
            \PY{n}{nn}\PY{o}{.}\PY{n}{ReLU}\PY{p}{(}\PY{p}{)}\PY{p}{,}
            \PY{n}{nn}\PY{o}{.}\PY{n}{Dropout}\PY{p}{(}\PY{l+m+mf}{0.4}\PY{p}{)}\PY{p}{,}
            \PY{n}{nn}\PY{o}{.}\PY{n}{Linear}\PY{p}{(}\PY{l+m+mi}{50}\PY{p}{,} \PY{l+m+mi}{1}\PY{p}{)}\PY{p}{)}

    \PY{k}{def}\PY{+w}{ }\PY{n+nf}{forward}\PY{p}{(}\PY{n+nb+bp}{self}\PY{p}{,} \PY{n}{x}\PY{p}{)}\PY{p}{:}
        \PY{n}{x} \PY{o}{=} \PY{n+nb+bp}{self}\PY{o}{.}\PY{n}{flatten}\PY{p}{(}\PY{n}{x}\PY{p}{)}
        \PY{k}{return} \PY{n}{torch}\PY{o}{.}\PY{n}{flatten}\PY{p}{(}\PY{n+nb+bp}{self}\PY{o}{.}\PY{n}{sequential}\PY{p}{(}\PY{n}{x}\PY{p}{)}\PY{p}{)}
\end{Verbatim}
\end{tcolorbox}

    The \texttt{class} statement identifies the code chunk as a declaration
for a class \texttt{HittersModel} that inherits from the base class
\texttt{nn.Module}. This base class is ubiquitous in \texttt{torch} and
represents the mappings in the neural networks.

Indented beneath the \texttt{class} statement are the methods of this
class: in this case \texttt{\_\_init\_\_} and \texttt{forward}. The
\texttt{\_\_init\_\_} method is called when an instance of the class is
created as in the cell below. In the methods, \texttt{self} always
refers to an instance of the class. In the \texttt{\_\_init\_\_} method,
we have attached two objects to \texttt{self} as attributes:
\texttt{flatten} and \texttt{sequential}. These are used in the
\texttt{forward} method to describe the map that this module implements.

There is one additional line in the \texttt{\_\_init\_\_} method, which
is a call to \texttt{super()}. This function allows subclasses
(i.e.~\texttt{HittersModel}) to access methods of the class they inherit
from. For example, the class \texttt{nn.Module} has its own
\texttt{\_\_init\_\_} method, which is different from the
\texttt{HittersModel.\_\_init\_\_()} method we've written above. Using
\texttt{super()} allows us to call the method of the base class. For
\texttt{torch} models, we will always be making this \texttt{super()}
call as it is necessary for the model to be properly interpreted by
\texttt{torch}.

The object \texttt{nn.Module} has more methods than simply
\texttt{\_\_init\_\_} and \texttt{forward}. These methods are directly
accessible to \texttt{HittersModel} instances because of this
inheritance. One such method we will see shortly is the \texttt{eval()}
method, used to disable dropout for when we want to evaluate the model
on test data.

    \begin{tcolorbox}[breakable, size=fbox, boxrule=1pt, pad at break*=1mm,colback=cellbackground, colframe=cellborder]
\prompt{In}{incolor}{19}{\boxspacing}
\begin{Verbatim}[commandchars=\\\{\}]
\PY{n}{hit\PYZus{}model} \PY{o}{=} \PY{n}{HittersModel}\PY{p}{(}\PY{n}{X}\PY{o}{.}\PY{n}{shape}\PY{p}{[}\PY{l+m+mi}{1}\PY{p}{]}\PY{p}{)}
\end{Verbatim}
\end{tcolorbox}

    The object \texttt{self.sequential} is a composition of four maps. The
first maps the 19 features of \texttt{Hitters} to 50 dimensions,
introducing \(50\times 19+50\) parameters for the weights and
\emph{intercept} of the map (often called the \emph{bias}). This layer
is then mapped to a ReLU layer followed by a 40\% dropout layer, and
finally a linear map down to 1 dimension, again with a bias. The total
number of trainable parameters is therefore
\(50\times 19+50+50+1=1051\).

    The package \texttt{torchinfo} provides a \texttt{summary()} function
that neatly summarizes this information. We specify the size of the
input and see the size of each tensor as it passes through layers of the
network.

    \begin{tcolorbox}[breakable, size=fbox, boxrule=1pt, pad at break*=1mm,colback=cellbackground, colframe=cellborder]
\prompt{In}{incolor}{20}{\boxspacing}
\begin{Verbatim}[commandchars=\\\{\}]
\PY{n}{summary}\PY{p}{(}\PY{n}{hit\PYZus{}model}\PY{p}{,} 
        \PY{n}{input\PYZus{}size}\PY{o}{=}\PY{n}{X\PYZus{}train}\PY{o}{.}\PY{n}{shape}\PY{p}{,}
        \PY{n}{col\PYZus{}names}\PY{o}{=}\PY{p}{[}\PY{l+s+s1}{\PYZsq{}}\PY{l+s+s1}{input\PYZus{}size}\PY{l+s+s1}{\PYZsq{}}\PY{p}{,}
                   \PY{l+s+s1}{\PYZsq{}}\PY{l+s+s1}{output\PYZus{}size}\PY{l+s+s1}{\PYZsq{}}\PY{p}{,}
                   \PY{l+s+s1}{\PYZsq{}}\PY{l+s+s1}{num\PYZus{}params}\PY{l+s+s1}{\PYZsq{}}\PY{p}{]}\PY{p}{)}
\end{Verbatim}
\end{tcolorbox}

            \begin{tcolorbox}[breakable, size=fbox, boxrule=.5pt, pad at break*=1mm, opacityfill=0]
\prompt{Out}{outcolor}{20}{\boxspacing}
\begin{Verbatim}[commandchars=\\\{\}]
================================================================================
===================================
Layer (type:depth-idx)                   Input Shape               Output Shape
Param \#
================================================================================
===================================
HittersModel                             [175, 19]                 [175]
--
├─Flatten: 1-1                           [175, 19]                 [175, 19]
--
├─Sequential: 1-2                        [175, 19]                 [175, 1]
--
│    └─Linear: 2-1                       [175, 19]                 [175, 50]
1,000
│    └─ReLU: 2-2                         [175, 50]                 [175, 50]
--
│    └─Dropout: 2-3                      [175, 50]                 [175, 50]
--
│    └─Linear: 2-4                       [175, 50]                 [175, 1]
51
================================================================================
===================================
Total params: 1,051
Trainable params: 1,051
Non-trainable params: 0
Total mult-adds (Units.MEGABYTES): 0.18
================================================================================
===================================
Input size (MB): 0.01
Forward/backward pass size (MB): 0.07
Params size (MB): 0.00
Estimated Total Size (MB): 0.09
================================================================================
===================================
\end{Verbatim}
\end{tcolorbox}
        
    We have truncated the end of the output slightly, here and in subsequent
uses.

We now need to transform our training data into a form accessible to
\texttt{torch}. The basic datatype in \texttt{torch} is a
\texttt{tensor}, which is very similar to an \texttt{ndarray} from early
chapters. We also note here that \texttt{torch} typically works with
32-bit (\emph{single precision}) rather than 64-bit (\emph{double
precision}) floating point numbers. We therefore convert our data to
\texttt{np.float32} before forming the tensor. The \(X\) and \(Y\)
tensors are then arranged into a \texttt{Dataset} recognized by
\texttt{torch} using \texttt{TensorDataset()}.

    \begin{tcolorbox}[breakable, size=fbox, boxrule=1pt, pad at break*=1mm,colback=cellbackground, colframe=cellborder]
\prompt{In}{incolor}{21}{\boxspacing}
\begin{Verbatim}[commandchars=\\\{\}]
\PY{n}{X\PYZus{}train\PYZus{}t} \PY{o}{=} \PY{n}{torch}\PY{o}{.}\PY{n}{tensor}\PY{p}{(}\PY{n}{X\PYZus{}train}\PY{o}{.}\PY{n}{astype}\PY{p}{(}\PY{n}{np}\PY{o}{.}\PY{n}{float32}\PY{p}{)}\PY{p}{)}
\PY{n}{Y\PYZus{}train\PYZus{}t} \PY{o}{=} \PY{n}{torch}\PY{o}{.}\PY{n}{tensor}\PY{p}{(}\PY{n}{Y\PYZus{}train}\PY{o}{.}\PY{n}{astype}\PY{p}{(}\PY{n}{np}\PY{o}{.}\PY{n}{float32}\PY{p}{)}\PY{p}{)}
\PY{n}{hit\PYZus{}train} \PY{o}{=} \PY{n}{TensorDataset}\PY{p}{(}\PY{n}{X\PYZus{}train\PYZus{}t}\PY{p}{,} \PY{n}{Y\PYZus{}train\PYZus{}t}\PY{p}{)}
\end{Verbatim}
\end{tcolorbox}

    We do the same for the test data.

    \begin{tcolorbox}[breakable, size=fbox, boxrule=1pt, pad at break*=1mm,colback=cellbackground, colframe=cellborder]
\prompt{In}{incolor}{22}{\boxspacing}
\begin{Verbatim}[commandchars=\\\{\}]
\PY{n}{X\PYZus{}test\PYZus{}t} \PY{o}{=} \PY{n}{torch}\PY{o}{.}\PY{n}{tensor}\PY{p}{(}\PY{n}{X\PYZus{}test}\PY{o}{.}\PY{n}{astype}\PY{p}{(}\PY{n}{np}\PY{o}{.}\PY{n}{float32}\PY{p}{)}\PY{p}{)}
\PY{n}{Y\PYZus{}test\PYZus{}t} \PY{o}{=} \PY{n}{torch}\PY{o}{.}\PY{n}{tensor}\PY{p}{(}\PY{n}{Y\PYZus{}test}\PY{o}{.}\PY{n}{astype}\PY{p}{(}\PY{n}{np}\PY{o}{.}\PY{n}{float32}\PY{p}{)}\PY{p}{)}
\PY{n}{hit\PYZus{}test} \PY{o}{=} \PY{n}{TensorDataset}\PY{p}{(}\PY{n}{X\PYZus{}test\PYZus{}t}\PY{p}{,} \PY{n}{Y\PYZus{}test\PYZus{}t}\PY{p}{)}
\end{Verbatim}
\end{tcolorbox}

    Finally, this dataset is passed to a \texttt{DataLoader()} which
ultimately passes data into our network. While this may seem like a lot
of overhead, this structure is helpful for more complex tasks where data
may live on different machines, or where data must be passed to a GPU.
We provide a helper function \texttt{SimpleDataModule()} in
\texttt{ISLP} to make this task easier for standard usage. One of its
arguments is \texttt{num\_workers}, which indicates how many processes
we will use for loading the data. For small data like \texttt{Hitters}
this will have little effect, but it does provide an advantage for the
\texttt{MNIST} and \texttt{CIFAR100} examples below. The \texttt{torch}
package will inspect the process running and determine a maximum number
of workers. \{This depends on the computing hardware and the number of
cores available.\} We've included a function
\texttt{rec\_num\_workers()} to compute this so we know how many workers
might be reasonable (here the max was 16).

    \begin{tcolorbox}[breakable, size=fbox, boxrule=1pt, pad at break*=1mm,colback=cellbackground, colframe=cellborder]
\prompt{In}{incolor}{23}{\boxspacing}
\begin{Verbatim}[commandchars=\\\{\}]
\PY{n}{max\PYZus{}num\PYZus{}workers} \PY{o}{=} \PY{n}{rec\PYZus{}num\PYZus{}workers}\PY{p}{(}\PY{p}{)}
\end{Verbatim}
\end{tcolorbox}

    The general training setup in \texttt{pytorch\_lightning} involves
training, validation and test data. These are each represented by
different data loaders. During each epoch, we run a training step to
learn the model and a validation step to track the error. The test data
is typically used at the end of training to evaluate the model.

In this case, as we had split only into test and training, we'll use the
test data as validation data with the argument
\texttt{validation=hit\_test}. The \texttt{validation} argument can be a
float between 0 and 1, an integer, or a \texttt{Dataset}. If a float
(respectively, integer), it is interpreted as a percentage (respectively
number) of the \emph{training} observations to be used for validation.
If it is a \texttt{Dataset}, it is passed directly to a data loader.

    \begin{tcolorbox}[breakable, size=fbox, boxrule=1pt, pad at break*=1mm,colback=cellbackground, colframe=cellborder]
\prompt{In}{incolor}{24}{\boxspacing}
\begin{Verbatim}[commandchars=\\\{\}]
\PY{n}{hit\PYZus{}dm} \PY{o}{=} \PY{n}{SimpleDataModule}\PY{p}{(}\PY{n}{hit\PYZus{}train}\PY{p}{,}
                          \PY{n}{hit\PYZus{}test}\PY{p}{,}
                          \PY{n}{batch\PYZus{}size}\PY{o}{=}\PY{l+m+mi}{32}\PY{p}{,}
                          \PY{n}{num\PYZus{}workers}\PY{o}{=}\PY{n+nb}{min}\PY{p}{(}\PY{l+m+mi}{4}\PY{p}{,} \PY{n}{max\PYZus{}num\PYZus{}workers}\PY{p}{)}\PY{p}{,}
                          \PY{n}{validation}\PY{o}{=}\PY{n}{hit\PYZus{}test}\PY{p}{)}
\end{Verbatim}
\end{tcolorbox}

    Next we must provide a \texttt{pytorch\_lightning} module that controls
the steps performed during the training process. We provide methods for
our \texttt{SimpleModule()} that simply record the value of the loss
function and any additional metrics at the end of each epoch. These
operations are controlled by the methods
\texttt{SimpleModule.{[}training/test/validation{]}\_step()}, though we
will not be modifying these in our examples.

    \begin{tcolorbox}[breakable, size=fbox, boxrule=1pt, pad at break*=1mm,colback=cellbackground, colframe=cellborder]
\prompt{In}{incolor}{25}{\boxspacing}
\begin{Verbatim}[commandchars=\\\{\}]
\PY{n}{hit\PYZus{}module} \PY{o}{=} \PY{n}{SimpleModule}\PY{o}{.}\PY{n}{regression}\PY{p}{(}\PY{n}{hit\PYZus{}model}\PY{p}{,}
                           \PY{n}{metrics}\PY{o}{=}\PY{p}{\PYZob{}}\PY{l+s+s1}{\PYZsq{}}\PY{l+s+s1}{mae}\PY{l+s+s1}{\PYZsq{}}\PY{p}{:}\PY{n}{MeanAbsoluteError}\PY{p}{(}\PY{p}{)}\PY{p}{\PYZcb{}}\PY{p}{)}
\end{Verbatim}
\end{tcolorbox}

    By using the \texttt{SimpleModule.regression()} method, we indicate that
we will use squared-error loss as in (10.23). We have also asked for
mean absolute error to be tracked as well in the metrics that are
logged.

We log our results via \texttt{CSVLogger()}, which in this case stores
the results in a CSV file within a directory \texttt{logs/hitters}.
After the fitting is complete, this allows us to load the results as a
\texttt{pd.DataFrame()} and visualize them below. There are several ways
to log the results within \texttt{pytorch\_lightning}, though we will
not cover those here in detail.

    \begin{tcolorbox}[breakable, size=fbox, boxrule=1pt, pad at break*=1mm,colback=cellbackground, colframe=cellborder]
\prompt{In}{incolor}{26}{\boxspacing}
\begin{Verbatim}[commandchars=\\\{\}]
\PY{n}{hit\PYZus{}logger} \PY{o}{=} \PY{n}{CSVLogger}\PY{p}{(}\PY{l+s+s1}{\PYZsq{}}\PY{l+s+s1}{logs}\PY{l+s+s1}{\PYZsq{}}\PY{p}{,} \PY{n}{name}\PY{o}{=}\PY{l+s+s1}{\PYZsq{}}\PY{l+s+s1}{hitters}\PY{l+s+s1}{\PYZsq{}}\PY{p}{)}
\end{Verbatim}
\end{tcolorbox}

    Finally we are ready to train our model and log the results. We use the
\texttt{Trainer()} object from \texttt{pytorch\_lightning} to do this
work. The argument \texttt{datamodule=hit\_dm} tells the trainer how
training/validation/test logs are produced, while the first argument
\texttt{hit\_module} specifies the network architecture as well as the
training/validation/test steps. The \texttt{callbacks} argument allows
for several tasks to be carried out at various points while training a
model. Here our \texttt{ErrorTracker()} callback will enable us to
compute validation error while training and, finally, the test error. We
now fit the model for 50 epochs.

    \begin{tcolorbox}[breakable, size=fbox, boxrule=1pt, pad at break*=1mm,colback=cellbackground, colframe=cellborder]
\prompt{In}{incolor}{27}{\boxspacing}
\begin{Verbatim}[commandchars=\\\{\}]
\PY{n}{hit\PYZus{}trainer} \PY{o}{=} \PY{n}{Trainer}\PY{p}{(}\PY{n}{deterministic}\PY{o}{=}\PY{k+kc}{True}\PY{p}{,}
                      \PY{n}{max\PYZus{}epochs}\PY{o}{=}\PY{l+m+mi}{50}\PY{p}{,}
                      \PY{n}{log\PYZus{}every\PYZus{}n\PYZus{}steps}\PY{o}{=}\PY{l+m+mi}{5}\PY{p}{,}
                      \PY{n}{logger}\PY{o}{=}\PY{n}{hit\PYZus{}logger}\PY{p}{,}
                      \PY{n}{callbacks}\PY{o}{=}\PY{p}{[}\PY{n}{ErrorTracker}\PY{p}{(}\PY{p}{)}\PY{p}{]}\PY{p}{)}
\PY{n}{hit\PYZus{}trainer}\PY{o}{.}\PY{n}{fit}\PY{p}{(}\PY{n}{hit\PYZus{}module}\PY{p}{,} \PY{n}{datamodule}\PY{o}{=}\PY{n}{hit\PYZus{}dm}\PY{p}{)}
\end{Verbatim}
\end{tcolorbox}

    \begin{Verbatim}[commandchars=\\\{\}]
GPU available: True (mps), used: True
TPU available: False, using: 0 TPU cores
IPU available: False, using: 0 IPUs
HPU available: False, using: 0 HPUs

  | Name  | Type         | Params
---------------------------------------
0 | model | HittersModel | 1.1 K
1 | loss  | MSELoss      | 0
---------------------------------------
1.1 K     Trainable params
0         Non-trainable params
1.1 K     Total params
0.004     Total estimated model params size (MB)
    \end{Verbatim}

    
    \begin{Verbatim}[commandchars=\\\{\}]
Sanity Checking: |                                                                                            …
    \end{Verbatim}

    
    
    \begin{Verbatim}[commandchars=\\\{\}]
Training: |                                                                                                   …
    \end{Verbatim}

    
    
    \begin{Verbatim}[commandchars=\\\{\}]
Validation: |                                                                                                 …
    \end{Verbatim}

    
    
    \begin{Verbatim}[commandchars=\\\{\}]
Validation: |                                                                                                 …
    \end{Verbatim}

    
    
    \begin{Verbatim}[commandchars=\\\{\}]
Validation: |                                                                                                 …
    \end{Verbatim}

    
    
    \begin{Verbatim}[commandchars=\\\{\}]
Validation: |                                                                                                 …
    \end{Verbatim}

    
    
    \begin{Verbatim}[commandchars=\\\{\}]
Validation: |                                                                                                 …
    \end{Verbatim}

    
    
    \begin{Verbatim}[commandchars=\\\{\}]
Validation: |                                                                                                 …
    \end{Verbatim}

    
    
    \begin{Verbatim}[commandchars=\\\{\}]
Validation: |                                                                                                 …
    \end{Verbatim}

    
    
    \begin{Verbatim}[commandchars=\\\{\}]
Validation: |                                                                                                 …
    \end{Verbatim}

    
    
    \begin{Verbatim}[commandchars=\\\{\}]
Validation: |                                                                                                 …
    \end{Verbatim}

    
    
    \begin{Verbatim}[commandchars=\\\{\}]
Validation: |                                                                                                 …
    \end{Verbatim}

    
    
    \begin{Verbatim}[commandchars=\\\{\}]
Validation: |                                                                                                 …
    \end{Verbatim}

    
    
    \begin{Verbatim}[commandchars=\\\{\}]
Validation: |                                                                                                 …
    \end{Verbatim}

    
    
    \begin{Verbatim}[commandchars=\\\{\}]
Validation: |                                                                                                 …
    \end{Verbatim}

    
    
    \begin{Verbatim}[commandchars=\\\{\}]
Validation: |                                                                                                 …
    \end{Verbatim}

    
    
    \begin{Verbatim}[commandchars=\\\{\}]
Validation: |                                                                                                 …
    \end{Verbatim}

    
    
    \begin{Verbatim}[commandchars=\\\{\}]
Validation: |                                                                                                 …
    \end{Verbatim}

    
    
    \begin{Verbatim}[commandchars=\\\{\}]
Validation: |                                                                                                 …
    \end{Verbatim}

    
    
    \begin{Verbatim}[commandchars=\\\{\}]
Validation: |                                                                                                 …
    \end{Verbatim}

    
    
    \begin{Verbatim}[commandchars=\\\{\}]
Validation: |                                                                                                 …
    \end{Verbatim}

    
    
    \begin{Verbatim}[commandchars=\\\{\}]
Validation: |                                                                                                 …
    \end{Verbatim}

    
    
    \begin{Verbatim}[commandchars=\\\{\}]
Validation: |                                                                                                 …
    \end{Verbatim}

    
    
    \begin{Verbatim}[commandchars=\\\{\}]
Validation: |                                                                                                 …
    \end{Verbatim}

    
    
    \begin{Verbatim}[commandchars=\\\{\}]
Validation: |                                                                                                 …
    \end{Verbatim}

    
    
    \begin{Verbatim}[commandchars=\\\{\}]
Validation: |                                                                                                 …
    \end{Verbatim}

    
    
    \begin{Verbatim}[commandchars=\\\{\}]
Validation: |                                                                                                 …
    \end{Verbatim}

    
    
    \begin{Verbatim}[commandchars=\\\{\}]
Validation: |                                                                                                 …
    \end{Verbatim}

    
    
    \begin{Verbatim}[commandchars=\\\{\}]
Validation: |                                                                                                 …
    \end{Verbatim}

    
    
    \begin{Verbatim}[commandchars=\\\{\}]
Validation: |                                                                                                 …
    \end{Verbatim}

    
    
    \begin{Verbatim}[commandchars=\\\{\}]
Validation: |                                                                                                 …
    \end{Verbatim}

    
    
    \begin{Verbatim}[commandchars=\\\{\}]
Validation: |                                                                                                 …
    \end{Verbatim}

    
    
    \begin{Verbatim}[commandchars=\\\{\}]
Validation: |                                                                                                 …
    \end{Verbatim}

    
    
    \begin{Verbatim}[commandchars=\\\{\}]
Validation: |                                                                                                 …
    \end{Verbatim}

    
    
    \begin{Verbatim}[commandchars=\\\{\}]
Validation: |                                                                                                 …
    \end{Verbatim}

    
    
    \begin{Verbatim}[commandchars=\\\{\}]
Validation: |                                                                                                 …
    \end{Verbatim}

    
    
    \begin{Verbatim}[commandchars=\\\{\}]
Validation: |                                                                                                 …
    \end{Verbatim}

    
    
    \begin{Verbatim}[commandchars=\\\{\}]
Validation: |                                                                                                 …
    \end{Verbatim}

    
    
    \begin{Verbatim}[commandchars=\\\{\}]
Validation: |                                                                                                 …
    \end{Verbatim}

    
    
    \begin{Verbatim}[commandchars=\\\{\}]
Validation: |                                                                                                 …
    \end{Verbatim}

    
    
    \begin{Verbatim}[commandchars=\\\{\}]
Validation: |                                                                                                 …
    \end{Verbatim}

    
    
    \begin{Verbatim}[commandchars=\\\{\}]
Validation: |                                                                                                 …
    \end{Verbatim}

    
    
    \begin{Verbatim}[commandchars=\\\{\}]
Validation: |                                                                                                 …
    \end{Verbatim}

    
    
    \begin{Verbatim}[commandchars=\\\{\}]
Validation: |                                                                                                 …
    \end{Verbatim}

    
    
    \begin{Verbatim}[commandchars=\\\{\}]
Validation: |                                                                                                 …
    \end{Verbatim}

    
    
    \begin{Verbatim}[commandchars=\\\{\}]
Validation: |                                                                                                 …
    \end{Verbatim}

    
    
    \begin{Verbatim}[commandchars=\\\{\}]
Validation: |                                                                                                 …
    \end{Verbatim}

    
    
    \begin{Verbatim}[commandchars=\\\{\}]
Validation: |                                                                                                 …
    \end{Verbatim}

    
    
    \begin{Verbatim}[commandchars=\\\{\}]
Validation: |                                                                                                 …
    \end{Verbatim}

    
    
    \begin{Verbatim}[commandchars=\\\{\}]
Validation: |                                                                                                 …
    \end{Verbatim}

    
    
    \begin{Verbatim}[commandchars=\\\{\}]
Validation: |                                                                                                 …
    \end{Verbatim}

    
    
    \begin{Verbatim}[commandchars=\\\{\}]
Validation: |                                                                                                 …
    \end{Verbatim}

    
    \begin{Verbatim}[commandchars=\\\{\}]
`Trainer.fit` stopped: `max\_epochs=50` reached.
    \end{Verbatim}

    At each step of SGD, the algorithm randomly selects 32 training
observations for the computation of the gradient. Recall from Section
10.7 that an epoch amounts to the number of SGD steps required to
process \(n\) observations. Since the training set has \(n=175\), and we
specified a \texttt{batch\_size} of 32 in the construction of
\texttt{hit\_dm}, an epoch is \(175/32=5.5\) SGD steps.

After having fit the model, we can evaluate performance on our test data
using the \texttt{test()} method of our trainer.

    \begin{tcolorbox}[breakable, size=fbox, boxrule=1pt, pad at break*=1mm,colback=cellbackground, colframe=cellborder]
\prompt{In}{incolor}{28}{\boxspacing}
\begin{Verbatim}[commandchars=\\\{\}]
\PY{n}{hit\PYZus{}trainer}\PY{o}{.}\PY{n}{test}\PY{p}{(}\PY{n}{hit\PYZus{}module}\PY{p}{,} \PY{n}{datamodule}\PY{o}{=}\PY{n}{hit\PYZus{}dm}\PY{p}{)}
\end{Verbatim}
\end{tcolorbox}

    
    \begin{Verbatim}[commandchars=\\\{\}]
Testing: |                                                                                                    …
    \end{Verbatim}

    
    \begin{Verbatim}[commandchars=\\\{\}]
────────────────────────────────────────────────────────────────────────────────
────────────────────────────────────────
       Test metric             DataLoader 0
────────────────────────────────────────────────────────────────────────────────
────────────────────────────────────────
        test\_loss              107904.640625
        test\_mae             221.8314971923828
────────────────────────────────────────────────────────────────────────────────
────────────────────────────────────────
    \end{Verbatim}

            \begin{tcolorbox}[breakable, size=fbox, boxrule=.5pt, pad at break*=1mm, opacityfill=0]
\prompt{Out}{outcolor}{28}{\boxspacing}
\begin{Verbatim}[commandchars=\\\{\}]
[\{'test\_loss': 107904.640625, 'test\_mae': 221.8314971923828\}]
\end{Verbatim}
\end{tcolorbox}
        
    The results of the fit have been logged into a CSV file. We can find the
results specific to this run in the
\texttt{experiment.metrics\_file\_path} attribute of our logger. Note
that each time the model is fit, the logger will output results into a
new subdirectory of our directory \texttt{logs/hitters}.

We now create a plot of the MAE (mean absolute error) as a function of
the number of epochs. First we retrieve the logged summaries.

    \begin{tcolorbox}[breakable, size=fbox, boxrule=1pt, pad at break*=1mm,colback=cellbackground, colframe=cellborder]
\prompt{In}{incolor}{29}{\boxspacing}
\begin{Verbatim}[commandchars=\\\{\}]
\PY{n}{hit\PYZus{}results} \PY{o}{=} \PY{n}{pd}\PY{o}{.}\PY{n}{read\PYZus{}csv}\PY{p}{(}\PY{n}{hit\PYZus{}logger}\PY{o}{.}\PY{n}{experiment}\PY{o}{.}\PY{n}{metrics\PYZus{}file\PYZus{}path}\PY{p}{)}
\end{Verbatim}
\end{tcolorbox}

    Since we will produce similar plots in later examples, we write a simple
generic function to produce this plot.

    \begin{tcolorbox}[breakable, size=fbox, boxrule=1pt, pad at break*=1mm,colback=cellbackground, colframe=cellborder]
\prompt{In}{incolor}{30}{\boxspacing}
\begin{Verbatim}[commandchars=\\\{\}]
\PY{k}{def}\PY{+w}{ }\PY{n+nf}{summary\PYZus{}plot}\PY{p}{(}\PY{n}{results}\PY{p}{,}
                 \PY{n}{ax}\PY{p}{,}
                 \PY{n}{col}\PY{o}{=}\PY{l+s+s1}{\PYZsq{}}\PY{l+s+s1}{loss}\PY{l+s+s1}{\PYZsq{}}\PY{p}{,}
                 \PY{n}{valid\PYZus{}legend}\PY{o}{=}\PY{l+s+s1}{\PYZsq{}}\PY{l+s+s1}{Validation}\PY{l+s+s1}{\PYZsq{}}\PY{p}{,}
                 \PY{n}{training\PYZus{}legend}\PY{o}{=}\PY{l+s+s1}{\PYZsq{}}\PY{l+s+s1}{Training}\PY{l+s+s1}{\PYZsq{}}\PY{p}{,}
                 \PY{n}{ylabel}\PY{o}{=}\PY{l+s+s1}{\PYZsq{}}\PY{l+s+s1}{Loss}\PY{l+s+s1}{\PYZsq{}}\PY{p}{,}
                 \PY{n}{fontsize}\PY{o}{=}\PY{l+m+mi}{20}\PY{p}{)}\PY{p}{:}
    \PY{k}{for} \PY{p}{(}\PY{n}{column}\PY{p}{,}
         \PY{n}{color}\PY{p}{,}
         \PY{n}{label}\PY{p}{)} \PY{o+ow}{in} \PY{n+nb}{zip}\PY{p}{(}\PY{p}{[}\PY{l+s+sa}{f}\PY{l+s+s1}{\PYZsq{}}\PY{l+s+s1}{train\PYZus{}}\PY{l+s+si}{\PYZob{}}\PY{n}{col}\PY{l+s+si}{\PYZcb{}}\PY{l+s+s1}{\PYZus{}epoch}\PY{l+s+s1}{\PYZsq{}}\PY{p}{,}
                        \PY{l+s+sa}{f}\PY{l+s+s1}{\PYZsq{}}\PY{l+s+s1}{valid\PYZus{}}\PY{l+s+si}{\PYZob{}}\PY{n}{col}\PY{l+s+si}{\PYZcb{}}\PY{l+s+s1}{\PYZsq{}}\PY{p}{]}\PY{p}{,}
                       \PY{p}{[}\PY{l+s+s1}{\PYZsq{}}\PY{l+s+s1}{black}\PY{l+s+s1}{\PYZsq{}}\PY{p}{,}
                        \PY{l+s+s1}{\PYZsq{}}\PY{l+s+s1}{red}\PY{l+s+s1}{\PYZsq{}}\PY{p}{]}\PY{p}{,}
                       \PY{p}{[}\PY{n}{training\PYZus{}legend}\PY{p}{,}
                        \PY{n}{valid\PYZus{}legend}\PY{p}{]}\PY{p}{)}\PY{p}{:}
        \PY{n}{results}\PY{o}{.}\PY{n}{plot}\PY{p}{(}\PY{n}{x}\PY{o}{=}\PY{l+s+s1}{\PYZsq{}}\PY{l+s+s1}{epoch}\PY{l+s+s1}{\PYZsq{}}\PY{p}{,}
                     \PY{n}{y}\PY{o}{=}\PY{n}{column}\PY{p}{,}
                     \PY{n}{label}\PY{o}{=}\PY{n}{label}\PY{p}{,}
                     \PY{n}{marker}\PY{o}{=}\PY{l+s+s1}{\PYZsq{}}\PY{l+s+s1}{o}\PY{l+s+s1}{\PYZsq{}}\PY{p}{,}
                     \PY{n}{color}\PY{o}{=}\PY{n}{color}\PY{p}{,}
                     \PY{n}{ax}\PY{o}{=}\PY{n}{ax}\PY{p}{)}
    \PY{n}{ax}\PY{o}{.}\PY{n}{set\PYZus{}xlabel}\PY{p}{(}\PY{l+s+s1}{\PYZsq{}}\PY{l+s+s1}{Epoch}\PY{l+s+s1}{\PYZsq{}}\PY{p}{)}
    \PY{n}{ax}\PY{o}{.}\PY{n}{set\PYZus{}ylabel}\PY{p}{(}\PY{n}{ylabel}\PY{p}{)}
    \PY{k}{return} \PY{n}{ax}
\end{Verbatim}
\end{tcolorbox}

    We now set up our axes, and use our function to produce the MAE plot.

    \begin{tcolorbox}[breakable, size=fbox, boxrule=1pt, pad at break*=1mm,colback=cellbackground, colframe=cellborder]
\prompt{In}{incolor}{31}{\boxspacing}
\begin{Verbatim}[commandchars=\\\{\}]
\PY{n}{fig}\PY{p}{,} \PY{n}{ax} \PY{o}{=} \PY{n}{subplots}\PY{p}{(}\PY{l+m+mi}{1}\PY{p}{,} \PY{l+m+mi}{1}\PY{p}{,} \PY{n}{figsize}\PY{o}{=}\PY{p}{(}\PY{l+m+mi}{6}\PY{p}{,} \PY{l+m+mi}{6}\PY{p}{)}\PY{p}{)}
\PY{n}{ax} \PY{o}{=} \PY{n}{summary\PYZus{}plot}\PY{p}{(}\PY{n}{hit\PYZus{}results}\PY{p}{,}
                  \PY{n}{ax}\PY{p}{,}
                  \PY{n}{col}\PY{o}{=}\PY{l+s+s1}{\PYZsq{}}\PY{l+s+s1}{mae}\PY{l+s+s1}{\PYZsq{}}\PY{p}{,}
                  \PY{n}{ylabel}\PY{o}{=}\PY{l+s+s1}{\PYZsq{}}\PY{l+s+s1}{MAE}\PY{l+s+s1}{\PYZsq{}}\PY{p}{,}
                  \PY{n}{valid\PYZus{}legend}\PY{o}{=}\PY{l+s+s1}{\PYZsq{}}\PY{l+s+s1}{Validation (=Test)}\PY{l+s+s1}{\PYZsq{}}\PY{p}{)}
\PY{n}{ax}\PY{o}{.}\PY{n}{set\PYZus{}ylim}\PY{p}{(}\PY{p}{[}\PY{l+m+mi}{0}\PY{p}{,} \PY{l+m+mi}{400}\PY{p}{]}\PY{p}{)}
\PY{n}{ax}\PY{o}{.}\PY{n}{set\PYZus{}xticks}\PY{p}{(}\PY{n}{np}\PY{o}{.}\PY{n}{linspace}\PY{p}{(}\PY{l+m+mi}{0}\PY{p}{,} \PY{l+m+mi}{50}\PY{p}{,} \PY{l+m+mi}{11}\PY{p}{)}\PY{o}{.}\PY{n}{astype}\PY{p}{(}\PY{n+nb}{int}\PY{p}{)}\PY{p}{)}\PY{p}{;}
\end{Verbatim}
\end{tcolorbox}

    \begin{center}
    \adjustimage{max size={0.9\linewidth}{0.9\paperheight}}{artifacts/latex/Ch10-deeplearning-lab_files/artifacts/latex/Ch10-deeplearning-lab_64_0.png}
    \end{center}
    { \hspace*{\fill} \\}
    
    We can predict directly from the final model, and evaluate its
performance on the test data. Before fitting, we call the
\texttt{eval()} method of \texttt{hit\_model}. This tells \texttt{torch}
to effectively consider this model to be fitted, so that we can use it
to predict on new data. For our model here, the biggest change is that
the dropout layers will be turned off, i.e.~no weights will be randomly
dropped in predicting on new data.

    \begin{tcolorbox}[breakable, size=fbox, boxrule=1pt, pad at break*=1mm,colback=cellbackground, colframe=cellborder]
\prompt{In}{incolor}{32}{\boxspacing}
\begin{Verbatim}[commandchars=\\\{\}]
\PY{n}{hit\PYZus{}model}\PY{o}{.}\PY{n}{eval}\PY{p}{(}\PY{p}{)} 
\PY{n}{preds} \PY{o}{=} \PY{n}{hit\PYZus{}module}\PY{p}{(}\PY{n}{X\PYZus{}test\PYZus{}t}\PY{p}{)}
\PY{n}{torch}\PY{o}{.}\PY{n}{abs}\PY{p}{(}\PY{n}{Y\PYZus{}test\PYZus{}t} \PY{o}{\PYZhy{}} \PY{n}{preds}\PY{p}{)}\PY{o}{.}\PY{n}{mean}\PY{p}{(}\PY{p}{)}
\end{Verbatim}
\end{tcolorbox}

            \begin{tcolorbox}[breakable, size=fbox, boxrule=.5pt, pad at break*=1mm, opacityfill=0]
\prompt{Out}{outcolor}{32}{\boxspacing}
\begin{Verbatim}[commandchars=\\\{\}]
tensor(221.8315, grad\_fn=<MeanBackward0>)
\end{Verbatim}
\end{tcolorbox}
        
    

    \hypertarget{cleanup}{%
\subsubsection{Cleanup}\label{cleanup}}

In setting up our data module, we had initiated several worker processes
that will remain running. We delete all references to the torch objects
to ensure these processes will be killed.

    \begin{tcolorbox}[breakable, size=fbox, boxrule=1pt, pad at break*=1mm,colback=cellbackground, colframe=cellborder]
\prompt{In}{incolor}{33}{\boxspacing}
\begin{Verbatim}[commandchars=\\\{\}]
\PY{k}{del}\PY{p}{(}\PY{n}{Hitters}\PY{p}{,}
    \PY{n}{hit\PYZus{}model}\PY{p}{,} \PY{n}{hit\PYZus{}dm}\PY{p}{,}
    \PY{n}{hit\PYZus{}logger}\PY{p}{,}
    \PY{n}{hit\PYZus{}test}\PY{p}{,} \PY{n}{hit\PYZus{}train}\PY{p}{,}
    \PY{n}{X}\PY{p}{,} \PY{n}{Y}\PY{p}{,}
    \PY{n}{X\PYZus{}test}\PY{p}{,} \PY{n}{X\PYZus{}train}\PY{p}{,}
    \PY{n}{Y\PYZus{}test}\PY{p}{,} \PY{n}{Y\PYZus{}train}\PY{p}{,}
    \PY{n}{X\PYZus{}test\PYZus{}t}\PY{p}{,} \PY{n}{Y\PYZus{}test\PYZus{}t}\PY{p}{,}
    \PY{n}{hit\PYZus{}trainer}\PY{p}{,} \PY{n}{hit\PYZus{}module}\PY{p}{)}
\end{Verbatim}
\end{tcolorbox}

    \hypertarget{multilayer-network-on-the-mnist-digit-data}{%
\subsection{Multilayer Network on the MNIST Digit
Data}\label{multilayer-network-on-the-mnist-digit-data}}

The \texttt{torchvision} package comes with a number of example
datasets, including the \texttt{MNIST} digit data. Our first step is to
retrieve the training and test data sets; the \texttt{MNIST()} function
within \texttt{torchvision.datasets} is provided for this purpose. The
data will be downloaded the first time this function is executed, and
stored in the directory \texttt{data/MNIST}.

    \begin{tcolorbox}[breakable, size=fbox, boxrule=1pt, pad at break*=1mm,colback=cellbackground, colframe=cellborder]
\prompt{In}{incolor}{34}{\boxspacing}
\begin{Verbatim}[commandchars=\\\{\}]
\PY{p}{(}\PY{n}{mnist\PYZus{}train}\PY{p}{,} 
 \PY{n}{mnist\PYZus{}test}\PY{p}{)} \PY{o}{=} \PY{p}{[}\PY{n}{MNIST}\PY{p}{(}\PY{n}{root}\PY{o}{=}\PY{l+s+s1}{\PYZsq{}}\PY{l+s+s1}{data}\PY{l+s+s1}{\PYZsq{}}\PY{p}{,}
                      \PY{n}{train}\PY{o}{=}\PY{n}{train}\PY{p}{,}
                      \PY{n}{download}\PY{o}{=}\PY{k+kc}{True}\PY{p}{,}
                      \PY{n}{transform}\PY{o}{=}\PY{n}{ToTensor}\PY{p}{(}\PY{p}{)}\PY{p}{)}
                \PY{k}{for} \PY{n}{train} \PY{o+ow}{in} \PY{p}{[}\PY{k+kc}{True}\PY{p}{,} \PY{k+kc}{False}\PY{p}{]}\PY{p}{]}
\PY{n}{mnist\PYZus{}train}
\end{Verbatim}
\end{tcolorbox}

            \begin{tcolorbox}[breakable, size=fbox, boxrule=.5pt, pad at break*=1mm, opacityfill=0]
\prompt{Out}{outcolor}{34}{\boxspacing}
\begin{Verbatim}[commandchars=\\\{\}]
Dataset MNIST
    Number of datapoints: 60000
    Root location: data
    Split: Train
    StandardTransform
Transform: ToTensor()
\end{Verbatim}
\end{tcolorbox}
        
    There are 60,000 images in the training data and 10,000 in the test
data. The images are \(28\times 28\), and stored as a matrix of pixels.
We need to transform each one into a vector.

Neural networks are somewhat sensitive to the scale of the inputs, much
as ridge and lasso regularization are affected by scaling. Here the
inputs are eight-bit grayscale values between 0 and 255, so we rescale
to the unit interval. \{Note: eight bits means \(2^8\), which equals
256. Since the convention is to start at \(0\), the possible values
range from \(0\) to \(255\).\} This transformation, along with some
reordering of the axes, is performed by the \texttt{ToTensor()}
transform from the \texttt{torchvision.transforms} package.

As in our \texttt{Hitters} example, we form a data module from the
training and test datasets, setting aside 20\% of the training images
for validation.

    \begin{tcolorbox}[breakable, size=fbox, boxrule=1pt, pad at break*=1mm,colback=cellbackground, colframe=cellborder]
\prompt{In}{incolor}{35}{\boxspacing}
\begin{Verbatim}[commandchars=\\\{\}]
\PY{n}{mnist\PYZus{}dm} \PY{o}{=} \PY{n}{SimpleDataModule}\PY{p}{(}\PY{n}{mnist\PYZus{}train}\PY{p}{,}
                            \PY{n}{mnist\PYZus{}test}\PY{p}{,}
                            \PY{n}{validation}\PY{o}{=}\PY{l+m+mf}{0.2}\PY{p}{,}
                            \PY{n}{num\PYZus{}workers}\PY{o}{=}\PY{n}{max\PYZus{}num\PYZus{}workers}\PY{p}{,}
                            \PY{n}{batch\PYZus{}size}\PY{o}{=}\PY{l+m+mi}{256}\PY{p}{)}
\end{Verbatim}
\end{tcolorbox}

    Let's take a look at the data that will get fed into our network. We
loop through the first few chunks of the test dataset, breaking after 2
batches:

    \begin{tcolorbox}[breakable, size=fbox, boxrule=1pt, pad at break*=1mm,colback=cellbackground, colframe=cellborder]
\prompt{In}{incolor}{36}{\boxspacing}
\begin{Verbatim}[commandchars=\\\{\}]
\PY{k}{for} \PY{n}{idx}\PY{p}{,} \PY{p}{(}\PY{n}{X\PYZus{}} \PY{p}{,}\PY{n}{Y\PYZus{}}\PY{p}{)} \PY{o+ow}{in} \PY{n+nb}{enumerate}\PY{p}{(}\PY{n}{mnist\PYZus{}dm}\PY{o}{.}\PY{n}{train\PYZus{}dataloader}\PY{p}{(}\PY{p}{)}\PY{p}{)}\PY{p}{:}
    \PY{n+nb}{print}\PY{p}{(}\PY{l+s+s1}{\PYZsq{}}\PY{l+s+s1}{X: }\PY{l+s+s1}{\PYZsq{}}\PY{p}{,} \PY{n}{X\PYZus{}}\PY{o}{.}\PY{n}{shape}\PY{p}{)}
    \PY{n+nb}{print}\PY{p}{(}\PY{l+s+s1}{\PYZsq{}}\PY{l+s+s1}{Y: }\PY{l+s+s1}{\PYZsq{}}\PY{p}{,} \PY{n}{Y\PYZus{}}\PY{o}{.}\PY{n}{shape}\PY{p}{)}
    \PY{k}{if} \PY{n}{idx} \PY{o}{\PYZgt{}}\PY{o}{=} \PY{l+m+mi}{1}\PY{p}{:}
        \PY{k}{break}
\end{Verbatim}
\end{tcolorbox}

    \begin{Verbatim}[commandchars=\\\{\}]
X:  torch.Size([256, 1, 28, 28])
Y:  torch.Size([256])
X:  torch.Size([256, 1, 28, 28])
Y:  torch.Size([256])
    \end{Verbatim}

    We see that the \(X\) for each batch consists of 256 images of size
\texttt{1x28x28}. Here the \texttt{1} indicates a single channel
(greyscale). For RGB images such as \texttt{CIFAR100} below, we will see
that the \texttt{1} in the size will be replaced by \texttt{3} for the
three RGB channels.

Now we are ready to specify our neural network.

    \begin{tcolorbox}[breakable, size=fbox, boxrule=1pt, pad at break*=1mm,colback=cellbackground, colframe=cellborder]
\prompt{In}{incolor}{37}{\boxspacing}
\begin{Verbatim}[commandchars=\\\{\}]
\PY{k}{class}\PY{+w}{ }\PY{n+nc}{MNISTModel}\PY{p}{(}\PY{n}{nn}\PY{o}{.}\PY{n}{Module}\PY{p}{)}\PY{p}{:}
    \PY{k}{def}\PY{+w}{ }\PY{n+nf+fm}{\PYZus{}\PYZus{}init\PYZus{}\PYZus{}}\PY{p}{(}\PY{n+nb+bp}{self}\PY{p}{)}\PY{p}{:}
        \PY{n+nb}{super}\PY{p}{(}\PY{n}{MNISTModel}\PY{p}{,} \PY{n+nb+bp}{self}\PY{p}{)}\PY{o}{.}\PY{n+nf+fm}{\PYZus{}\PYZus{}init\PYZus{}\PYZus{}}\PY{p}{(}\PY{p}{)}
        \PY{n+nb+bp}{self}\PY{o}{.}\PY{n}{layer1} \PY{o}{=} \PY{n}{nn}\PY{o}{.}\PY{n}{Sequential}\PY{p}{(}
            \PY{n}{nn}\PY{o}{.}\PY{n}{Flatten}\PY{p}{(}\PY{p}{)}\PY{p}{,}
            \PY{n}{nn}\PY{o}{.}\PY{n}{Linear}\PY{p}{(}\PY{l+m+mi}{28}\PY{o}{*}\PY{l+m+mi}{28}\PY{p}{,} \PY{l+m+mi}{256}\PY{p}{)}\PY{p}{,}
            \PY{n}{nn}\PY{o}{.}\PY{n}{ReLU}\PY{p}{(}\PY{p}{)}\PY{p}{,}
            \PY{n}{nn}\PY{o}{.}\PY{n}{Dropout}\PY{p}{(}\PY{l+m+mf}{0.4}\PY{p}{)}\PY{p}{)}
        \PY{n+nb+bp}{self}\PY{o}{.}\PY{n}{layer2} \PY{o}{=} \PY{n}{nn}\PY{o}{.}\PY{n}{Sequential}\PY{p}{(}
            \PY{n}{nn}\PY{o}{.}\PY{n}{Linear}\PY{p}{(}\PY{l+m+mi}{256}\PY{p}{,} \PY{l+m+mi}{128}\PY{p}{)}\PY{p}{,}
            \PY{n}{nn}\PY{o}{.}\PY{n}{ReLU}\PY{p}{(}\PY{p}{)}\PY{p}{,}
            \PY{n}{nn}\PY{o}{.}\PY{n}{Dropout}\PY{p}{(}\PY{l+m+mf}{0.3}\PY{p}{)}\PY{p}{)}
        \PY{n+nb+bp}{self}\PY{o}{.}\PY{n}{\PYZus{}forward} \PY{o}{=} \PY{n}{nn}\PY{o}{.}\PY{n}{Sequential}\PY{p}{(}
            \PY{n+nb+bp}{self}\PY{o}{.}\PY{n}{layer1}\PY{p}{,}
            \PY{n+nb+bp}{self}\PY{o}{.}\PY{n}{layer2}\PY{p}{,}
            \PY{n}{nn}\PY{o}{.}\PY{n}{Linear}\PY{p}{(}\PY{l+m+mi}{128}\PY{p}{,} \PY{l+m+mi}{10}\PY{p}{)}\PY{p}{)}
    \PY{k}{def}\PY{+w}{ }\PY{n+nf}{forward}\PY{p}{(}\PY{n+nb+bp}{self}\PY{p}{,} \PY{n}{x}\PY{p}{)}\PY{p}{:}
        \PY{k}{return} \PY{n+nb+bp}{self}\PY{o}{.}\PY{n}{\PYZus{}forward}\PY{p}{(}\PY{n}{x}\PY{p}{)}
\end{Verbatim}
\end{tcolorbox}

    We see that in the first layer, each \texttt{1x28x28} image is
flattened, then mapped to 256 dimensions where we apply a ReLU
activation with 40\% dropout. A second layer maps the first layer's
output down to 128 dimensions, applying a ReLU activation with 30\%
dropout. Finally, the 128 dimensions are mapped down to 10, the number
of classes in the \texttt{MNIST} data.

    \begin{tcolorbox}[breakable, size=fbox, boxrule=1pt, pad at break*=1mm,colback=cellbackground, colframe=cellborder]
\prompt{In}{incolor}{38}{\boxspacing}
\begin{Verbatim}[commandchars=\\\{\}]
\PY{n}{mnist\PYZus{}model} \PY{o}{=} \PY{n}{MNISTModel}\PY{p}{(}\PY{p}{)}
\end{Verbatim}
\end{tcolorbox}

    We can check that the model produces output of expected size based on
our existing batch \texttt{X\_} above.

    \begin{tcolorbox}[breakable, size=fbox, boxrule=1pt, pad at break*=1mm,colback=cellbackground, colframe=cellborder]
\prompt{In}{incolor}{39}{\boxspacing}
\begin{Verbatim}[commandchars=\\\{\}]
\PY{n}{mnist\PYZus{}model}\PY{p}{(}\PY{n}{X\PYZus{}}\PY{p}{)}\PY{o}{.}\PY{n}{size}\PY{p}{(}\PY{p}{)}
\end{Verbatim}
\end{tcolorbox}

            \begin{tcolorbox}[breakable, size=fbox, boxrule=.5pt, pad at break*=1mm, opacityfill=0]
\prompt{Out}{outcolor}{39}{\boxspacing}
\begin{Verbatim}[commandchars=\\\{\}]
torch.Size([256, 10])
\end{Verbatim}
\end{tcolorbox}
        
    Let's take a look at the summary of the model. Instead of an
\texttt{input\_size} we can pass a tensor of correct shape. In this
case, we pass through the final batched \texttt{X\_} from above.

    \begin{tcolorbox}[breakable, size=fbox, boxrule=1pt, pad at break*=1mm,colback=cellbackground, colframe=cellborder]
\prompt{In}{incolor}{40}{\boxspacing}
\begin{Verbatim}[commandchars=\\\{\}]
\PY{n}{summary}\PY{p}{(}\PY{n}{mnist\PYZus{}model}\PY{p}{,}
        \PY{n}{input\PYZus{}data}\PY{o}{=}\PY{n}{X\PYZus{}}\PY{p}{,}
        \PY{n}{col\PYZus{}names}\PY{o}{=}\PY{p}{[}\PY{l+s+s1}{\PYZsq{}}\PY{l+s+s1}{input\PYZus{}size}\PY{l+s+s1}{\PYZsq{}}\PY{p}{,}
                   \PY{l+s+s1}{\PYZsq{}}\PY{l+s+s1}{output\PYZus{}size}\PY{l+s+s1}{\PYZsq{}}\PY{p}{,}
                   \PY{l+s+s1}{\PYZsq{}}\PY{l+s+s1}{num\PYZus{}params}\PY{l+s+s1}{\PYZsq{}}\PY{p}{]}\PY{p}{)}
\end{Verbatim}
\end{tcolorbox}

            \begin{tcolorbox}[breakable, size=fbox, boxrule=.5pt, pad at break*=1mm, opacityfill=0]
\prompt{Out}{outcolor}{40}{\boxspacing}
\begin{Verbatim}[commandchars=\\\{\}]
================================================================================
===================================
Layer (type:depth-idx)                   Input Shape               Output Shape
Param \#
================================================================================
===================================
MNISTModel                               [256, 1, 28, 28]          [256, 10]
--
├─Sequential: 1-1                        [256, 1, 28, 28]          [256, 10]
--
│    └─Sequential: 2-1                   [256, 1, 28, 28]          [256, 256]
--
│    │    └─Flatten: 3-1                 [256, 1, 28, 28]          [256, 784]
--
│    │    └─Linear: 3-2                  [256, 784]                [256, 256]
200,960
│    │    └─ReLU: 3-3                    [256, 256]                [256, 256]
--
│    │    └─Dropout: 3-4                 [256, 256]                [256, 256]
--
│    └─Sequential: 2-2                   [256, 256]                [256, 128]
--
│    │    └─Linear: 3-5                  [256, 256]                [256, 128]
32,896
│    │    └─ReLU: 3-6                    [256, 128]                [256, 128]
--
│    │    └─Dropout: 3-7                 [256, 128]                [256, 128]
--
│    └─Linear: 2-3                       [256, 128]                [256, 10]
1,290
================================================================================
===================================
Total params: 235,146
Trainable params: 235,146
Non-trainable params: 0
Total mult-adds (Units.MEGABYTES): 60.20
================================================================================
===================================
Input size (MB): 0.80
Forward/backward pass size (MB): 0.81
Params size (MB): 0.94
Estimated Total Size (MB): 2.55
================================================================================
===================================
\end{Verbatim}
\end{tcolorbox}
        
    Having set up both the model and the data module, fitting this model is
now almost identical to the \texttt{Hitters} example. In contrast to our
regression model, here we will use the
\texttt{SimpleModule.classification()} method which uses the
cross-entropy loss function instead of mean squared error. It must be
supplied with the number of classes in the problem.

    \begin{tcolorbox}[breakable, size=fbox, boxrule=1pt, pad at break*=1mm,colback=cellbackground, colframe=cellborder]
\prompt{In}{incolor}{41}{\boxspacing}
\begin{Verbatim}[commandchars=\\\{\}]
\PY{n}{mnist\PYZus{}module} \PY{o}{=} \PY{n}{SimpleModule}\PY{o}{.}\PY{n}{classification}\PY{p}{(}\PY{n}{mnist\PYZus{}model}\PY{p}{,}
                                           \PY{n}{num\PYZus{}classes}\PY{o}{=}\PY{l+m+mi}{10}\PY{p}{)}
\PY{n}{mnist\PYZus{}logger} \PY{o}{=} \PY{n}{CSVLogger}\PY{p}{(}\PY{l+s+s1}{\PYZsq{}}\PY{l+s+s1}{logs}\PY{l+s+s1}{\PYZsq{}}\PY{p}{,} \PY{n}{name}\PY{o}{=}\PY{l+s+s1}{\PYZsq{}}\PY{l+s+s1}{MNIST}\PY{l+s+s1}{\PYZsq{}}\PY{p}{)}
\end{Verbatim}
\end{tcolorbox}

    Now we are ready to go. The final step is to supply training data, and
fit the model. We disable the progress bar below to avoid lengthy output
in the browser when running.

    \begin{tcolorbox}[breakable, size=fbox, boxrule=1pt, pad at break*=1mm,colback=cellbackground, colframe=cellborder]
\prompt{In}{incolor}{42}{\boxspacing}
\begin{Verbatim}[commandchars=\\\{\}]
\PY{n}{mnist\PYZus{}trainer} \PY{o}{=} \PY{n}{Trainer}\PY{p}{(}\PY{n}{deterministic}\PY{o}{=}\PY{k+kc}{True}\PY{p}{,}
                        \PY{n}{max\PYZus{}epochs}\PY{o}{=}\PY{l+m+mi}{30}\PY{p}{,}
                        \PY{n}{logger}\PY{o}{=}\PY{n}{mnist\PYZus{}logger}\PY{p}{,}
                        \PY{n}{enable\PYZus{}progress\PYZus{}bar}\PY{o}{=}\PY{k+kc}{False}\PY{p}{,}
                        \PY{n}{callbacks}\PY{o}{=}\PY{p}{[}\PY{n}{ErrorTracker}\PY{p}{(}\PY{p}{)}\PY{p}{]}\PY{p}{)}
\PY{n}{mnist\PYZus{}trainer}\PY{o}{.}\PY{n}{fit}\PY{p}{(}\PY{n}{mnist\PYZus{}module}\PY{p}{,}
                  \PY{n}{datamodule}\PY{o}{=}\PY{n}{mnist\PYZus{}dm}\PY{p}{)}
\end{Verbatim}
\end{tcolorbox}

    \begin{Verbatim}[commandchars=\\\{\}]
GPU available: True (mps), used: True
TPU available: False, using: 0 TPU cores
IPU available: False, using: 0 IPUs
HPU available: False, using: 0 HPUs

  | Name  | Type             | Params
-------------------------------------------
0 | model | MNISTModel       | 235 K
1 | loss  | CrossEntropyLoss | 0
-------------------------------------------
235 K     Trainable params
0         Non-trainable params
235 K     Total params
0.941     Total estimated model params size (MB)
`Trainer.fit` stopped: `max\_epochs=30` reached.
    \end{Verbatim}

    We have suppressed the output here, which is a progress report on the
fitting of the model, grouped by epoch. This is very useful, since on
large datasets fitting can take time. Fitting this model took 245
seconds on a MacBook Pro with an Apple M1 Pro chip with 10 cores and 16
GB of RAM. Here we specified a validation split of 20\%, so training is
actually performed on 80\% of the 60,000 observations in the training
set. This is an alternative to actually supplying validation data, like
we did for the \texttt{Hitters} data. SGD uses batches of 256
observations in computing the gradient, and doing the arithmetic, we see
that an epoch corresponds to 188 gradient steps.

    \texttt{SimpleModule.classification()} includes an accuracy metric by
default. Other classification metrics can be added from
\texttt{torchmetrics}. We will use our \texttt{summary\_plot()} function
to display accuracy across epochs.

    \begin{tcolorbox}[breakable, size=fbox, boxrule=1pt, pad at break*=1mm,colback=cellbackground, colframe=cellborder]
\prompt{In}{incolor}{43}{\boxspacing}
\begin{Verbatim}[commandchars=\\\{\}]
\PY{n}{mnist\PYZus{}results} \PY{o}{=} \PY{n}{pd}\PY{o}{.}\PY{n}{read\PYZus{}csv}\PY{p}{(}\PY{n}{mnist\PYZus{}logger}\PY{o}{.}\PY{n}{experiment}\PY{o}{.}\PY{n}{metrics\PYZus{}file\PYZus{}path}\PY{p}{)}
\PY{n}{fig}\PY{p}{,} \PY{n}{ax} \PY{o}{=} \PY{n}{subplots}\PY{p}{(}\PY{l+m+mi}{1}\PY{p}{,} \PY{l+m+mi}{1}\PY{p}{,} \PY{n}{figsize}\PY{o}{=}\PY{p}{(}\PY{l+m+mi}{6}\PY{p}{,} \PY{l+m+mi}{6}\PY{p}{)}\PY{p}{)}
\PY{n}{summary\PYZus{}plot}\PY{p}{(}\PY{n}{mnist\PYZus{}results}\PY{p}{,}
             \PY{n}{ax}\PY{p}{,}
             \PY{n}{col}\PY{o}{=}\PY{l+s+s1}{\PYZsq{}}\PY{l+s+s1}{accuracy}\PY{l+s+s1}{\PYZsq{}}\PY{p}{,}
             \PY{n}{ylabel}\PY{o}{=}\PY{l+s+s1}{\PYZsq{}}\PY{l+s+s1}{Accuracy}\PY{l+s+s1}{\PYZsq{}}\PY{p}{)}
\PY{n}{ax}\PY{o}{.}\PY{n}{set\PYZus{}ylim}\PY{p}{(}\PY{p}{[}\PY{l+m+mf}{0.5}\PY{p}{,} \PY{l+m+mi}{1}\PY{p}{]}\PY{p}{)}
\PY{n}{ax}\PY{o}{.}\PY{n}{set\PYZus{}ylabel}\PY{p}{(}\PY{l+s+s1}{\PYZsq{}}\PY{l+s+s1}{Accuracy}\PY{l+s+s1}{\PYZsq{}}\PY{p}{)}
\PY{n}{ax}\PY{o}{.}\PY{n}{set\PYZus{}xticks}\PY{p}{(}\PY{n}{np}\PY{o}{.}\PY{n}{linspace}\PY{p}{(}\PY{l+m+mi}{0}\PY{p}{,} \PY{l+m+mi}{30}\PY{p}{,} \PY{l+m+mi}{7}\PY{p}{)}\PY{o}{.}\PY{n}{astype}\PY{p}{(}\PY{n+nb}{int}\PY{p}{)}\PY{p}{)}\PY{p}{;}
\end{Verbatim}
\end{tcolorbox}

    \begin{center}
    \adjustimage{max size={0.9\linewidth}{0.9\paperheight}}{artifacts/latex/Ch10-deeplearning-lab_files/artifacts/latex/Ch10-deeplearning-lab_90_0.png}
    \end{center}
    { \hspace*{\fill} \\}
    
    Once again we evaluate the accuracy using the \texttt{test()} method of
our trainer. This model achieves 97\% accuracy on the test data.

    \begin{tcolorbox}[breakable, size=fbox, boxrule=1pt, pad at break*=1mm,colback=cellbackground, colframe=cellborder]
\prompt{In}{incolor}{44}{\boxspacing}
\begin{Verbatim}[commandchars=\\\{\}]
\PY{n}{mnist\PYZus{}trainer}\PY{o}{.}\PY{n}{test}\PY{p}{(}\PY{n}{mnist\PYZus{}module}\PY{p}{,}
                   \PY{n}{datamodule}\PY{o}{=}\PY{n}{mnist\PYZus{}dm}\PY{p}{)}
\end{Verbatim}
\end{tcolorbox}

    \begin{Verbatim}[commandchars=\\\{\}]
────────────────────────────────────────────────────────────────────────────────
────────────────────────────────────────
       Test metric             DataLoader 0
────────────────────────────────────────────────────────────────────────────────
────────────────────────────────────────
      test\_accuracy         0.9620000123977661
        test\_loss           0.15120187401771545
────────────────────────────────────────────────────────────────────────────────
────────────────────────────────────────
    \end{Verbatim}

            \begin{tcolorbox}[breakable, size=fbox, boxrule=.5pt, pad at break*=1mm, opacityfill=0]
\prompt{Out}{outcolor}{44}{\boxspacing}
\begin{Verbatim}[commandchars=\\\{\}]
[\{'test\_loss': 0.15120187401771545, 'test\_accuracy': 0.9620000123977661\}]
\end{Verbatim}
\end{tcolorbox}
        
    Table 10.1 also reports the error rates resulting from LDA (Chapter 4)
and multiclass logistic regression. For LDA we refer the reader to
Section 4.7.3. Although we could use the \texttt{sklearn} function
\texttt{LogisticRegression()} to fit\\
multiclass logistic regression, we are set up here to fit such a model
with \texttt{torch}. We just have an input layer and an output layer,
and omit the hidden layers!

    \begin{tcolorbox}[breakable, size=fbox, boxrule=1pt, pad at break*=1mm,colback=cellbackground, colframe=cellborder]
\prompt{In}{incolor}{45}{\boxspacing}
\begin{Verbatim}[commandchars=\\\{\}]
\PY{k}{class}\PY{+w}{ }\PY{n+nc}{MNIST\PYZus{}MLR}\PY{p}{(}\PY{n}{nn}\PY{o}{.}\PY{n}{Module}\PY{p}{)}\PY{p}{:}
    \PY{k}{def}\PY{+w}{ }\PY{n+nf+fm}{\PYZus{}\PYZus{}init\PYZus{}\PYZus{}}\PY{p}{(}\PY{n+nb+bp}{self}\PY{p}{)}\PY{p}{:}
        \PY{n+nb}{super}\PY{p}{(}\PY{n}{MNIST\PYZus{}MLR}\PY{p}{,} \PY{n+nb+bp}{self}\PY{p}{)}\PY{o}{.}\PY{n+nf+fm}{\PYZus{}\PYZus{}init\PYZus{}\PYZus{}}\PY{p}{(}\PY{p}{)}
        \PY{n+nb+bp}{self}\PY{o}{.}\PY{n}{linear} \PY{o}{=} \PY{n}{nn}\PY{o}{.}\PY{n}{Sequential}\PY{p}{(}\PY{n}{nn}\PY{o}{.}\PY{n}{Flatten}\PY{p}{(}\PY{p}{)}\PY{p}{,}
                                    \PY{n}{nn}\PY{o}{.}\PY{n}{Linear}\PY{p}{(}\PY{l+m+mi}{784}\PY{p}{,} \PY{l+m+mi}{10}\PY{p}{)}\PY{p}{)}
    \PY{k}{def}\PY{+w}{ }\PY{n+nf}{forward}\PY{p}{(}\PY{n+nb+bp}{self}\PY{p}{,} \PY{n}{x}\PY{p}{)}\PY{p}{:}
        \PY{k}{return} \PY{n+nb+bp}{self}\PY{o}{.}\PY{n}{linear}\PY{p}{(}\PY{n}{x}\PY{p}{)}

\PY{n}{mlr\PYZus{}model} \PY{o}{=} \PY{n}{MNIST\PYZus{}MLR}\PY{p}{(}\PY{p}{)}
\PY{n}{mlr\PYZus{}module} \PY{o}{=} \PY{n}{SimpleModule}\PY{o}{.}\PY{n}{classification}\PY{p}{(}\PY{n}{mlr\PYZus{}model}\PY{p}{,}
                                         \PY{n}{num\PYZus{}classes}\PY{o}{=}\PY{l+m+mi}{10}\PY{p}{)}
\PY{n}{mlr\PYZus{}logger} \PY{o}{=} \PY{n}{CSVLogger}\PY{p}{(}\PY{l+s+s1}{\PYZsq{}}\PY{l+s+s1}{logs}\PY{l+s+s1}{\PYZsq{}}\PY{p}{,} \PY{n}{name}\PY{o}{=}\PY{l+s+s1}{\PYZsq{}}\PY{l+s+s1}{MNIST\PYZus{}MLR}\PY{l+s+s1}{\PYZsq{}}\PY{p}{)}
\end{Verbatim}
\end{tcolorbox}

    \begin{tcolorbox}[breakable, size=fbox, boxrule=1pt, pad at break*=1mm,colback=cellbackground, colframe=cellborder]
\prompt{In}{incolor}{46}{\boxspacing}
\begin{Verbatim}[commandchars=\\\{\}]
\PY{n}{mlr\PYZus{}trainer} \PY{o}{=} \PY{n}{Trainer}\PY{p}{(}\PY{n}{deterministic}\PY{o}{=}\PY{k+kc}{True}\PY{p}{,}
                      \PY{n}{max\PYZus{}epochs}\PY{o}{=}\PY{l+m+mi}{30}\PY{p}{,}
                      \PY{n}{enable\PYZus{}progress\PYZus{}bar}\PY{o}{=}\PY{k+kc}{False}\PY{p}{,}
                      \PY{n}{callbacks}\PY{o}{=}\PY{p}{[}\PY{n}{ErrorTracker}\PY{p}{(}\PY{p}{)}\PY{p}{]}\PY{p}{)}
\PY{n}{mlr\PYZus{}trainer}\PY{o}{.}\PY{n}{fit}\PY{p}{(}\PY{n}{mlr\PYZus{}module}\PY{p}{,} \PY{n}{datamodule}\PY{o}{=}\PY{n}{mnist\PYZus{}dm}\PY{p}{)}
\end{Verbatim}
\end{tcolorbox}

    \begin{Verbatim}[commandchars=\\\{\}]
GPU available: True (mps), used: True
TPU available: False, using: 0 TPU cores
IPU available: False, using: 0 IPUs
HPU available: False, using: 0 HPUs
/Users/jonathantaylor/anaconda3/envs/ISLP\_312/lib/python3.12/site-packages/pytor
ch\_lightning/trainer/connectors/logger\_connector/logger\_connector.py:75:
Starting from v1.9.0, `tensorboardX` has been removed as a dependency of the
`pytorch\_lightning` package, due to potential conflicts with other packages in
the ML ecosystem. For this reason, `logger=True` will use `CSVLogger` as the
default logger, unless the `tensorboard` or `tensorboardX` packages are found.
Please `pip install lightning[extra]` or one of them to enable TensorBoard
support by default

  | Name  | Type             | Params
-------------------------------------------
0 | model | MNIST\_MLR        | 7.9 K
1 | loss  | CrossEntropyLoss | 0
-------------------------------------------
7.9 K     Trainable params
0         Non-trainable params
7.9 K     Total params
0.031     Total estimated model params size (MB)
`Trainer.fit` stopped: `max\_epochs=30` reached.
    \end{Verbatim}

    We fit the model just as before and compute the test results.

    \begin{tcolorbox}[breakable, size=fbox, boxrule=1pt, pad at break*=1mm,colback=cellbackground, colframe=cellborder]
\prompt{In}{incolor}{47}{\boxspacing}
\begin{Verbatim}[commandchars=\\\{\}]
\PY{n}{mlr\PYZus{}trainer}\PY{o}{.}\PY{n}{test}\PY{p}{(}\PY{n}{mlr\PYZus{}module}\PY{p}{,}
                 \PY{n}{datamodule}\PY{o}{=}\PY{n}{mnist\PYZus{}dm}\PY{p}{)}
\end{Verbatim}
\end{tcolorbox}

    \begin{Verbatim}[commandchars=\\\{\}]
────────────────────────────────────────────────────────────────────────────────
────────────────────────────────────────
       Test metric             DataLoader 0
────────────────────────────────────────────────────────────────────────────────
────────────────────────────────────────
      test\_accuracy          0.916100025177002
        test\_loss           0.3469300866127014
────────────────────────────────────────────────────────────────────────────────
────────────────────────────────────────
    \end{Verbatim}

            \begin{tcolorbox}[breakable, size=fbox, boxrule=.5pt, pad at break*=1mm, opacityfill=0]
\prompt{Out}{outcolor}{47}{\boxspacing}
\begin{Verbatim}[commandchars=\\\{\}]
[\{'test\_loss': 0.3469300866127014, 'test\_accuracy': 0.916100025177002\}]
\end{Verbatim}
\end{tcolorbox}
        
    The accuracy is above 90\% even for this pretty simple model.

As in the \texttt{Hitters} example, we delete some of the objects we
created above.

    \begin{tcolorbox}[breakable, size=fbox, boxrule=1pt, pad at break*=1mm,colback=cellbackground, colframe=cellborder]
\prompt{In}{incolor}{48}{\boxspacing}
\begin{Verbatim}[commandchars=\\\{\}]
\PY{k}{del}\PY{p}{(}\PY{n}{mnist\PYZus{}test}\PY{p}{,}
    \PY{n}{mnist\PYZus{}train}\PY{p}{,}
    \PY{n}{mnist\PYZus{}model}\PY{p}{,}
    \PY{n}{mnist\PYZus{}dm}\PY{p}{,}
    \PY{n}{mnist\PYZus{}trainer}\PY{p}{,}
    \PY{n}{mnist\PYZus{}module}\PY{p}{,}
    \PY{n}{mnist\PYZus{}results}\PY{p}{,}
    \PY{n}{mlr\PYZus{}model}\PY{p}{,}
    \PY{n}{mlr\PYZus{}module}\PY{p}{,}
    \PY{n}{mlr\PYZus{}trainer}\PY{p}{)}
\end{Verbatim}
\end{tcolorbox}

    \hypertarget{convolutional-neural-networks}{%
\subsection{Convolutional Neural
Networks}\label{convolutional-neural-networks}}

In this section we fit a CNN to the \texttt{CIFAR100} data, which is
available in the \texttt{torchvision} package. It is arranged in a
similar fashion as the \texttt{MNIST} data.

    \begin{tcolorbox}[breakable, size=fbox, boxrule=1pt, pad at break*=1mm,colback=cellbackground, colframe=cellborder]
\prompt{In}{incolor}{49}{\boxspacing}
\begin{Verbatim}[commandchars=\\\{\}]
\PY{p}{(}\PY{n}{cifar\PYZus{}train}\PY{p}{,} 
 \PY{n}{cifar\PYZus{}test}\PY{p}{)} \PY{o}{=} \PY{p}{[}\PY{n}{CIFAR100}\PY{p}{(}\PY{n}{root}\PY{o}{=}\PY{l+s+s2}{\PYZdq{}}\PY{l+s+s2}{data}\PY{l+s+s2}{\PYZdq{}}\PY{p}{,}
                         \PY{n}{train}\PY{o}{=}\PY{n}{train}\PY{p}{,}
                         \PY{n}{download}\PY{o}{=}\PY{k+kc}{True}\PY{p}{)}
             \PY{k}{for} \PY{n}{train} \PY{o+ow}{in} \PY{p}{[}\PY{k+kc}{True}\PY{p}{,} \PY{k+kc}{False}\PY{p}{]}\PY{p}{]}
\end{Verbatim}
\end{tcolorbox}

    \begin{Verbatim}[commandchars=\\\{\}]
Files already downloaded and verified
Files already downloaded and verified
    \end{Verbatim}

    \begin{tcolorbox}[breakable, size=fbox, boxrule=1pt, pad at break*=1mm,colback=cellbackground, colframe=cellborder]
\prompt{In}{incolor}{50}{\boxspacing}
\begin{Verbatim}[commandchars=\\\{\}]
\PY{n}{transform} \PY{o}{=} \PY{n}{ToTensor}\PY{p}{(}\PY{p}{)}
\PY{n}{cifar\PYZus{}train\PYZus{}X} \PY{o}{=} \PY{n}{torch}\PY{o}{.}\PY{n}{stack}\PY{p}{(}\PY{p}{[}\PY{n}{transform}\PY{p}{(}\PY{n}{x}\PY{p}{)} \PY{k}{for} \PY{n}{x} \PY{o+ow}{in}
                            \PY{n}{cifar\PYZus{}train}\PY{o}{.}\PY{n}{data}\PY{p}{]}\PY{p}{)}
\PY{n}{cifar\PYZus{}test\PYZus{}X} \PY{o}{=} \PY{n}{torch}\PY{o}{.}\PY{n}{stack}\PY{p}{(}\PY{p}{[}\PY{n}{transform}\PY{p}{(}\PY{n}{x}\PY{p}{)} \PY{k}{for} \PY{n}{x} \PY{o+ow}{in}
                            \PY{n}{cifar\PYZus{}test}\PY{o}{.}\PY{n}{data}\PY{p}{]}\PY{p}{)}
\PY{n}{cifar\PYZus{}train} \PY{o}{=} \PY{n}{TensorDataset}\PY{p}{(}\PY{n}{cifar\PYZus{}train\PYZus{}X}\PY{p}{,}
                            \PY{n}{torch}\PY{o}{.}\PY{n}{tensor}\PY{p}{(}\PY{n}{cifar\PYZus{}train}\PY{o}{.}\PY{n}{targets}\PY{p}{)}\PY{p}{)}
\PY{n}{cifar\PYZus{}test} \PY{o}{=} \PY{n}{TensorDataset}\PY{p}{(}\PY{n}{cifar\PYZus{}test\PYZus{}X}\PY{p}{,}
                            \PY{n}{torch}\PY{o}{.}\PY{n}{tensor}\PY{p}{(}\PY{n}{cifar\PYZus{}test}\PY{o}{.}\PY{n}{targets}\PY{p}{)}\PY{p}{)}
\end{Verbatim}
\end{tcolorbox}

    The \texttt{CIFAR100} dataset consists of 50,000 training images, each
represented by a three-dimensional tensor: each three-color image is
represented as a set of three channels, each of which consists of
\(32\times 32\) eight-bit pixels. We standardize as we did for the
digits, but keep the array structure. This is accomplished with the
\texttt{ToTensor()} transform.

Creating the data module is similar to the \texttt{MNIST} example.

    \begin{tcolorbox}[breakable, size=fbox, boxrule=1pt, pad at break*=1mm,colback=cellbackground, colframe=cellborder]
\prompt{In}{incolor}{51}{\boxspacing}
\begin{Verbatim}[commandchars=\\\{\}]
\PY{n}{cifar\PYZus{}dm} \PY{o}{=} \PY{n}{SimpleDataModule}\PY{p}{(}\PY{n}{cifar\PYZus{}train}\PY{p}{,}
                            \PY{n}{cifar\PYZus{}test}\PY{p}{,}
                            \PY{n}{validation}\PY{o}{=}\PY{l+m+mf}{0.2}\PY{p}{,}
                            \PY{n}{num\PYZus{}workers}\PY{o}{=}\PY{n}{max\PYZus{}num\PYZus{}workers}\PY{p}{,}
                            \PY{n}{batch\PYZus{}size}\PY{o}{=}\PY{l+m+mi}{128}\PY{p}{)}
\end{Verbatim}
\end{tcolorbox}

    We again look at the shape of typical batches in our data loaders.

    \begin{tcolorbox}[breakable, size=fbox, boxrule=1pt, pad at break*=1mm,colback=cellbackground, colframe=cellborder]
\prompt{In}{incolor}{52}{\boxspacing}
\begin{Verbatim}[commandchars=\\\{\}]
\PY{k}{for} \PY{n}{idx}\PY{p}{,} \PY{p}{(}\PY{n}{X\PYZus{}} \PY{p}{,}\PY{n}{Y\PYZus{}}\PY{p}{)} \PY{o+ow}{in} \PY{n+nb}{enumerate}\PY{p}{(}\PY{n}{cifar\PYZus{}dm}\PY{o}{.}\PY{n}{train\PYZus{}dataloader}\PY{p}{(}\PY{p}{)}\PY{p}{)}\PY{p}{:}
    \PY{n+nb}{print}\PY{p}{(}\PY{l+s+s1}{\PYZsq{}}\PY{l+s+s1}{X: }\PY{l+s+s1}{\PYZsq{}}\PY{p}{,} \PY{n}{X\PYZus{}}\PY{o}{.}\PY{n}{shape}\PY{p}{)}
    \PY{n+nb}{print}\PY{p}{(}\PY{l+s+s1}{\PYZsq{}}\PY{l+s+s1}{Y: }\PY{l+s+s1}{\PYZsq{}}\PY{p}{,} \PY{n}{Y\PYZus{}}\PY{o}{.}\PY{n}{shape}\PY{p}{)}
    \PY{k}{if} \PY{n}{idx} \PY{o}{\PYZgt{}}\PY{o}{=} \PY{l+m+mi}{1}\PY{p}{:}
        \PY{k}{break}
\end{Verbatim}
\end{tcolorbox}

    \begin{Verbatim}[commandchars=\\\{\}]
X:  torch.Size([128, 3, 32, 32])
Y:  torch.Size([128])
X:  torch.Size([128, 3, 32, 32])
Y:  torch.Size([128])
    \end{Verbatim}

    Before we start, we look at some of the training images; similar code
produced Figure 10.5 on page 406. The example below also illustrates
that \texttt{TensorDataset} objects can be indexed with integers --- we
are choosing random images from the training data by indexing
\texttt{cifar\_train}. In order to display correctly, we must reorder
the dimensions by a call to \texttt{np.transpose()}.

    \begin{tcolorbox}[breakable, size=fbox, boxrule=1pt, pad at break*=1mm,colback=cellbackground, colframe=cellborder]
\prompt{In}{incolor}{53}{\boxspacing}
\begin{Verbatim}[commandchars=\\\{\}]
\PY{n}{fig}\PY{p}{,} \PY{n}{axes} \PY{o}{=} \PY{n}{subplots}\PY{p}{(}\PY{l+m+mi}{5}\PY{p}{,} \PY{l+m+mi}{5}\PY{p}{,} \PY{n}{figsize}\PY{o}{=}\PY{p}{(}\PY{l+m+mi}{10}\PY{p}{,}\PY{l+m+mi}{10}\PY{p}{)}\PY{p}{)}
\PY{n}{rng} \PY{o}{=} \PY{n}{np}\PY{o}{.}\PY{n}{random}\PY{o}{.}\PY{n}{default\PYZus{}rng}\PY{p}{(}\PY{l+m+mi}{4}\PY{p}{)}
\PY{n}{indices} \PY{o}{=} \PY{n}{rng}\PY{o}{.}\PY{n}{choice}\PY{p}{(}\PY{n}{np}\PY{o}{.}\PY{n}{arange}\PY{p}{(}\PY{n+nb}{len}\PY{p}{(}\PY{n}{cifar\PYZus{}train}\PY{p}{)}\PY{p}{)}\PY{p}{,} \PY{l+m+mi}{25}\PY{p}{,}
                     \PY{n}{replace}\PY{o}{=}\PY{k+kc}{False}\PY{p}{)}\PY{o}{.}\PY{n}{reshape}\PY{p}{(}\PY{p}{(}\PY{l+m+mi}{5}\PY{p}{,}\PY{l+m+mi}{5}\PY{p}{)}\PY{p}{)}
\PY{k}{for} \PY{n}{i} \PY{o+ow}{in} \PY{n+nb}{range}\PY{p}{(}\PY{l+m+mi}{5}\PY{p}{)}\PY{p}{:}
    \PY{k}{for} \PY{n}{j} \PY{o+ow}{in} \PY{n+nb}{range}\PY{p}{(}\PY{l+m+mi}{5}\PY{p}{)}\PY{p}{:}
        \PY{n}{idx} \PY{o}{=} \PY{n}{indices}\PY{p}{[}\PY{n}{i}\PY{p}{,}\PY{n}{j}\PY{p}{]}
        \PY{n}{axes}\PY{p}{[}\PY{n}{i}\PY{p}{,}\PY{n}{j}\PY{p}{]}\PY{o}{.}\PY{n}{imshow}\PY{p}{(}\PY{n}{np}\PY{o}{.}\PY{n}{transpose}\PY{p}{(}\PY{n}{cifar\PYZus{}train}\PY{p}{[}\PY{n}{idx}\PY{p}{]}\PY{p}{[}\PY{l+m+mi}{0}\PY{p}{]}\PY{p}{,}
                                      \PY{p}{[}\PY{l+m+mi}{1}\PY{p}{,}\PY{l+m+mi}{2}\PY{p}{,}\PY{l+m+mi}{0}\PY{p}{]}\PY{p}{)}\PY{p}{,}
                                      \PY{n}{interpolation}\PY{o}{=}\PY{k+kc}{None}\PY{p}{)}
        \PY{n}{axes}\PY{p}{[}\PY{n}{i}\PY{p}{,}\PY{n}{j}\PY{p}{]}\PY{o}{.}\PY{n}{set\PYZus{}xticks}\PY{p}{(}\PY{p}{[}\PY{p}{]}\PY{p}{)}
        \PY{n}{axes}\PY{p}{[}\PY{n}{i}\PY{p}{,}\PY{n}{j}\PY{p}{]}\PY{o}{.}\PY{n}{set\PYZus{}yticks}\PY{p}{(}\PY{p}{[}\PY{p}{]}\PY{p}{)}
\end{Verbatim}
\end{tcolorbox}

    \begin{center}
    \adjustimage{max size={0.9\linewidth}{0.9\paperheight}}{artifacts/latex/Ch10-deeplearning-lab_files/artifacts/latex/Ch10-deeplearning-lab_108_0.png}
    \end{center}
    { \hspace*{\fill} \\}
    
    Here the \texttt{imshow()} method recognizes from the shape of its
argument that it is a 3-dimensional array, with the last dimension
indexing the three RGB color channels.

We specify a moderately-sized CNN for demonstration purposes, similar in
structure to Figure 10.8. We use several layers, each consisting of
convolution, ReLU, and max-pooling steps. We first define a module that
defines one of these layers. As in our previous examples, we overwrite
the \texttt{\_\_init\_\_()} and \texttt{forward()} methods of
\texttt{nn.Module}. This user-defined module can now be used in ways
just like \texttt{nn.Linear()} or \texttt{nn.Dropout()}.

    \begin{tcolorbox}[breakable, size=fbox, boxrule=1pt, pad at break*=1mm,colback=cellbackground, colframe=cellborder]
\prompt{In}{incolor}{54}{\boxspacing}
\begin{Verbatim}[commandchars=\\\{\}]
\PY{k}{class}\PY{+w}{ }\PY{n+nc}{BuildingBlock}\PY{p}{(}\PY{n}{nn}\PY{o}{.}\PY{n}{Module}\PY{p}{)}\PY{p}{:}

    \PY{k}{def}\PY{+w}{ }\PY{n+nf+fm}{\PYZus{}\PYZus{}init\PYZus{}\PYZus{}}\PY{p}{(}\PY{n+nb+bp}{self}\PY{p}{,}
                 \PY{n}{in\PYZus{}channels}\PY{p}{,}
                 \PY{n}{out\PYZus{}channels}\PY{p}{)}\PY{p}{:}

        \PY{n+nb}{super}\PY{p}{(}\PY{n}{BuildingBlock}\PY{p}{,} \PY{n+nb+bp}{self}\PY{p}{)}\PY{o}{.}\PY{n+nf+fm}{\PYZus{}\PYZus{}init\PYZus{}\PYZus{}}\PY{p}{(}\PY{p}{)}
        \PY{n+nb+bp}{self}\PY{o}{.}\PY{n}{conv} \PY{o}{=} \PY{n}{nn}\PY{o}{.}\PY{n}{Conv2d}\PY{p}{(}\PY{n}{in\PYZus{}channels}\PY{o}{=}\PY{n}{in\PYZus{}channels}\PY{p}{,}
                              \PY{n}{out\PYZus{}channels}\PY{o}{=}\PY{n}{out\PYZus{}channels}\PY{p}{,}
                              \PY{n}{kernel\PYZus{}size}\PY{o}{=}\PY{p}{(}\PY{l+m+mi}{3}\PY{p}{,}\PY{l+m+mi}{3}\PY{p}{)}\PY{p}{,}
                              \PY{n}{padding}\PY{o}{=}\PY{l+s+s1}{\PYZsq{}}\PY{l+s+s1}{same}\PY{l+s+s1}{\PYZsq{}}\PY{p}{)}
        \PY{n+nb+bp}{self}\PY{o}{.}\PY{n}{activation} \PY{o}{=} \PY{n}{nn}\PY{o}{.}\PY{n}{ReLU}\PY{p}{(}\PY{p}{)}
        \PY{n+nb+bp}{self}\PY{o}{.}\PY{n}{pool} \PY{o}{=} \PY{n}{nn}\PY{o}{.}\PY{n}{MaxPool2d}\PY{p}{(}\PY{n}{kernel\PYZus{}size}\PY{o}{=}\PY{p}{(}\PY{l+m+mi}{2}\PY{p}{,}\PY{l+m+mi}{2}\PY{p}{)}\PY{p}{)}

    \PY{k}{def}\PY{+w}{ }\PY{n+nf}{forward}\PY{p}{(}\PY{n+nb+bp}{self}\PY{p}{,} \PY{n}{x}\PY{p}{)}\PY{p}{:}
        \PY{k}{return} \PY{n+nb+bp}{self}\PY{o}{.}\PY{n}{pool}\PY{p}{(}\PY{n+nb+bp}{self}\PY{o}{.}\PY{n}{activation}\PY{p}{(}\PY{n+nb+bp}{self}\PY{o}{.}\PY{n}{conv}\PY{p}{(}\PY{n}{x}\PY{p}{)}\PY{p}{)}\PY{p}{)}
\end{Verbatim}
\end{tcolorbox}

    Notice that we used the \texttt{padding\ =\ "same"} argument to
\texttt{nn.Conv2d()}, which ensures that the output channels have the
same dimension as the input channels. There are 32 channels in the first
hidden layer, in contrast to the three channels in the input layer. We
use a \(3\times 3\) convolution filter for each channel in all the
layers. Each convolution is followed by a max-pooling layer over
\(2\times2\) blocks.

In forming our deep learning model for the \texttt{CIFAR100} data, we
use several of our \texttt{BuildingBlock()} modules sequentially. This
simple example illustrates some of the power of \texttt{torch}. Users
can define modules of their own, which can be combined in other modules.
Ultimately, everything is fit by a generic trainer.

    \begin{tcolorbox}[breakable, size=fbox, boxrule=1pt, pad at break*=1mm,colback=cellbackground, colframe=cellborder]
\prompt{In}{incolor}{55}{\boxspacing}
\begin{Verbatim}[commandchars=\\\{\}]
\PY{k}{class}\PY{+w}{ }\PY{n+nc}{CIFARModel}\PY{p}{(}\PY{n}{nn}\PY{o}{.}\PY{n}{Module}\PY{p}{)}\PY{p}{:}

    \PY{k}{def}\PY{+w}{ }\PY{n+nf+fm}{\PYZus{}\PYZus{}init\PYZus{}\PYZus{}}\PY{p}{(}\PY{n+nb+bp}{self}\PY{p}{)}\PY{p}{:}
        \PY{n+nb}{super}\PY{p}{(}\PY{n}{CIFARModel}\PY{p}{,} \PY{n+nb+bp}{self}\PY{p}{)}\PY{o}{.}\PY{n+nf+fm}{\PYZus{}\PYZus{}init\PYZus{}\PYZus{}}\PY{p}{(}\PY{p}{)}
        \PY{n}{sizes} \PY{o}{=} \PY{p}{[}\PY{p}{(}\PY{l+m+mi}{3}\PY{p}{,}\PY{l+m+mi}{32}\PY{p}{)}\PY{p}{,}
                 \PY{p}{(}\PY{l+m+mi}{32}\PY{p}{,}\PY{l+m+mi}{64}\PY{p}{)}\PY{p}{,}
                 \PY{p}{(}\PY{l+m+mi}{64}\PY{p}{,}\PY{l+m+mi}{128}\PY{p}{)}\PY{p}{,}
                 \PY{p}{(}\PY{l+m+mi}{128}\PY{p}{,}\PY{l+m+mi}{256}\PY{p}{)}\PY{p}{]}
        \PY{n+nb+bp}{self}\PY{o}{.}\PY{n}{conv} \PY{o}{=} \PY{n}{nn}\PY{o}{.}\PY{n}{Sequential}\PY{p}{(}\PY{o}{*}\PY{p}{[}\PY{n}{BuildingBlock}\PY{p}{(}\PY{n}{in\PYZus{}}\PY{p}{,} \PY{n}{out\PYZus{}}\PY{p}{)}
                                    \PY{k}{for} \PY{n}{in\PYZus{}}\PY{p}{,} \PY{n}{out\PYZus{}} \PY{o+ow}{in} \PY{n}{sizes}\PY{p}{]}\PY{p}{)}

        \PY{n+nb+bp}{self}\PY{o}{.}\PY{n}{output} \PY{o}{=} \PY{n}{nn}\PY{o}{.}\PY{n}{Sequential}\PY{p}{(}\PY{n}{nn}\PY{o}{.}\PY{n}{Dropout}\PY{p}{(}\PY{l+m+mf}{0.5}\PY{p}{)}\PY{p}{,}
                                    \PY{n}{nn}\PY{o}{.}\PY{n}{Linear}\PY{p}{(}\PY{l+m+mi}{2}\PY{o}{*}\PY{l+m+mi}{2}\PY{o}{*}\PY{l+m+mi}{256}\PY{p}{,} \PY{l+m+mi}{512}\PY{p}{)}\PY{p}{,}
                                    \PY{n}{nn}\PY{o}{.}\PY{n}{ReLU}\PY{p}{(}\PY{p}{)}\PY{p}{,}
                                    \PY{n}{nn}\PY{o}{.}\PY{n}{Linear}\PY{p}{(}\PY{l+m+mi}{512}\PY{p}{,} \PY{l+m+mi}{100}\PY{p}{)}\PY{p}{)}
    \PY{k}{def}\PY{+w}{ }\PY{n+nf}{forward}\PY{p}{(}\PY{n+nb+bp}{self}\PY{p}{,} \PY{n}{x}\PY{p}{)}\PY{p}{:}
        \PY{n}{val} \PY{o}{=} \PY{n+nb+bp}{self}\PY{o}{.}\PY{n}{conv}\PY{p}{(}\PY{n}{x}\PY{p}{)}
        \PY{n}{val} \PY{o}{=} \PY{n}{torch}\PY{o}{.}\PY{n}{flatten}\PY{p}{(}\PY{n}{val}\PY{p}{,} \PY{n}{start\PYZus{}dim}\PY{o}{=}\PY{l+m+mi}{1}\PY{p}{)}
        \PY{k}{return} \PY{n+nb+bp}{self}\PY{o}{.}\PY{n}{output}\PY{p}{(}\PY{n}{val}\PY{p}{)}
\end{Verbatim}
\end{tcolorbox}

    We build the model and look at the summary. (We had created examples of
\texttt{X\_} earlier.)

    \begin{tcolorbox}[breakable, size=fbox, boxrule=1pt, pad at break*=1mm,colback=cellbackground, colframe=cellborder]
\prompt{In}{incolor}{56}{\boxspacing}
\begin{Verbatim}[commandchars=\\\{\}]
\PY{n}{cifar\PYZus{}model} \PY{o}{=} \PY{n}{CIFARModel}\PY{p}{(}\PY{p}{)}
\PY{n}{summary}\PY{p}{(}\PY{n}{cifar\PYZus{}model}\PY{p}{,}
        \PY{n}{input\PYZus{}data}\PY{o}{=}\PY{n}{X\PYZus{}}\PY{p}{,}
        \PY{n}{col\PYZus{}names}\PY{o}{=}\PY{p}{[}\PY{l+s+s1}{\PYZsq{}}\PY{l+s+s1}{input\PYZus{}size}\PY{l+s+s1}{\PYZsq{}}\PY{p}{,}
                   \PY{l+s+s1}{\PYZsq{}}\PY{l+s+s1}{output\PYZus{}size}\PY{l+s+s1}{\PYZsq{}}\PY{p}{,}
                   \PY{l+s+s1}{\PYZsq{}}\PY{l+s+s1}{num\PYZus{}params}\PY{l+s+s1}{\PYZsq{}}\PY{p}{]}\PY{p}{)}
\end{Verbatim}
\end{tcolorbox}

            \begin{tcolorbox}[breakable, size=fbox, boxrule=.5pt, pad at break*=1mm, opacityfill=0]
\prompt{Out}{outcolor}{56}{\boxspacing}
\begin{Verbatim}[commandchars=\\\{\}]
================================================================================
===================================
Layer (type:depth-idx)                   Input Shape               Output Shape
Param \#
================================================================================
===================================
CIFARModel                               [128, 3, 32, 32]          [128, 100]
--
├─Sequential: 1-1                        [128, 3, 32, 32]          [128, 256, 2,
2]          --
│    └─BuildingBlock: 2-1                [128, 3, 32, 32]          [128, 32, 16,
16]         --
│    │    └─Conv2d: 3-1                  [128, 3, 32, 32]          [128, 32, 32,
32]         896
│    │    └─ReLU: 3-2                    [128, 32, 32, 32]         [128, 32, 32,
32]         --
│    │    └─MaxPool2d: 3-3               [128, 32, 32, 32]         [128, 32, 16,
16]         --
│    └─BuildingBlock: 2-2                [128, 32, 16, 16]         [128, 64, 8,
8]           --
│    │    └─Conv2d: 3-4                  [128, 32, 16, 16]         [128, 64, 16,
16]         18,496
│    │    └─ReLU: 3-5                    [128, 64, 16, 16]         [128, 64, 16,
16]         --
│    │    └─MaxPool2d: 3-6               [128, 64, 16, 16]         [128, 64, 8,
8]           --
│    └─BuildingBlock: 2-3                [128, 64, 8, 8]           [128, 128, 4,
4]          --
│    │    └─Conv2d: 3-7                  [128, 64, 8, 8]           [128, 128, 8,
8]          73,856
│    │    └─ReLU: 3-8                    [128, 128, 8, 8]          [128, 128, 8,
8]          --
│    │    └─MaxPool2d: 3-9               [128, 128, 8, 8]          [128, 128, 4,
4]          --
│    └─BuildingBlock: 2-4                [128, 128, 4, 4]          [128, 256, 2,
2]          --
│    │    └─Conv2d: 3-10                 [128, 128, 4, 4]          [128, 256, 4,
4]          295,168
│    │    └─ReLU: 3-11                   [128, 256, 4, 4]          [128, 256, 4,
4]          --
│    │    └─MaxPool2d: 3-12              [128, 256, 4, 4]          [128, 256, 2,
2]          --
├─Sequential: 1-2                        [128, 1024]               [128, 100]
--
│    └─Dropout: 2-5                      [128, 1024]               [128, 1024]
--
│    └─Linear: 2-6                       [128, 1024]               [128, 512]
524,800
│    └─ReLU: 2-7                         [128, 512]                [128, 512]
--
│    └─Linear: 2-8                       [128, 512]                [128, 100]
51,300
================================================================================
===================================
Total params: 964,516
Trainable params: 964,516
Non-trainable params: 0
Total mult-adds (Units.GIGABYTES): 2.01
================================================================================
===================================
Input size (MB): 1.57
Forward/backward pass size (MB): 63.54
Params size (MB): 3.86
Estimated Total Size (MB): 68.97
================================================================================
===================================
\end{Verbatim}
\end{tcolorbox}
        
    The total number of trainable parameters is 964,516. By studying the
size of the parameters, we can see that the channels halve in both
dimensions after each of these max-pooling operations. After the last of
these we have a layer with 256 channels of dimension \(2\times 2\).
These are then flattened to a dense layer of size 1,024; in other words,
each of the \(2\times 2\) matrices is turned into a \(4\)-vector, and
put side-by-side in one layer. This is followed by a dropout
regularization layer, then another dense layer of size 512, and finally,
the output layer.

Up to now, we have been using a default optimizer in
\texttt{SimpleModule()}. For these data, experiments show that a smaller
learning rate performs better than the default 0.01. We use a custom
optimizer here with a learning rate of 0.001. Besides this, the logging
and training follow a similar pattern to our previous examples. The
optimizer takes an argument \texttt{params} that informs the optimizer
which parameters are involved in SGD (stochastic gradient descent).

We saw earlier that entries of a module's parameters are tensors. In
passing the parameters to the optimizer we are doing more than simply
passing arrays; part of the structure of the graph is encoded in the
tensors themselves.

    \begin{tcolorbox}[breakable, size=fbox, boxrule=1pt, pad at break*=1mm,colback=cellbackground, colframe=cellborder]
\prompt{In}{incolor}{57}{\boxspacing}
\begin{Verbatim}[commandchars=\\\{\}]
\PY{n}{cifar\PYZus{}optimizer} \PY{o}{=} \PY{n}{RMSprop}\PY{p}{(}\PY{n}{cifar\PYZus{}model}\PY{o}{.}\PY{n}{parameters}\PY{p}{(}\PY{p}{)}\PY{p}{,} \PY{n}{lr}\PY{o}{=}\PY{l+m+mf}{0.001}\PY{p}{)}
\PY{n}{cifar\PYZus{}module} \PY{o}{=} \PY{n}{SimpleModule}\PY{o}{.}\PY{n}{classification}\PY{p}{(}\PY{n}{cifar\PYZus{}model}\PY{p}{,}
                                    \PY{n}{num\PYZus{}classes}\PY{o}{=}\PY{l+m+mi}{100}\PY{p}{,}
                                    \PY{n}{optimizer}\PY{o}{=}\PY{n}{cifar\PYZus{}optimizer}\PY{p}{)}
\PY{n}{cifar\PYZus{}logger} \PY{o}{=} \PY{n}{CSVLogger}\PY{p}{(}\PY{l+s+s1}{\PYZsq{}}\PY{l+s+s1}{logs}\PY{l+s+s1}{\PYZsq{}}\PY{p}{,} \PY{n}{name}\PY{o}{=}\PY{l+s+s1}{\PYZsq{}}\PY{l+s+s1}{CIFAR100}\PY{l+s+s1}{\PYZsq{}}\PY{p}{)}
\end{Verbatim}
\end{tcolorbox}

    \begin{tcolorbox}[breakable, size=fbox, boxrule=1pt, pad at break*=1mm,colback=cellbackground, colframe=cellborder]
\prompt{In}{incolor}{58}{\boxspacing}
\begin{Verbatim}[commandchars=\\\{\}]
\PY{n}{cifar\PYZus{}trainer} \PY{o}{=} \PY{n}{Trainer}\PY{p}{(}\PY{n}{deterministic}\PY{o}{=}\PY{k+kc}{True}\PY{p}{,}
                        \PY{n}{max\PYZus{}epochs}\PY{o}{=}\PY{l+m+mi}{30}\PY{p}{,}
                        \PY{n}{logger}\PY{o}{=}\PY{n}{cifar\PYZus{}logger}\PY{p}{,}
                        \PY{n}{enable\PYZus{}progress\PYZus{}bar}\PY{o}{=}\PY{k+kc}{False}\PY{p}{,}
                        \PY{n}{callbacks}\PY{o}{=}\PY{p}{[}\PY{n}{ErrorTracker}\PY{p}{(}\PY{p}{)}\PY{p}{]}\PY{p}{)}
\PY{n}{cifar\PYZus{}trainer}\PY{o}{.}\PY{n}{fit}\PY{p}{(}\PY{n}{cifar\PYZus{}module}\PY{p}{,}
                  \PY{n}{datamodule}\PY{o}{=}\PY{n}{cifar\PYZus{}dm}\PY{p}{)}
\end{Verbatim}
\end{tcolorbox}

    \begin{Verbatim}[commandchars=\\\{\}]
GPU available: True (mps), used: True
TPU available: False, using: 0 TPU cores
IPU available: False, using: 0 IPUs
HPU available: False, using: 0 HPUs

  | Name  | Type             | Params
-------------------------------------------
0 | model | CIFARModel       | 964 K
1 | loss  | CrossEntropyLoss | 0
-------------------------------------------
964 K     Trainable params
0         Non-trainable params
964 K     Total params
3.858     Total estimated model params size (MB)
`Trainer.fit` stopped: `max\_epochs=30` reached.
    \end{Verbatim}

    This model can take 10 minutes or more to run and achieves about 42\%
accuracy on the test data. Although this is not terrible for 100-class
data (a random classifier gets 1\% accuracy), searching the web we see
results around 75\%. Typically it takes a lot of architecture carpentry,
fiddling with regularization, and time, to achieve such results.

Let's take a look at the validation and training accuracy across epochs.

    \begin{tcolorbox}[breakable, size=fbox, boxrule=1pt, pad at break*=1mm,colback=cellbackground, colframe=cellborder]
\prompt{In}{incolor}{59}{\boxspacing}
\begin{Verbatim}[commandchars=\\\{\}]
\PY{n}{log\PYZus{}path} \PY{o}{=} \PY{n}{cifar\PYZus{}logger}\PY{o}{.}\PY{n}{experiment}\PY{o}{.}\PY{n}{metrics\PYZus{}file\PYZus{}path}
\PY{n}{cifar\PYZus{}results} \PY{o}{=} \PY{n}{pd}\PY{o}{.}\PY{n}{read\PYZus{}csv}\PY{p}{(}\PY{n}{log\PYZus{}path}\PY{p}{)}
\PY{n}{fig}\PY{p}{,} \PY{n}{ax} \PY{o}{=} \PY{n}{subplots}\PY{p}{(}\PY{l+m+mi}{1}\PY{p}{,} \PY{l+m+mi}{1}\PY{p}{,} \PY{n}{figsize}\PY{o}{=}\PY{p}{(}\PY{l+m+mi}{6}\PY{p}{,} \PY{l+m+mi}{6}\PY{p}{)}\PY{p}{)}
\PY{n}{summary\PYZus{}plot}\PY{p}{(}\PY{n}{cifar\PYZus{}results}\PY{p}{,}
             \PY{n}{ax}\PY{p}{,}
             \PY{n}{col}\PY{o}{=}\PY{l+s+s1}{\PYZsq{}}\PY{l+s+s1}{accuracy}\PY{l+s+s1}{\PYZsq{}}\PY{p}{,}
             \PY{n}{ylabel}\PY{o}{=}\PY{l+s+s1}{\PYZsq{}}\PY{l+s+s1}{Accuracy}\PY{l+s+s1}{\PYZsq{}}\PY{p}{)}
\PY{n}{ax}\PY{o}{.}\PY{n}{set\PYZus{}xticks}\PY{p}{(}\PY{n}{np}\PY{o}{.}\PY{n}{linspace}\PY{p}{(}\PY{l+m+mi}{0}\PY{p}{,} \PY{l+m+mi}{10}\PY{p}{,} \PY{l+m+mi}{6}\PY{p}{)}\PY{o}{.}\PY{n}{astype}\PY{p}{(}\PY{n+nb}{int}\PY{p}{)}\PY{p}{)}
\PY{n}{ax}\PY{o}{.}\PY{n}{set\PYZus{}ylabel}\PY{p}{(}\PY{l+s+s1}{\PYZsq{}}\PY{l+s+s1}{Accuracy}\PY{l+s+s1}{\PYZsq{}}\PY{p}{)}
\PY{n}{ax}\PY{o}{.}\PY{n}{set\PYZus{}ylim}\PY{p}{(}\PY{p}{[}\PY{l+m+mi}{0}\PY{p}{,} \PY{l+m+mi}{1}\PY{p}{]}\PY{p}{)}\PY{p}{;}
\end{Verbatim}
\end{tcolorbox}

    \begin{center}
    \adjustimage{max size={0.9\linewidth}{0.9\paperheight}}{artifacts/latex/Ch10-deeplearning-lab_files/artifacts/latex/Ch10-deeplearning-lab_119_0.png}
    \end{center}
    { \hspace*{\fill} \\}
    
    Finally, we evaluate our model on our test data.

    \begin{tcolorbox}[breakable, size=fbox, boxrule=1pt, pad at break*=1mm,colback=cellbackground, colframe=cellborder]
\prompt{In}{incolor}{60}{\boxspacing}
\begin{Verbatim}[commandchars=\\\{\}]
\PY{n}{cifar\PYZus{}trainer}\PY{o}{.}\PY{n}{test}\PY{p}{(}\PY{n}{cifar\PYZus{}module}\PY{p}{,}
                   \PY{n}{datamodule}\PY{o}{=}\PY{n}{cifar\PYZus{}dm}\PY{p}{)}
\end{Verbatim}
\end{tcolorbox}

    \begin{Verbatim}[commandchars=\\\{\}]
────────────────────────────────────────────────────────────────────────────────
────────────────────────────────────────
       Test metric             DataLoader 0
────────────────────────────────────────────────────────────────────────────────
────────────────────────────────────────
      test\_accuracy         0.4269999861717224
        test\_loss            2.45865797996521
────────────────────────────────────────────────────────────────────────────────
────────────────────────────────────────
    \end{Verbatim}

            \begin{tcolorbox}[breakable, size=fbox, boxrule=.5pt, pad at break*=1mm, opacityfill=0]
\prompt{Out}{outcolor}{60}{\boxspacing}
\begin{Verbatim}[commandchars=\\\{\}]
[\{'test\_loss': 2.45865797996521, 'test\_accuracy': 0.4269999861717224\}]
\end{Verbatim}
\end{tcolorbox}
        
    \hypertarget{hardware-acceleration}{%
\subsubsection{Hardware Acceleration}\label{hardware-acceleration}}

As deep learning has become ubiquitous in machine learning, hardware
manufacturers have produced special libraries that can often speed up
the gradient-descent steps.

For instance, Mac OS devices with the M1 chip may have the \emph{Metal}
programming framework enabled, which can speed up the \texttt{torch}
computations. We present an example of how to use this acceleration.

The main changes are to the \texttt{Trainer()} call as well as to the
metrics that will be evaluated on the data. These metrics must be told
where the data will be located at evaluation time. This is accomplished
with a call to the \texttt{to()} method of the metrics.

    \begin{tcolorbox}[breakable, size=fbox, boxrule=1pt, pad at break*=1mm,colback=cellbackground, colframe=cellborder]
\prompt{In}{incolor}{61}{\boxspacing}
\begin{Verbatim}[commandchars=\\\{\}]
\PY{k}{try}\PY{p}{:}
    \PY{k}{for} \PY{n}{name}\PY{p}{,} \PY{n}{metric} \PY{o+ow}{in} \PY{n}{cifar\PYZus{}module}\PY{o}{.}\PY{n}{metrics}\PY{o}{.}\PY{n}{items}\PY{p}{(}\PY{p}{)}\PY{p}{:}
        \PY{n}{cifar\PYZus{}module}\PY{o}{.}\PY{n}{metrics}\PY{p}{[}\PY{n}{name}\PY{p}{]} \PY{o}{=} \PY{n}{metric}\PY{o}{.}\PY{n}{to}\PY{p}{(}\PY{l+s+s1}{\PYZsq{}}\PY{l+s+s1}{mps}\PY{l+s+s1}{\PYZsq{}}\PY{p}{)}
    \PY{n}{cifar\PYZus{}trainer\PYZus{}mps} \PY{o}{=} \PY{n}{Trainer}\PY{p}{(}\PY{n}{accelerator}\PY{o}{=}\PY{l+s+s1}{\PYZsq{}}\PY{l+s+s1}{mps}\PY{l+s+s1}{\PYZsq{}}\PY{p}{,}
                                \PY{n}{deterministic}\PY{o}{=}\PY{k+kc}{True}\PY{p}{,}
                                \PY{n}{enable\PYZus{}progress\PYZus{}bar}\PY{o}{=}\PY{k+kc}{False}\PY{p}{,}
                                \PY{n}{max\PYZus{}epochs}\PY{o}{=}\PY{l+m+mi}{30}\PY{p}{)}
    \PY{n}{cifar\PYZus{}trainer\PYZus{}mps}\PY{o}{.}\PY{n}{fit}\PY{p}{(}\PY{n}{cifar\PYZus{}module}\PY{p}{,}
                          \PY{n}{datamodule}\PY{o}{=}\PY{n}{cifar\PYZus{}dm}\PY{p}{)}
    \PY{n}{cifar\PYZus{}trainer\PYZus{}mps}\PY{o}{.}\PY{n}{test}\PY{p}{(}\PY{n}{cifar\PYZus{}module}\PY{p}{,}
                          \PY{n}{datamodule}\PY{o}{=}\PY{n}{cifar\PYZus{}dm}\PY{p}{)}
\PY{k}{except}\PY{p}{:}
    \PY{k}{pass}
\end{Verbatim}
\end{tcolorbox}

    \begin{Verbatim}[commandchars=\\\{\}]
GPU available: True (mps), used: True
TPU available: False, using: 0 TPU cores
IPU available: False, using: 0 IPUs
HPU available: False, using: 0 HPUs

  | Name  | Type             | Params
-------------------------------------------
0 | model | CIFARModel       | 964 K
1 | loss  | CrossEntropyLoss | 0
-------------------------------------------
964 K     Trainable params
0         Non-trainable params
964 K     Total params
3.858     Total estimated model params size (MB)
`Trainer.fit` stopped: `max\_epochs=30` reached.
    \end{Verbatim}

    This yields approximately two- or three-fold acceleration for each
epoch. We have protected this code block using \texttt{try:} and
\texttt{except:} clauses; if it works, we get the speedup, if it fails,
nothing happens.

    \hypertarget{using-pretrained-cnn-models}{%
\subsection{Using Pretrained CNN
Models}\label{using-pretrained-cnn-models}}

We now show how to use a CNN pretrained on the \texttt{imagenet}
database to classify natural images, and demonstrate how we produced
Figure 10.10. We copied six JPEG images from a digital photo album into
the directory \texttt{book\_images}. These images are available from the
data section of \textless www.statlearning.com\textgreater, the ISLP
book website. Download \texttt{book\_images.zip}; when clicked it
creates the \texttt{book\_images} directory.

The pretrained network we use is called \texttt{resnet50}; specification
details can be found on the web. We will read in the images, and convert
them into the array format expected by the \texttt{torch} software to
match the specifications in \texttt{resnet50}. The conversion involves a
resize, a crop and then a predefined standardization for each of the
three channels. We now read in the images and preprocess them.

    \begin{tcolorbox}[breakable, size=fbox, boxrule=1pt, pad at break*=1mm,colback=cellbackground, colframe=cellborder]
\prompt{In}{incolor}{62}{\boxspacing}
\begin{Verbatim}[commandchars=\\\{\}]
\PY{n}{resize} \PY{o}{=} \PY{n}{Resize}\PY{p}{(}\PY{p}{(}\PY{l+m+mi}{232}\PY{p}{,}\PY{l+m+mi}{232}\PY{p}{)}\PY{p}{,} \PY{n}{antialias}\PY{o}{=}\PY{k+kc}{True}\PY{p}{)}
\PY{n}{crop} \PY{o}{=} \PY{n}{CenterCrop}\PY{p}{(}\PY{l+m+mi}{224}\PY{p}{)}
\PY{n}{normalize} \PY{o}{=} \PY{n}{Normalize}\PY{p}{(}\PY{p}{[}\PY{l+m+mf}{0.485}\PY{p}{,}\PY{l+m+mf}{0.456}\PY{p}{,}\PY{l+m+mf}{0.406}\PY{p}{]}\PY{p}{,}
                      \PY{p}{[}\PY{l+m+mf}{0.229}\PY{p}{,}\PY{l+m+mf}{0.224}\PY{p}{,}\PY{l+m+mf}{0.225}\PY{p}{]}\PY{p}{)}
\PY{n}{imgfiles} \PY{o}{=} \PY{n+nb}{sorted}\PY{p}{(}\PY{p}{[}\PY{n}{f} \PY{k}{for} \PY{n}{f} \PY{o+ow}{in} \PY{n}{glob}\PY{p}{(}\PY{l+s+s1}{\PYZsq{}}\PY{l+s+s1}{book\PYZus{}images/*}\PY{l+s+s1}{\PYZsq{}}\PY{p}{)}\PY{p}{]}\PY{p}{)}
\PY{n}{imgs} \PY{o}{=} \PY{n}{torch}\PY{o}{.}\PY{n}{stack}\PY{p}{(}\PY{p}{[}\PY{n}{torch}\PY{o}{.}\PY{n}{div}\PY{p}{(}\PY{n}{crop}\PY{p}{(}\PY{n}{resize}\PY{p}{(}\PY{n}{read\PYZus{}image}\PY{p}{(}\PY{n}{f}\PY{p}{)}\PY{p}{)}\PY{p}{)}\PY{p}{,} \PY{l+m+mi}{255}\PY{p}{)}
                    \PY{k}{for} \PY{n}{f} \PY{o+ow}{in} \PY{n}{imgfiles}\PY{p}{]}\PY{p}{)}
\PY{n}{imgs} \PY{o}{=} \PY{n}{normalize}\PY{p}{(}\PY{n}{imgs}\PY{p}{)}
\PY{n}{imgs}\PY{o}{.}\PY{n}{size}\PY{p}{(}\PY{p}{)}
\end{Verbatim}
\end{tcolorbox}

            \begin{tcolorbox}[breakable, size=fbox, boxrule=.5pt, pad at break*=1mm, opacityfill=0]
\prompt{Out}{outcolor}{62}{\boxspacing}
\begin{Verbatim}[commandchars=\\\{\}]
torch.Size([6, 3, 224, 224])
\end{Verbatim}
\end{tcolorbox}
        
    We now set up the trained network with the weights we read in code
block\textasciitilde6. The model has 50 layers, with a fair bit of
complexity.

    \begin{tcolorbox}[breakable, size=fbox, boxrule=1pt, pad at break*=1mm,colback=cellbackground, colframe=cellborder]
\prompt{In}{incolor}{63}{\boxspacing}
\begin{Verbatim}[commandchars=\\\{\}]
\PY{n}{resnet\PYZus{}model} \PY{o}{=} \PY{n}{resnet50}\PY{p}{(}\PY{n}{weights}\PY{o}{=}\PY{n}{ResNet50\PYZus{}Weights}\PY{o}{.}\PY{n}{DEFAULT}\PY{p}{)}
\PY{n}{summary}\PY{p}{(}\PY{n}{resnet\PYZus{}model}\PY{p}{,}
        \PY{n}{input\PYZus{}data}\PY{o}{=}\PY{n}{imgs}\PY{p}{,}
        \PY{n}{col\PYZus{}names}\PY{o}{=}\PY{p}{[}\PY{l+s+s1}{\PYZsq{}}\PY{l+s+s1}{input\PYZus{}size}\PY{l+s+s1}{\PYZsq{}}\PY{p}{,}
                   \PY{l+s+s1}{\PYZsq{}}\PY{l+s+s1}{output\PYZus{}size}\PY{l+s+s1}{\PYZsq{}}\PY{p}{,}
                   \PY{l+s+s1}{\PYZsq{}}\PY{l+s+s1}{num\PYZus{}params}\PY{l+s+s1}{\PYZsq{}}\PY{p}{]}\PY{p}{)}
\end{Verbatim}
\end{tcolorbox}

            \begin{tcolorbox}[breakable, size=fbox, boxrule=.5pt, pad at break*=1mm, opacityfill=0]
\prompt{Out}{outcolor}{63}{\boxspacing}
\begin{Verbatim}[commandchars=\\\{\}]
================================================================================
===================================
Layer (type:depth-idx)                   Input Shape               Output Shape
Param \#
================================================================================
===================================
ResNet                                   [6, 3, 224, 224]          [6, 1000]
--
├─Conv2d: 1-1                            [6, 3, 224, 224]          [6, 64, 112,
112]         9,408
├─BatchNorm2d: 1-2                       [6, 64, 112, 112]         [6, 64, 112,
112]         128
├─ReLU: 1-3                              [6, 64, 112, 112]         [6, 64, 112,
112]         --
├─MaxPool2d: 1-4                         [6, 64, 112, 112]         [6, 64, 56,
56]           --
├─Sequential: 1-5                        [6, 64, 56, 56]           [6, 256, 56,
56]          --
│    └─Bottleneck: 2-1                   [6, 64, 56, 56]           [6, 256, 56,
56]          --
│    │    └─Conv2d: 3-1                  [6, 64, 56, 56]           [6, 64, 56,
56]           4,096
│    │    └─BatchNorm2d: 3-2             [6, 64, 56, 56]           [6, 64, 56,
56]           128
│    │    └─ReLU: 3-3                    [6, 64, 56, 56]           [6, 64, 56,
56]           --
│    │    └─Conv2d: 3-4                  [6, 64, 56, 56]           [6, 64, 56,
56]           36,864
│    │    └─BatchNorm2d: 3-5             [6, 64, 56, 56]           [6, 64, 56,
56]           128
│    │    └─ReLU: 3-6                    [6, 64, 56, 56]           [6, 64, 56,
56]           --
│    │    └─Conv2d: 3-7                  [6, 64, 56, 56]           [6, 256, 56,
56]          16,384
│    │    └─BatchNorm2d: 3-8             [6, 256, 56, 56]          [6, 256, 56,
56]          512
│    │    └─Sequential: 3-9              [6, 64, 56, 56]           [6, 256, 56,
56]          16,896
│    │    └─ReLU: 3-10                   [6, 256, 56, 56]          [6, 256, 56,
56]          --
│    └─Bottleneck: 2-2                   [6, 256, 56, 56]          [6, 256, 56,
56]          --
│    │    └─Conv2d: 3-11                 [6, 256, 56, 56]          [6, 64, 56,
56]           16,384
│    │    └─BatchNorm2d: 3-12            [6, 64, 56, 56]           [6, 64, 56,
56]           128
│    │    └─ReLU: 3-13                   [6, 64, 56, 56]           [6, 64, 56,
56]           --
│    │    └─Conv2d: 3-14                 [6, 64, 56, 56]           [6, 64, 56,
56]           36,864
│    │    └─BatchNorm2d: 3-15            [6, 64, 56, 56]           [6, 64, 56,
56]           128
│    │    └─ReLU: 3-16                   [6, 64, 56, 56]           [6, 64, 56,
56]           --
│    │    └─Conv2d: 3-17                 [6, 64, 56, 56]           [6, 256, 56,
56]          16,384
│    │    └─BatchNorm2d: 3-18            [6, 256, 56, 56]          [6, 256, 56,
56]          512
│    │    └─ReLU: 3-19                   [6, 256, 56, 56]          [6, 256, 56,
56]          --
│    └─Bottleneck: 2-3                   [6, 256, 56, 56]          [6, 256, 56,
56]          --
│    │    └─Conv2d: 3-20                 [6, 256, 56, 56]          [6, 64, 56,
56]           16,384
│    │    └─BatchNorm2d: 3-21            [6, 64, 56, 56]           [6, 64, 56,
56]           128
│    │    └─ReLU: 3-22                   [6, 64, 56, 56]           [6, 64, 56,
56]           --
│    │    └─Conv2d: 3-23                 [6, 64, 56, 56]           [6, 64, 56,
56]           36,864
│    │    └─BatchNorm2d: 3-24            [6, 64, 56, 56]           [6, 64, 56,
56]           128
│    │    └─ReLU: 3-25                   [6, 64, 56, 56]           [6, 64, 56,
56]           --
│    │    └─Conv2d: 3-26                 [6, 64, 56, 56]           [6, 256, 56,
56]          16,384
│    │    └─BatchNorm2d: 3-27            [6, 256, 56, 56]          [6, 256, 56,
56]          512
│    │    └─ReLU: 3-28                   [6, 256, 56, 56]          [6, 256, 56,
56]          --
├─Sequential: 1-6                        [6, 256, 56, 56]          [6, 512, 28,
28]          --
│    └─Bottleneck: 2-4                   [6, 256, 56, 56]          [6, 512, 28,
28]          --
│    │    └─Conv2d: 3-29                 [6, 256, 56, 56]          [6, 128, 56,
56]          32,768
│    │    └─BatchNorm2d: 3-30            [6, 128, 56, 56]          [6, 128, 56,
56]          256
│    │    └─ReLU: 3-31                   [6, 128, 56, 56]          [6, 128, 56,
56]          --
│    │    └─Conv2d: 3-32                 [6, 128, 56, 56]          [6, 128, 28,
28]          147,456
│    │    └─BatchNorm2d: 3-33            [6, 128, 28, 28]          [6, 128, 28,
28]          256
│    │    └─ReLU: 3-34                   [6, 128, 28, 28]          [6, 128, 28,
28]          --
│    │    └─Conv2d: 3-35                 [6, 128, 28, 28]          [6, 512, 28,
28]          65,536
│    │    └─BatchNorm2d: 3-36            [6, 512, 28, 28]          [6, 512, 28,
28]          1,024
│    │    └─Sequential: 3-37             [6, 256, 56, 56]          [6, 512, 28,
28]          132,096
│    │    └─ReLU: 3-38                   [6, 512, 28, 28]          [6, 512, 28,
28]          --
│    └─Bottleneck: 2-5                   [6, 512, 28, 28]          [6, 512, 28,
28]          --
│    │    └─Conv2d: 3-39                 [6, 512, 28, 28]          [6, 128, 28,
28]          65,536
│    │    └─BatchNorm2d: 3-40            [6, 128, 28, 28]          [6, 128, 28,
28]          256
│    │    └─ReLU: 3-41                   [6, 128, 28, 28]          [6, 128, 28,
28]          --
│    │    └─Conv2d: 3-42                 [6, 128, 28, 28]          [6, 128, 28,
28]          147,456
│    │    └─BatchNorm2d: 3-43            [6, 128, 28, 28]          [6, 128, 28,
28]          256
│    │    └─ReLU: 3-44                   [6, 128, 28, 28]          [6, 128, 28,
28]          --
│    │    └─Conv2d: 3-45                 [6, 128, 28, 28]          [6, 512, 28,
28]          65,536
│    │    └─BatchNorm2d: 3-46            [6, 512, 28, 28]          [6, 512, 28,
28]          1,024
│    │    └─ReLU: 3-47                   [6, 512, 28, 28]          [6, 512, 28,
28]          --
│    └─Bottleneck: 2-6                   [6, 512, 28, 28]          [6, 512, 28,
28]          --
│    │    └─Conv2d: 3-48                 [6, 512, 28, 28]          [6, 128, 28,
28]          65,536
│    │    └─BatchNorm2d: 3-49            [6, 128, 28, 28]          [6, 128, 28,
28]          256
│    │    └─ReLU: 3-50                   [6, 128, 28, 28]          [6, 128, 28,
28]          --
│    │    └─Conv2d: 3-51                 [6, 128, 28, 28]          [6, 128, 28,
28]          147,456
│    │    └─BatchNorm2d: 3-52            [6, 128, 28, 28]          [6, 128, 28,
28]          256
│    │    └─ReLU: 3-53                   [6, 128, 28, 28]          [6, 128, 28,
28]          --
│    │    └─Conv2d: 3-54                 [6, 128, 28, 28]          [6, 512, 28,
28]          65,536
│    │    └─BatchNorm2d: 3-55            [6, 512, 28, 28]          [6, 512, 28,
28]          1,024
│    │    └─ReLU: 3-56                   [6, 512, 28, 28]          [6, 512, 28,
28]          --
│    └─Bottleneck: 2-7                   [6, 512, 28, 28]          [6, 512, 28,
28]          --
│    │    └─Conv2d: 3-57                 [6, 512, 28, 28]          [6, 128, 28,
28]          65,536
│    │    └─BatchNorm2d: 3-58            [6, 128, 28, 28]          [6, 128, 28,
28]          256
│    │    └─ReLU: 3-59                   [6, 128, 28, 28]          [6, 128, 28,
28]          --
│    │    └─Conv2d: 3-60                 [6, 128, 28, 28]          [6, 128, 28,
28]          147,456
│    │    └─BatchNorm2d: 3-61            [6, 128, 28, 28]          [6, 128, 28,
28]          256
│    │    └─ReLU: 3-62                   [6, 128, 28, 28]          [6, 128, 28,
28]          --
│    │    └─Conv2d: 3-63                 [6, 128, 28, 28]          [6, 512, 28,
28]          65,536
│    │    └─BatchNorm2d: 3-64            [6, 512, 28, 28]          [6, 512, 28,
28]          1,024
│    │    └─ReLU: 3-65                   [6, 512, 28, 28]          [6, 512, 28,
28]          --
├─Sequential: 1-7                        [6, 512, 28, 28]          [6, 1024, 14,
14]         --
│    └─Bottleneck: 2-8                   [6, 512, 28, 28]          [6, 1024, 14,
14]         --
│    │    └─Conv2d: 3-66                 [6, 512, 28, 28]          [6, 256, 28,
28]          131,072
│    │    └─BatchNorm2d: 3-67            [6, 256, 28, 28]          [6, 256, 28,
28]          512
│    │    └─ReLU: 3-68                   [6, 256, 28, 28]          [6, 256, 28,
28]          --
│    │    └─Conv2d: 3-69                 [6, 256, 28, 28]          [6, 256, 14,
14]          589,824
│    │    └─BatchNorm2d: 3-70            [6, 256, 14, 14]          [6, 256, 14,
14]          512
│    │    └─ReLU: 3-71                   [6, 256, 14, 14]          [6, 256, 14,
14]          --
│    │    └─Conv2d: 3-72                 [6, 256, 14, 14]          [6, 1024, 14,
14]         262,144
│    │    └─BatchNorm2d: 3-73            [6, 1024, 14, 14]         [6, 1024, 14,
14]         2,048
│    │    └─Sequential: 3-74             [6, 512, 28, 28]          [6, 1024, 14,
14]         526,336
│    │    └─ReLU: 3-75                   [6, 1024, 14, 14]         [6, 1024, 14,
14]         --
│    └─Bottleneck: 2-9                   [6, 1024, 14, 14]         [6, 1024, 14,
14]         --
│    │    └─Conv2d: 3-76                 [6, 1024, 14, 14]         [6, 256, 14,
14]          262,144
│    │    └─BatchNorm2d: 3-77            [6, 256, 14, 14]          [6, 256, 14,
14]          512
│    │    └─ReLU: 3-78                   [6, 256, 14, 14]          [6, 256, 14,
14]          --
│    │    └─Conv2d: 3-79                 [6, 256, 14, 14]          [6, 256, 14,
14]          589,824
│    │    └─BatchNorm2d: 3-80            [6, 256, 14, 14]          [6, 256, 14,
14]          512
│    │    └─ReLU: 3-81                   [6, 256, 14, 14]          [6, 256, 14,
14]          --
│    │    └─Conv2d: 3-82                 [6, 256, 14, 14]          [6, 1024, 14,
14]         262,144
│    │    └─BatchNorm2d: 3-83            [6, 1024, 14, 14]         [6, 1024, 14,
14]         2,048
│    │    └─ReLU: 3-84                   [6, 1024, 14, 14]         [6, 1024, 14,
14]         --
│    └─Bottleneck: 2-10                  [6, 1024, 14, 14]         [6, 1024, 14,
14]         --
│    │    └─Conv2d: 3-85                 [6, 1024, 14, 14]         [6, 256, 14,
14]          262,144
│    │    └─BatchNorm2d: 3-86            [6, 256, 14, 14]          [6, 256, 14,
14]          512
│    │    └─ReLU: 3-87                   [6, 256, 14, 14]          [6, 256, 14,
14]          --
│    │    └─Conv2d: 3-88                 [6, 256, 14, 14]          [6, 256, 14,
14]          589,824
│    │    └─BatchNorm2d: 3-89            [6, 256, 14, 14]          [6, 256, 14,
14]          512
│    │    └─ReLU: 3-90                   [6, 256, 14, 14]          [6, 256, 14,
14]          --
│    │    └─Conv2d: 3-91                 [6, 256, 14, 14]          [6, 1024, 14,
14]         262,144
│    │    └─BatchNorm2d: 3-92            [6, 1024, 14, 14]         [6, 1024, 14,
14]         2,048
│    │    └─ReLU: 3-93                   [6, 1024, 14, 14]         [6, 1024, 14,
14]         --
│    └─Bottleneck: 2-11                  [6, 1024, 14, 14]         [6, 1024, 14,
14]         --
│    │    └─Conv2d: 3-94                 [6, 1024, 14, 14]         [6, 256, 14,
14]          262,144
│    │    └─BatchNorm2d: 3-95            [6, 256, 14, 14]          [6, 256, 14,
14]          512
│    │    └─ReLU: 3-96                   [6, 256, 14, 14]          [6, 256, 14,
14]          --
│    │    └─Conv2d: 3-97                 [6, 256, 14, 14]          [6, 256, 14,
14]          589,824
│    │    └─BatchNorm2d: 3-98            [6, 256, 14, 14]          [6, 256, 14,
14]          512
│    │    └─ReLU: 3-99                   [6, 256, 14, 14]          [6, 256, 14,
14]          --
│    │    └─Conv2d: 3-100                [6, 256, 14, 14]          [6, 1024, 14,
14]         262,144
│    │    └─BatchNorm2d: 3-101           [6, 1024, 14, 14]         [6, 1024, 14,
14]         2,048
│    │    └─ReLU: 3-102                  [6, 1024, 14, 14]         [6, 1024, 14,
14]         --
│    └─Bottleneck: 2-12                  [6, 1024, 14, 14]         [6, 1024, 14,
14]         --
│    │    └─Conv2d: 3-103                [6, 1024, 14, 14]         [6, 256, 14,
14]          262,144
│    │    └─BatchNorm2d: 3-104           [6, 256, 14, 14]          [6, 256, 14,
14]          512
│    │    └─ReLU: 3-105                  [6, 256, 14, 14]          [6, 256, 14,
14]          --
│    │    └─Conv2d: 3-106                [6, 256, 14, 14]          [6, 256, 14,
14]          589,824
│    │    └─BatchNorm2d: 3-107           [6, 256, 14, 14]          [6, 256, 14,
14]          512
│    │    └─ReLU: 3-108                  [6, 256, 14, 14]          [6, 256, 14,
14]          --
│    │    └─Conv2d: 3-109                [6, 256, 14, 14]          [6, 1024, 14,
14]         262,144
│    │    └─BatchNorm2d: 3-110           [6, 1024, 14, 14]         [6, 1024, 14,
14]         2,048
│    │    └─ReLU: 3-111                  [6, 1024, 14, 14]         [6, 1024, 14,
14]         --
│    └─Bottleneck: 2-13                  [6, 1024, 14, 14]         [6, 1024, 14,
14]         --
│    │    └─Conv2d: 3-112                [6, 1024, 14, 14]         [6, 256, 14,
14]          262,144
│    │    └─BatchNorm2d: 3-113           [6, 256, 14, 14]          [6, 256, 14,
14]          512
│    │    └─ReLU: 3-114                  [6, 256, 14, 14]          [6, 256, 14,
14]          --
│    │    └─Conv2d: 3-115                [6, 256, 14, 14]          [6, 256, 14,
14]          589,824
│    │    └─BatchNorm2d: 3-116           [6, 256, 14, 14]          [6, 256, 14,
14]          512
│    │    └─ReLU: 3-117                  [6, 256, 14, 14]          [6, 256, 14,
14]          --
│    │    └─Conv2d: 3-118                [6, 256, 14, 14]          [6, 1024, 14,
14]         262,144
│    │    └─BatchNorm2d: 3-119           [6, 1024, 14, 14]         [6, 1024, 14,
14]         2,048
│    │    └─ReLU: 3-120                  [6, 1024, 14, 14]         [6, 1024, 14,
14]         --
├─Sequential: 1-8                        [6, 1024, 14, 14]         [6, 2048, 7,
7]           --
│    └─Bottleneck: 2-14                  [6, 1024, 14, 14]         [6, 2048, 7,
7]           --
│    │    └─Conv2d: 3-121                [6, 1024, 14, 14]         [6, 512, 14,
14]          524,288
│    │    └─BatchNorm2d: 3-122           [6, 512, 14, 14]          [6, 512, 14,
14]          1,024
│    │    └─ReLU: 3-123                  [6, 512, 14, 14]          [6, 512, 14,
14]          --
│    │    └─Conv2d: 3-124                [6, 512, 14, 14]          [6, 512, 7,
7]            2,359,296
│    │    └─BatchNorm2d: 3-125           [6, 512, 7, 7]            [6, 512, 7,
7]            1,024
│    │    └─ReLU: 3-126                  [6, 512, 7, 7]            [6, 512, 7,
7]            --
│    │    └─Conv2d: 3-127                [6, 512, 7, 7]            [6, 2048, 7,
7]           1,048,576
│    │    └─BatchNorm2d: 3-128           [6, 2048, 7, 7]           [6, 2048, 7,
7]           4,096
│    │    └─Sequential: 3-129            [6, 1024, 14, 14]         [6, 2048, 7,
7]           2,101,248
│    │    └─ReLU: 3-130                  [6, 2048, 7, 7]           [6, 2048, 7,
7]           --
│    └─Bottleneck: 2-15                  [6, 2048, 7, 7]           [6, 2048, 7,
7]           --
│    │    └─Conv2d: 3-131                [6, 2048, 7, 7]           [6, 512, 7,
7]            1,048,576
│    │    └─BatchNorm2d: 3-132           [6, 512, 7, 7]            [6, 512, 7,
7]            1,024
│    │    └─ReLU: 3-133                  [6, 512, 7, 7]            [6, 512, 7,
7]            --
│    │    └─Conv2d: 3-134                [6, 512, 7, 7]            [6, 512, 7,
7]            2,359,296
│    │    └─BatchNorm2d: 3-135           [6, 512, 7, 7]            [6, 512, 7,
7]            1,024
│    │    └─ReLU: 3-136                  [6, 512, 7, 7]            [6, 512, 7,
7]            --
│    │    └─Conv2d: 3-137                [6, 512, 7, 7]            [6, 2048, 7,
7]           1,048,576
│    │    └─BatchNorm2d: 3-138           [6, 2048, 7, 7]           [6, 2048, 7,
7]           4,096
│    │    └─ReLU: 3-139                  [6, 2048, 7, 7]           [6, 2048, 7,
7]           --
│    └─Bottleneck: 2-16                  [6, 2048, 7, 7]           [6, 2048, 7,
7]           --
│    │    └─Conv2d: 3-140                [6, 2048, 7, 7]           [6, 512, 7,
7]            1,048,576
│    │    └─BatchNorm2d: 3-141           [6, 512, 7, 7]            [6, 512, 7,
7]            1,024
│    │    └─ReLU: 3-142                  [6, 512, 7, 7]            [6, 512, 7,
7]            --
│    │    └─Conv2d: 3-143                [6, 512, 7, 7]            [6, 512, 7,
7]            2,359,296
│    │    └─BatchNorm2d: 3-144           [6, 512, 7, 7]            [6, 512, 7,
7]            1,024
│    │    └─ReLU: 3-145                  [6, 512, 7, 7]            [6, 512, 7,
7]            --
│    │    └─Conv2d: 3-146                [6, 512, 7, 7]            [6, 2048, 7,
7]           1,048,576
│    │    └─BatchNorm2d: 3-147           [6, 2048, 7, 7]           [6, 2048, 7,
7]           4,096
│    │    └─ReLU: 3-148                  [6, 2048, 7, 7]           [6, 2048, 7,
7]           --
├─AdaptiveAvgPool2d: 1-9                 [6, 2048, 7, 7]           [6, 2048, 1,
1]           --
├─Linear: 1-10                           [6, 2048]                 [6, 1000]
2,049,000
================================================================================
===================================
Total params: 25,557,032
Trainable params: 25,557,032
Non-trainable params: 0
Total mult-adds (Units.GIGABYTES): 24.54
================================================================================
===================================
Input size (MB): 3.61
Forward/backward pass size (MB): 1066.99
Params size (MB): 102.23
Estimated Total Size (MB): 1172.83
================================================================================
===================================
\end{Verbatim}
\end{tcolorbox}
        
    We set the mode to \texttt{eval()} to ensure that the model is ready to
predict on new data.

    \begin{tcolorbox}[breakable, size=fbox, boxrule=1pt, pad at break*=1mm,colback=cellbackground, colframe=cellborder]
\prompt{In}{incolor}{64}{\boxspacing}
\begin{Verbatim}[commandchars=\\\{\}]
\PY{n}{resnet\PYZus{}model}\PY{o}{.}\PY{n}{eval}\PY{p}{(}\PY{p}{)}
\end{Verbatim}
\end{tcolorbox}

            \begin{tcolorbox}[breakable, size=fbox, boxrule=.5pt, pad at break*=1mm, opacityfill=0]
\prompt{Out}{outcolor}{64}{\boxspacing}
\begin{Verbatim}[commandchars=\\\{\}]
ResNet(
  (conv1): Conv2d(3, 64, kernel\_size=(7, 7), stride=(2, 2), padding=(3, 3),
bias=False)
  (bn1): BatchNorm2d(64, eps=1e-05, momentum=0.1, affine=True,
track\_running\_stats=True)
  (relu): ReLU(inplace=True)
  (maxpool): MaxPool2d(kernel\_size=3, stride=2, padding=1, dilation=1,
ceil\_mode=False)
  (layer1): Sequential(
    (0): Bottleneck(
      (conv1): Conv2d(64, 64, kernel\_size=(1, 1), stride=(1, 1), bias=False)
      (bn1): BatchNorm2d(64, eps=1e-05, momentum=0.1, affine=True,
track\_running\_stats=True)
      (conv2): Conv2d(64, 64, kernel\_size=(3, 3), stride=(1, 1), padding=(1, 1),
bias=False)
      (bn2): BatchNorm2d(64, eps=1e-05, momentum=0.1, affine=True,
track\_running\_stats=True)
      (conv3): Conv2d(64, 256, kernel\_size=(1, 1), stride=(1, 1), bias=False)
      (bn3): BatchNorm2d(256, eps=1e-05, momentum=0.1, affine=True,
track\_running\_stats=True)
      (relu): ReLU(inplace=True)
      (downsample): Sequential(
        (0): Conv2d(64, 256, kernel\_size=(1, 1), stride=(1, 1), bias=False)
        (1): BatchNorm2d(256, eps=1e-05, momentum=0.1, affine=True,
track\_running\_stats=True)
      )
    )
    (1): Bottleneck(
      (conv1): Conv2d(256, 64, kernel\_size=(1, 1), stride=(1, 1), bias=False)
      (bn1): BatchNorm2d(64, eps=1e-05, momentum=0.1, affine=True,
track\_running\_stats=True)
      (conv2): Conv2d(64, 64, kernel\_size=(3, 3), stride=(1, 1), padding=(1, 1),
bias=False)
      (bn2): BatchNorm2d(64, eps=1e-05, momentum=0.1, affine=True,
track\_running\_stats=True)
      (conv3): Conv2d(64, 256, kernel\_size=(1, 1), stride=(1, 1), bias=False)
      (bn3): BatchNorm2d(256, eps=1e-05, momentum=0.1, affine=True,
track\_running\_stats=True)
      (relu): ReLU(inplace=True)
    )
    (2): Bottleneck(
      (conv1): Conv2d(256, 64, kernel\_size=(1, 1), stride=(1, 1), bias=False)
      (bn1): BatchNorm2d(64, eps=1e-05, momentum=0.1, affine=True,
track\_running\_stats=True)
      (conv2): Conv2d(64, 64, kernel\_size=(3, 3), stride=(1, 1), padding=(1, 1),
bias=False)
      (bn2): BatchNorm2d(64, eps=1e-05, momentum=0.1, affine=True,
track\_running\_stats=True)
      (conv3): Conv2d(64, 256, kernel\_size=(1, 1), stride=(1, 1), bias=False)
      (bn3): BatchNorm2d(256, eps=1e-05, momentum=0.1, affine=True,
track\_running\_stats=True)
      (relu): ReLU(inplace=True)
    )
  )
  (layer2): Sequential(
    (0): Bottleneck(
      (conv1): Conv2d(256, 128, kernel\_size=(1, 1), stride=(1, 1), bias=False)
      (bn1): BatchNorm2d(128, eps=1e-05, momentum=0.1, affine=True,
track\_running\_stats=True)
      (conv2): Conv2d(128, 128, kernel\_size=(3, 3), stride=(2, 2), padding=(1,
1), bias=False)
      (bn2): BatchNorm2d(128, eps=1e-05, momentum=0.1, affine=True,
track\_running\_stats=True)
      (conv3): Conv2d(128, 512, kernel\_size=(1, 1), stride=(1, 1), bias=False)
      (bn3): BatchNorm2d(512, eps=1e-05, momentum=0.1, affine=True,
track\_running\_stats=True)
      (relu): ReLU(inplace=True)
      (downsample): Sequential(
        (0): Conv2d(256, 512, kernel\_size=(1, 1), stride=(2, 2), bias=False)
        (1): BatchNorm2d(512, eps=1e-05, momentum=0.1, affine=True,
track\_running\_stats=True)
      )
    )
    (1): Bottleneck(
      (conv1): Conv2d(512, 128, kernel\_size=(1, 1), stride=(1, 1), bias=False)
      (bn1): BatchNorm2d(128, eps=1e-05, momentum=0.1, affine=True,
track\_running\_stats=True)
      (conv2): Conv2d(128, 128, kernel\_size=(3, 3), stride=(1, 1), padding=(1,
1), bias=False)
      (bn2): BatchNorm2d(128, eps=1e-05, momentum=0.1, affine=True,
track\_running\_stats=True)
      (conv3): Conv2d(128, 512, kernel\_size=(1, 1), stride=(1, 1), bias=False)
      (bn3): BatchNorm2d(512, eps=1e-05, momentum=0.1, affine=True,
track\_running\_stats=True)
      (relu): ReLU(inplace=True)
    )
    (2): Bottleneck(
      (conv1): Conv2d(512, 128, kernel\_size=(1, 1), stride=(1, 1), bias=False)
      (bn1): BatchNorm2d(128, eps=1e-05, momentum=0.1, affine=True,
track\_running\_stats=True)
      (conv2): Conv2d(128, 128, kernel\_size=(3, 3), stride=(1, 1), padding=(1,
1), bias=False)
      (bn2): BatchNorm2d(128, eps=1e-05, momentum=0.1, affine=True,
track\_running\_stats=True)
      (conv3): Conv2d(128, 512, kernel\_size=(1, 1), stride=(1, 1), bias=False)
      (bn3): BatchNorm2d(512, eps=1e-05, momentum=0.1, affine=True,
track\_running\_stats=True)
      (relu): ReLU(inplace=True)
    )
    (3): Bottleneck(
      (conv1): Conv2d(512, 128, kernel\_size=(1, 1), stride=(1, 1), bias=False)
      (bn1): BatchNorm2d(128, eps=1e-05, momentum=0.1, affine=True,
track\_running\_stats=True)
      (conv2): Conv2d(128, 128, kernel\_size=(3, 3), stride=(1, 1), padding=(1,
1), bias=False)
      (bn2): BatchNorm2d(128, eps=1e-05, momentum=0.1, affine=True,
track\_running\_stats=True)
      (conv3): Conv2d(128, 512, kernel\_size=(1, 1), stride=(1, 1), bias=False)
      (bn3): BatchNorm2d(512, eps=1e-05, momentum=0.1, affine=True,
track\_running\_stats=True)
      (relu): ReLU(inplace=True)
    )
  )
  (layer3): Sequential(
    (0): Bottleneck(
      (conv1): Conv2d(512, 256, kernel\_size=(1, 1), stride=(1, 1), bias=False)
      (bn1): BatchNorm2d(256, eps=1e-05, momentum=0.1, affine=True,
track\_running\_stats=True)
      (conv2): Conv2d(256, 256, kernel\_size=(3, 3), stride=(2, 2), padding=(1,
1), bias=False)
      (bn2): BatchNorm2d(256, eps=1e-05, momentum=0.1, affine=True,
track\_running\_stats=True)
      (conv3): Conv2d(256, 1024, kernel\_size=(1, 1), stride=(1, 1), bias=False)
      (bn3): BatchNorm2d(1024, eps=1e-05, momentum=0.1, affine=True,
track\_running\_stats=True)
      (relu): ReLU(inplace=True)
      (downsample): Sequential(
        (0): Conv2d(512, 1024, kernel\_size=(1, 1), stride=(2, 2), bias=False)
        (1): BatchNorm2d(1024, eps=1e-05, momentum=0.1, affine=True,
track\_running\_stats=True)
      )
    )
    (1): Bottleneck(
      (conv1): Conv2d(1024, 256, kernel\_size=(1, 1), stride=(1, 1), bias=False)
      (bn1): BatchNorm2d(256, eps=1e-05, momentum=0.1, affine=True,
track\_running\_stats=True)
      (conv2): Conv2d(256, 256, kernel\_size=(3, 3), stride=(1, 1), padding=(1,
1), bias=False)
      (bn2): BatchNorm2d(256, eps=1e-05, momentum=0.1, affine=True,
track\_running\_stats=True)
      (conv3): Conv2d(256, 1024, kernel\_size=(1, 1), stride=(1, 1), bias=False)
      (bn3): BatchNorm2d(1024, eps=1e-05, momentum=0.1, affine=True,
track\_running\_stats=True)
      (relu): ReLU(inplace=True)
    )
    (2): Bottleneck(
      (conv1): Conv2d(1024, 256, kernel\_size=(1, 1), stride=(1, 1), bias=False)
      (bn1): BatchNorm2d(256, eps=1e-05, momentum=0.1, affine=True,
track\_running\_stats=True)
      (conv2): Conv2d(256, 256, kernel\_size=(3, 3), stride=(1, 1), padding=(1,
1), bias=False)
      (bn2): BatchNorm2d(256, eps=1e-05, momentum=0.1, affine=True,
track\_running\_stats=True)
      (conv3): Conv2d(256, 1024, kernel\_size=(1, 1), stride=(1, 1), bias=False)
      (bn3): BatchNorm2d(1024, eps=1e-05, momentum=0.1, affine=True,
track\_running\_stats=True)
      (relu): ReLU(inplace=True)
    )
    (3): Bottleneck(
      (conv1): Conv2d(1024, 256, kernel\_size=(1, 1), stride=(1, 1), bias=False)
      (bn1): BatchNorm2d(256, eps=1e-05, momentum=0.1, affine=True,
track\_running\_stats=True)
      (conv2): Conv2d(256, 256, kernel\_size=(3, 3), stride=(1, 1), padding=(1,
1), bias=False)
      (bn2): BatchNorm2d(256, eps=1e-05, momentum=0.1, affine=True,
track\_running\_stats=True)
      (conv3): Conv2d(256, 1024, kernel\_size=(1, 1), stride=(1, 1), bias=False)
      (bn3): BatchNorm2d(1024, eps=1e-05, momentum=0.1, affine=True,
track\_running\_stats=True)
      (relu): ReLU(inplace=True)
    )
    (4): Bottleneck(
      (conv1): Conv2d(1024, 256, kernel\_size=(1, 1), stride=(1, 1), bias=False)
      (bn1): BatchNorm2d(256, eps=1e-05, momentum=0.1, affine=True,
track\_running\_stats=True)
      (conv2): Conv2d(256, 256, kernel\_size=(3, 3), stride=(1, 1), padding=(1,
1), bias=False)
      (bn2): BatchNorm2d(256, eps=1e-05, momentum=0.1, affine=True,
track\_running\_stats=True)
      (conv3): Conv2d(256, 1024, kernel\_size=(1, 1), stride=(1, 1), bias=False)
      (bn3): BatchNorm2d(1024, eps=1e-05, momentum=0.1, affine=True,
track\_running\_stats=True)
      (relu): ReLU(inplace=True)
    )
    (5): Bottleneck(
      (conv1): Conv2d(1024, 256, kernel\_size=(1, 1), stride=(1, 1), bias=False)
      (bn1): BatchNorm2d(256, eps=1e-05, momentum=0.1, affine=True,
track\_running\_stats=True)
      (conv2): Conv2d(256, 256, kernel\_size=(3, 3), stride=(1, 1), padding=(1,
1), bias=False)
      (bn2): BatchNorm2d(256, eps=1e-05, momentum=0.1, affine=True,
track\_running\_stats=True)
      (conv3): Conv2d(256, 1024, kernel\_size=(1, 1), stride=(1, 1), bias=False)
      (bn3): BatchNorm2d(1024, eps=1e-05, momentum=0.1, affine=True,
track\_running\_stats=True)
      (relu): ReLU(inplace=True)
    )
  )
  (layer4): Sequential(
    (0): Bottleneck(
      (conv1): Conv2d(1024, 512, kernel\_size=(1, 1), stride=(1, 1), bias=False)
      (bn1): BatchNorm2d(512, eps=1e-05, momentum=0.1, affine=True,
track\_running\_stats=True)
      (conv2): Conv2d(512, 512, kernel\_size=(3, 3), stride=(2, 2), padding=(1,
1), bias=False)
      (bn2): BatchNorm2d(512, eps=1e-05, momentum=0.1, affine=True,
track\_running\_stats=True)
      (conv3): Conv2d(512, 2048, kernel\_size=(1, 1), stride=(1, 1), bias=False)
      (bn3): BatchNorm2d(2048, eps=1e-05, momentum=0.1, affine=True,
track\_running\_stats=True)
      (relu): ReLU(inplace=True)
      (downsample): Sequential(
        (0): Conv2d(1024, 2048, kernel\_size=(1, 1), stride=(2, 2), bias=False)
        (1): BatchNorm2d(2048, eps=1e-05, momentum=0.1, affine=True,
track\_running\_stats=True)
      )
    )
    (1): Bottleneck(
      (conv1): Conv2d(2048, 512, kernel\_size=(1, 1), stride=(1, 1), bias=False)
      (bn1): BatchNorm2d(512, eps=1e-05, momentum=0.1, affine=True,
track\_running\_stats=True)
      (conv2): Conv2d(512, 512, kernel\_size=(3, 3), stride=(1, 1), padding=(1,
1), bias=False)
      (bn2): BatchNorm2d(512, eps=1e-05, momentum=0.1, affine=True,
track\_running\_stats=True)
      (conv3): Conv2d(512, 2048, kernel\_size=(1, 1), stride=(1, 1), bias=False)
      (bn3): BatchNorm2d(2048, eps=1e-05, momentum=0.1, affine=True,
track\_running\_stats=True)
      (relu): ReLU(inplace=True)
    )
    (2): Bottleneck(
      (conv1): Conv2d(2048, 512, kernel\_size=(1, 1), stride=(1, 1), bias=False)
      (bn1): BatchNorm2d(512, eps=1e-05, momentum=0.1, affine=True,
track\_running\_stats=True)
      (conv2): Conv2d(512, 512, kernel\_size=(3, 3), stride=(1, 1), padding=(1,
1), bias=False)
      (bn2): BatchNorm2d(512, eps=1e-05, momentum=0.1, affine=True,
track\_running\_stats=True)
      (conv3): Conv2d(512, 2048, kernel\_size=(1, 1), stride=(1, 1), bias=False)
      (bn3): BatchNorm2d(2048, eps=1e-05, momentum=0.1, affine=True,
track\_running\_stats=True)
      (relu): ReLU(inplace=True)
    )
  )
  (avgpool): AdaptiveAvgPool2d(output\_size=(1, 1))
  (fc): Linear(in\_features=2048, out\_features=1000, bias=True)
)
\end{Verbatim}
\end{tcolorbox}
        
    Inspecting the output above, we see that when setting up the
\texttt{resnet\_model}, the authors defined a \texttt{Bottleneck}, much
like our \texttt{BuildingBlock} module.

We now feed our six images through the fitted network.

    \begin{tcolorbox}[breakable, size=fbox, boxrule=1pt, pad at break*=1mm,colback=cellbackground, colframe=cellborder]
\prompt{In}{incolor}{65}{\boxspacing}
\begin{Verbatim}[commandchars=\\\{\}]
\PY{n}{img\PYZus{}preds} \PY{o}{=} \PY{n}{resnet\PYZus{}model}\PY{p}{(}\PY{n}{imgs}\PY{p}{)}
\end{Verbatim}
\end{tcolorbox}

    Let's look at the predicted probabilities for each of the top 3 choices.
First we compute the probabilities by applying the softmax to the logits
in \texttt{img\_preds}. Note that we have had to call the
\texttt{detach()} method on the tensor \texttt{img\_preds} in order to
convert it to a more familiar \texttt{ndarray}.

    \begin{tcolorbox}[breakable, size=fbox, boxrule=1pt, pad at break*=1mm,colback=cellbackground, colframe=cellborder]
\prompt{In}{incolor}{66}{\boxspacing}
\begin{Verbatim}[commandchars=\\\{\}]
\PY{n}{img\PYZus{}probs} \PY{o}{=} \PY{n}{np}\PY{o}{.}\PY{n}{exp}\PY{p}{(}\PY{n}{np}\PY{o}{.}\PY{n}{asarray}\PY{p}{(}\PY{n}{img\PYZus{}preds}\PY{o}{.}\PY{n}{detach}\PY{p}{(}\PY{p}{)}\PY{p}{)}\PY{p}{)}
\PY{n}{img\PYZus{}probs} \PY{o}{/}\PY{o}{=} \PY{n}{img\PYZus{}probs}\PY{o}{.}\PY{n}{sum}\PY{p}{(}\PY{l+m+mi}{1}\PY{p}{)}\PY{p}{[}\PY{p}{:}\PY{p}{,}\PY{k+kc}{None}\PY{p}{]}
\end{Verbatim}
\end{tcolorbox}

    In order to see the class labels, we must download the index file
associated with \texttt{imagenet}. \{This is avalable from the book
website and
\href{https://s3.amazonaws.com/deep-learning-models/image-models/imagenet_class_index.json}{s3.amazonaws.com/deep-learning-models/image-models/imagenet\_class\_index.json}.\}

    \begin{tcolorbox}[breakable, size=fbox, boxrule=1pt, pad at break*=1mm,colback=cellbackground, colframe=cellborder]
\prompt{In}{incolor}{67}{\boxspacing}
\begin{Verbatim}[commandchars=\\\{\}]
\PY{n}{labs} \PY{o}{=} \PY{n}{json}\PY{o}{.}\PY{n}{load}\PY{p}{(}\PY{n+nb}{open}\PY{p}{(}\PY{l+s+s1}{\PYZsq{}}\PY{l+s+s1}{imagenet\PYZus{}class\PYZus{}index.json}\PY{l+s+s1}{\PYZsq{}}\PY{p}{)}\PY{p}{)}
\PY{n}{class\PYZus{}labels} \PY{o}{=} \PY{n}{pd}\PY{o}{.}\PY{n}{DataFrame}\PY{p}{(}\PY{p}{[}\PY{p}{(}\PY{n+nb}{int}\PY{p}{(}\PY{n}{k}\PY{p}{)}\PY{p}{,} \PY{n}{v}\PY{p}{[}\PY{l+m+mi}{1}\PY{p}{]}\PY{p}{)} \PY{k}{for} \PY{n}{k}\PY{p}{,} \PY{n}{v} \PY{o+ow}{in} 
                           \PY{n}{labs}\PY{o}{.}\PY{n}{items}\PY{p}{(}\PY{p}{)}\PY{p}{]}\PY{p}{,}
                           \PY{n}{columns}\PY{o}{=}\PY{p}{[}\PY{l+s+s1}{\PYZsq{}}\PY{l+s+s1}{idx}\PY{l+s+s1}{\PYZsq{}}\PY{p}{,} \PY{l+s+s1}{\PYZsq{}}\PY{l+s+s1}{label}\PY{l+s+s1}{\PYZsq{}}\PY{p}{]}\PY{p}{)}
\PY{n}{class\PYZus{}labels} \PY{o}{=} \PY{n}{class\PYZus{}labels}\PY{o}{.}\PY{n}{set\PYZus{}index}\PY{p}{(}\PY{l+s+s1}{\PYZsq{}}\PY{l+s+s1}{idx}\PY{l+s+s1}{\PYZsq{}}\PY{p}{)}
\PY{n}{class\PYZus{}labels} \PY{o}{=} \PY{n}{class\PYZus{}labels}\PY{o}{.}\PY{n}{sort\PYZus{}index}\PY{p}{(}\PY{p}{)}
\end{Verbatim}
\end{tcolorbox}

    We'll now construct a data frame for each image file with the labels
with the three highest probabilities as estimated by the model above.

    \begin{tcolorbox}[breakable, size=fbox, boxrule=1pt, pad at break*=1mm,colback=cellbackground, colframe=cellborder]
\prompt{In}{incolor}{68}{\boxspacing}
\begin{Verbatim}[commandchars=\\\{\}]
\PY{k}{for} \PY{n}{i}\PY{p}{,} \PY{n}{imgfile} \PY{o+ow}{in} \PY{n+nb}{enumerate}\PY{p}{(}\PY{n}{imgfiles}\PY{p}{)}\PY{p}{:}
    \PY{n}{img\PYZus{}df} \PY{o}{=} \PY{n}{class\PYZus{}labels}\PY{o}{.}\PY{n}{copy}\PY{p}{(}\PY{p}{)}
    \PY{n}{img\PYZus{}df}\PY{p}{[}\PY{l+s+s1}{\PYZsq{}}\PY{l+s+s1}{prob}\PY{l+s+s1}{\PYZsq{}}\PY{p}{]} \PY{o}{=} \PY{n}{img\PYZus{}probs}\PY{p}{[}\PY{n}{i}\PY{p}{]}
    \PY{n}{img\PYZus{}df} \PY{o}{=} \PY{n}{img\PYZus{}df}\PY{o}{.}\PY{n}{sort\PYZus{}values}\PY{p}{(}\PY{n}{by}\PY{o}{=}\PY{l+s+s1}{\PYZsq{}}\PY{l+s+s1}{prob}\PY{l+s+s1}{\PYZsq{}}\PY{p}{,} \PY{n}{ascending}\PY{o}{=}\PY{k+kc}{False}\PY{p}{)}\PY{p}{[}\PY{p}{:}\PY{l+m+mi}{3}\PY{p}{]}
    \PY{n+nb}{print}\PY{p}{(}\PY{l+s+sa}{f}\PY{l+s+s1}{\PYZsq{}}\PY{l+s+s1}{Image: }\PY{l+s+si}{\PYZob{}}\PY{n}{imgfile}\PY{l+s+si}{\PYZcb{}}\PY{l+s+s1}{\PYZsq{}}\PY{p}{)}
    \PY{n+nb}{print}\PY{p}{(}\PY{n}{img\PYZus{}df}\PY{o}{.}\PY{n}{reset\PYZus{}index}\PY{p}{(}\PY{p}{)}\PY{o}{.}\PY{n}{drop}\PY{p}{(}\PY{n}{columns}\PY{o}{=}\PY{p}{[}\PY{l+s+s1}{\PYZsq{}}\PY{l+s+s1}{idx}\PY{l+s+s1}{\PYZsq{}}\PY{p}{]}\PY{p}{)}\PY{p}{)}
\end{Verbatim}
\end{tcolorbox}

    \begin{Verbatim}[commandchars=\\\{\}]
Image: book\_images/Cape\_Weaver.jpg
      label      prob
0   jacamar  0.297499
1     macaw  0.068107
2  lorikeet  0.051105
Image: book\_images/Flamingo.jpg
            label      prob
0        flamingo  0.609515
1       spoonbill  0.013586
2  American\_egret  0.002132
Image: book\_images/Hawk\_Fountain.jpg
            label      prob
0            kite  0.184682
1           robin  0.084021
2  great\_grey\_owl  0.061274
Image: book\_images/Hawk\_cropped.jpg
            label      prob
0            kite  0.453833
1  great\_grey\_owl  0.015914
2             jay  0.012210
Image: book\_images/Lhasa\_Apso.jpg
             label      prob
0            Lhasa  0.260317
1         Shih-Tzu  0.097196
2  Tibetan\_terrier  0.032820
Image: book\_images/Sleeping\_Cat.jpg
         label      prob
0  Persian\_cat  0.163070
1        tabby  0.074143
2    tiger\_cat  0.042578
    \end{Verbatim}

    We see that the model is quite confident about \texttt{Flamingo.jpg},
but a little less so for the other images.

We end this section with our usual cleanup.

    \begin{tcolorbox}[breakable, size=fbox, boxrule=1pt, pad at break*=1mm,colback=cellbackground, colframe=cellborder]
\prompt{In}{incolor}{69}{\boxspacing}
\begin{Verbatim}[commandchars=\\\{\}]
\PY{k}{del}\PY{p}{(}\PY{n}{cifar\PYZus{}test}\PY{p}{,}
    \PY{n}{cifar\PYZus{}train}\PY{p}{,}
    \PY{n}{cifar\PYZus{}dm}\PY{p}{,}
    \PY{n}{cifar\PYZus{}module}\PY{p}{,}
    \PY{n}{cifar\PYZus{}logger}\PY{p}{,}
    \PY{n}{cifar\PYZus{}optimizer}\PY{p}{,}
    \PY{n}{cifar\PYZus{}trainer}\PY{p}{)}
\end{Verbatim}
\end{tcolorbox}

    \hypertarget{imdb-document-classification}{%
\subsection{IMDB Document
Classification}\label{imdb-document-classification}}

We now implement models for sentiment classification (Section 10.4) on
the \texttt{IMDB} dataset. As mentioned above code
block\textasciitilde8, we are using a preprocessed version of the
\texttt{IMDB} dataset found in the \texttt{keras} package. As
\texttt{keras} uses \texttt{tensorflow}, a different tensor and deep
learning library, we have converted the data to be suitable for
\texttt{torch}. The code used to convert from \texttt{keras} is
available in the module \texttt{ISLP.torch.\_make\_imdb}. It requires
some of the \texttt{keras} packages to run. These data use a dictionary
of size 10,000.

We have stored three different representations of the review data for
this lab:

\begin{itemize}
\tightlist
\item
  \texttt{load\_tensor()}, a sparse tensor version usable by
  \texttt{torch};
\item
  \texttt{load\_sparse()}, a sparse matrix version usable by
  \texttt{sklearn}, since we will compare with a lasso fit;
\item
  \texttt{load\_sequential()}, a padded version of the original sequence
  representation, limited to the last 500 words of each review.
\end{itemize}

    \begin{tcolorbox}[breakable, size=fbox, boxrule=1pt, pad at break*=1mm,colback=cellbackground, colframe=cellborder]
\prompt{In}{incolor}{70}{\boxspacing}
\begin{Verbatim}[commandchars=\\\{\}]
\PY{p}{(}\PY{n}{imdb\PYZus{}seq\PYZus{}train}\PY{p}{,}
 \PY{n}{imdb\PYZus{}seq\PYZus{}test}\PY{p}{)} \PY{o}{=} \PY{n}{load\PYZus{}sequential}\PY{p}{(}\PY{n}{root}\PY{o}{=}\PY{l+s+s1}{\PYZsq{}}\PY{l+s+s1}{data/IMDB}\PY{l+s+s1}{\PYZsq{}}\PY{p}{)}
\PY{n}{padded\PYZus{}sample} \PY{o}{=} \PY{n}{np}\PY{o}{.}\PY{n}{asarray}\PY{p}{(}\PY{n}{imdb\PYZus{}seq\PYZus{}train}\PY{o}{.}\PY{n}{tensors}\PY{p}{[}\PY{l+m+mi}{0}\PY{p}{]}\PY{p}{[}\PY{l+m+mi}{0}\PY{p}{]}\PY{p}{)}
\PY{n}{sample\PYZus{}review} \PY{o}{=} \PY{n}{padded\PYZus{}sample}\PY{p}{[}\PY{n}{padded\PYZus{}sample} \PY{o}{\PYZgt{}} \PY{l+m+mi}{0}\PY{p}{]}\PY{p}{[}\PY{p}{:}\PY{l+m+mi}{12}\PY{p}{]}
\PY{n}{sample\PYZus{}review}\PY{p}{[}\PY{p}{:}\PY{l+m+mi}{12}\PY{p}{]}
\end{Verbatim}
\end{tcolorbox}

            \begin{tcolorbox}[breakable, size=fbox, boxrule=.5pt, pad at break*=1mm, opacityfill=0]
\prompt{Out}{outcolor}{70}{\boxspacing}
\begin{Verbatim}[commandchars=\\\{\}]
array([   1,   14,   22,   16,   43,  530,  973, 1622, 1385,   65,  458,
       4468], dtype=int32)
\end{Verbatim}
\end{tcolorbox}
        
    The datasets \texttt{imdb\_seq\_train} and \texttt{imdb\_seq\_test} are
both instances of the class \texttt{TensorDataset}. The tensors used to
construct them can be found in the \texttt{tensors} attribute, with the
first tensor the features \texttt{X} and the second the outcome
\texttt{Y}. We have taken the first row of features and stored it as
\texttt{padded\_sample}. In the preprocessing used to form these data,
sequences were padded with 0s in the beginning if they were not long
enough, hence we remove this padding by restricting to entries where
\texttt{padded\_sample\ \textgreater{}\ 0}. We then provide the first 12
words of the sample review.

We can find these words in the \texttt{lookup} dictionary from the
\texttt{ISLP.torch.imdb} module.

    \begin{tcolorbox}[breakable, size=fbox, boxrule=1pt, pad at break*=1mm,colback=cellbackground, colframe=cellborder]
\prompt{In}{incolor}{71}{\boxspacing}
\begin{Verbatim}[commandchars=\\\{\}]
\PY{n}{lookup} \PY{o}{=} \PY{n}{load\PYZus{}lookup}\PY{p}{(}\PY{n}{root}\PY{o}{=}\PY{l+s+s1}{\PYZsq{}}\PY{l+s+s1}{data/IMDB}\PY{l+s+s1}{\PYZsq{}}\PY{p}{)}
\PY{l+s+s1}{\PYZsq{}}\PY{l+s+s1}{ }\PY{l+s+s1}{\PYZsq{}}\PY{o}{.}\PY{n}{join}\PY{p}{(}\PY{n}{lookup}\PY{p}{[}\PY{n}{i}\PY{p}{]} \PY{k}{for} \PY{n}{i} \PY{o+ow}{in} \PY{n}{sample\PYZus{}review}\PY{p}{)}
\end{Verbatim}
\end{tcolorbox}

            \begin{tcolorbox}[breakable, size=fbox, boxrule=.5pt, pad at break*=1mm, opacityfill=0]
\prompt{Out}{outcolor}{71}{\boxspacing}
\begin{Verbatim}[commandchars=\\\{\}]
"<START> this film was just brilliant casting location scenery story direction
everyone's"
\end{Verbatim}
\end{tcolorbox}
        
    For our first model, we have created a binary feature for each of the
10,000 possible words in the dataset, with an entry of one in the
\(i,j\) entry if word \(j\) appears in review \(i\). As most reviews are
quite short, such a feature matrix has over 98\% zeros. These data are
accessed using \texttt{load\_tensor()} from the \texttt{ISLP} library.

    \begin{tcolorbox}[breakable, size=fbox, boxrule=1pt, pad at break*=1mm,colback=cellbackground, colframe=cellborder]
\prompt{In}{incolor}{72}{\boxspacing}
\begin{Verbatim}[commandchars=\\\{\}]
\PY{n}{max\PYZus{}num\PYZus{}workers}\PY{o}{=}\PY{l+m+mi}{10}
\PY{p}{(}\PY{n}{imdb\PYZus{}train}\PY{p}{,}
 \PY{n}{imdb\PYZus{}test}\PY{p}{)} \PY{o}{=} \PY{n}{load\PYZus{}tensor}\PY{p}{(}\PY{n}{root}\PY{o}{=}\PY{l+s+s1}{\PYZsq{}}\PY{l+s+s1}{data/IMDB}\PY{l+s+s1}{\PYZsq{}}\PY{p}{)}
\PY{n}{imdb\PYZus{}dm} \PY{o}{=} \PY{n}{SimpleDataModule}\PY{p}{(}\PY{n}{imdb\PYZus{}train}\PY{p}{,}
                           \PY{n}{imdb\PYZus{}test}\PY{p}{,}
                           \PY{n}{validation}\PY{o}{=}\PY{l+m+mi}{2000}\PY{p}{,}
                           \PY{n}{num\PYZus{}workers}\PY{o}{=}\PY{n+nb}{min}\PY{p}{(}\PY{l+m+mi}{6}\PY{p}{,} \PY{n}{max\PYZus{}num\PYZus{}workers}\PY{p}{)}\PY{p}{,}
                           \PY{n}{batch\PYZus{}size}\PY{o}{=}\PY{l+m+mi}{512}\PY{p}{)}
\end{Verbatim}
\end{tcolorbox}

    We'll use a two-layer model for our first model.

    \begin{tcolorbox}[breakable, size=fbox, boxrule=1pt, pad at break*=1mm,colback=cellbackground, colframe=cellborder]
\prompt{In}{incolor}{73}{\boxspacing}
\begin{Verbatim}[commandchars=\\\{\}]
\PY{k}{class}\PY{+w}{ }\PY{n+nc}{IMDBModel}\PY{p}{(}\PY{n}{nn}\PY{o}{.}\PY{n}{Module}\PY{p}{)}\PY{p}{:}

    \PY{k}{def}\PY{+w}{ }\PY{n+nf+fm}{\PYZus{}\PYZus{}init\PYZus{}\PYZus{}}\PY{p}{(}\PY{n+nb+bp}{self}\PY{p}{,} \PY{n}{input\PYZus{}size}\PY{p}{)}\PY{p}{:}
        \PY{n+nb}{super}\PY{p}{(}\PY{n}{IMDBModel}\PY{p}{,} \PY{n+nb+bp}{self}\PY{p}{)}\PY{o}{.}\PY{n+nf+fm}{\PYZus{}\PYZus{}init\PYZus{}\PYZus{}}\PY{p}{(}\PY{p}{)}
        \PY{n+nb+bp}{self}\PY{o}{.}\PY{n}{dense1} \PY{o}{=} \PY{n}{nn}\PY{o}{.}\PY{n}{Linear}\PY{p}{(}\PY{n}{input\PYZus{}size}\PY{p}{,} \PY{l+m+mi}{16}\PY{p}{)}
        \PY{n+nb+bp}{self}\PY{o}{.}\PY{n}{activation} \PY{o}{=} \PY{n}{nn}\PY{o}{.}\PY{n}{ReLU}\PY{p}{(}\PY{p}{)}
        \PY{n+nb+bp}{self}\PY{o}{.}\PY{n}{dense2} \PY{o}{=} \PY{n}{nn}\PY{o}{.}\PY{n}{Linear}\PY{p}{(}\PY{l+m+mi}{16}\PY{p}{,} \PY{l+m+mi}{16}\PY{p}{)}
        \PY{n+nb+bp}{self}\PY{o}{.}\PY{n}{output} \PY{o}{=} \PY{n}{nn}\PY{o}{.}\PY{n}{Linear}\PY{p}{(}\PY{l+m+mi}{16}\PY{p}{,} \PY{l+m+mi}{1}\PY{p}{)}

    \PY{k}{def}\PY{+w}{ }\PY{n+nf}{forward}\PY{p}{(}\PY{n+nb+bp}{self}\PY{p}{,} \PY{n}{x}\PY{p}{)}\PY{p}{:}
        \PY{n}{val} \PY{o}{=} \PY{n}{x}
        \PY{k}{for} \PY{n}{\PYZus{}map} \PY{o+ow}{in} \PY{p}{[}\PY{n+nb+bp}{self}\PY{o}{.}\PY{n}{dense1}\PY{p}{,}
                     \PY{n+nb+bp}{self}\PY{o}{.}\PY{n}{activation}\PY{p}{,}
                     \PY{n+nb+bp}{self}\PY{o}{.}\PY{n}{dense2}\PY{p}{,}
                     \PY{n+nb+bp}{self}\PY{o}{.}\PY{n}{activation}\PY{p}{,}
                     \PY{n+nb+bp}{self}\PY{o}{.}\PY{n}{output}\PY{p}{]}\PY{p}{:}
            \PY{n}{val} \PY{o}{=} \PY{n}{\PYZus{}map}\PY{p}{(}\PY{n}{val}\PY{p}{)}
        \PY{k}{return} \PY{n}{torch}\PY{o}{.}\PY{n}{flatten}\PY{p}{(}\PY{n}{val}\PY{p}{)}
\end{Verbatim}
\end{tcolorbox}

    We now instantiate our model and look at a summary.

    \begin{tcolorbox}[breakable, size=fbox, boxrule=1pt, pad at break*=1mm,colback=cellbackground, colframe=cellborder]
\prompt{In}{incolor}{74}{\boxspacing}
\begin{Verbatim}[commandchars=\\\{\}]
\PY{n}{imdb\PYZus{}model} \PY{o}{=} \PY{n}{IMDBModel}\PY{p}{(}\PY{n}{imdb\PYZus{}test}\PY{o}{.}\PY{n}{tensors}\PY{p}{[}\PY{l+m+mi}{0}\PY{p}{]}\PY{o}{.}\PY{n}{size}\PY{p}{(}\PY{p}{)}\PY{p}{[}\PY{l+m+mi}{1}\PY{p}{]}\PY{p}{)}
\PY{n}{summary}\PY{p}{(}\PY{n}{imdb\PYZus{}model}\PY{p}{,}
        \PY{n}{input\PYZus{}size}\PY{o}{=}\PY{n}{imdb\PYZus{}test}\PY{o}{.}\PY{n}{tensors}\PY{p}{[}\PY{l+m+mi}{0}\PY{p}{]}\PY{o}{.}\PY{n}{size}\PY{p}{(}\PY{p}{)}\PY{p}{,}
        \PY{n}{col\PYZus{}names}\PY{o}{=}\PY{p}{[}\PY{l+s+s1}{\PYZsq{}}\PY{l+s+s1}{input\PYZus{}size}\PY{l+s+s1}{\PYZsq{}}\PY{p}{,}
                   \PY{l+s+s1}{\PYZsq{}}\PY{l+s+s1}{output\PYZus{}size}\PY{l+s+s1}{\PYZsq{}}\PY{p}{,}
                   \PY{l+s+s1}{\PYZsq{}}\PY{l+s+s1}{num\PYZus{}params}\PY{l+s+s1}{\PYZsq{}}\PY{p}{]}\PY{p}{)}
\end{Verbatim}
\end{tcolorbox}

            \begin{tcolorbox}[breakable, size=fbox, boxrule=.5pt, pad at break*=1mm, opacityfill=0]
\prompt{Out}{outcolor}{74}{\boxspacing}
\begin{Verbatim}[commandchars=\\\{\}]
================================================================================
===================================
Layer (type:depth-idx)                   Input Shape               Output Shape
Param \#
================================================================================
===================================
IMDBModel                                [25000, 10003]            [25000]
--
├─Linear: 1-1                            [25000, 10003]            [25000, 16]
160,064
├─ReLU: 1-2                              [25000, 16]               [25000, 16]
--
├─Linear: 1-3                            [25000, 16]               [25000, 16]
272
├─ReLU: 1-4                              [25000, 16]               [25000, 16]
--
├─Linear: 1-5                            [25000, 16]               [25000, 1]
17
================================================================================
===================================
Total params: 160,353
Trainable params: 160,353
Non-trainable params: 0
Total mult-adds (Units.GIGABYTES): 4.01
================================================================================
===================================
Input size (MB): 1000.30
Forward/backward pass size (MB): 6.60
Params size (MB): 0.64
Estimated Total Size (MB): 1007.54
================================================================================
===================================
\end{Verbatim}
\end{tcolorbox}
        
    We'll again use a smaller learning rate for these data, hence we pass an
\texttt{optimizer} to the \texttt{SimpleModule}. Since the reviews are
classified into positive or negative sentiment, we use
\texttt{SimpleModule.binary\_classification()}. \{Our use of
\texttt{binary\_classification()} instead of \texttt{classification()}
is due to some subtlety in how \texttt{torchmetrics.Accuracy()} works,
as well as the data type of the targets.\}

    \begin{tcolorbox}[breakable, size=fbox, boxrule=1pt, pad at break*=1mm,colback=cellbackground, colframe=cellborder]
\prompt{In}{incolor}{75}{\boxspacing}
\begin{Verbatim}[commandchars=\\\{\}]
\PY{n}{imdb\PYZus{}optimizer} \PY{o}{=} \PY{n}{RMSprop}\PY{p}{(}\PY{n}{imdb\PYZus{}model}\PY{o}{.}\PY{n}{parameters}\PY{p}{(}\PY{p}{)}\PY{p}{,} \PY{n}{lr}\PY{o}{=}\PY{l+m+mf}{0.001}\PY{p}{)}
\PY{n}{imdb\PYZus{}module} \PY{o}{=} \PY{n}{SimpleModule}\PY{o}{.}\PY{n}{binary\PYZus{}classification}\PY{p}{(}
                         \PY{n}{imdb\PYZus{}model}\PY{p}{,}
                         \PY{n}{optimizer}\PY{o}{=}\PY{n}{imdb\PYZus{}optimizer}\PY{p}{)}
\end{Verbatim}
\end{tcolorbox}

    Having loaded the datasets into a data module and created a
\texttt{SimpleModule}, the remaining steps are familiar.

    \begin{tcolorbox}[breakable, size=fbox, boxrule=1pt, pad at break*=1mm,colback=cellbackground, colframe=cellborder]
\prompt{In}{incolor}{76}{\boxspacing}
\begin{Verbatim}[commandchars=\\\{\}]
\PY{n}{imdb\PYZus{}logger} \PY{o}{=} \PY{n}{CSVLogger}\PY{p}{(}\PY{l+s+s1}{\PYZsq{}}\PY{l+s+s1}{logs}\PY{l+s+s1}{\PYZsq{}}\PY{p}{,} \PY{n}{name}\PY{o}{=}\PY{l+s+s1}{\PYZsq{}}\PY{l+s+s1}{IMDB}\PY{l+s+s1}{\PYZsq{}}\PY{p}{)}
\PY{n}{imdb\PYZus{}trainer} \PY{o}{=} \PY{n}{Trainer}\PY{p}{(}\PY{n}{deterministic}\PY{o}{=}\PY{k+kc}{True}\PY{p}{,}
                       \PY{n}{max\PYZus{}epochs}\PY{o}{=}\PY{l+m+mi}{30}\PY{p}{,}
                       \PY{n}{logger}\PY{o}{=}\PY{n}{imdb\PYZus{}logger}\PY{p}{,}
                       \PY{n}{enable\PYZus{}progress\PYZus{}bar}\PY{o}{=}\PY{k+kc}{False}\PY{p}{,}
                       \PY{n}{callbacks}\PY{o}{=}\PY{p}{[}\PY{n}{ErrorTracker}\PY{p}{(}\PY{p}{)}\PY{p}{]}\PY{p}{)}
\PY{n}{imdb\PYZus{}trainer}\PY{o}{.}\PY{n}{fit}\PY{p}{(}\PY{n}{imdb\PYZus{}module}\PY{p}{,}
                 \PY{n}{datamodule}\PY{o}{=}\PY{n}{imdb\PYZus{}dm}\PY{p}{)}
\end{Verbatim}
\end{tcolorbox}

    \begin{Verbatim}[commandchars=\\\{\}]
GPU available: True (mps), used: True
TPU available: False, using: 0 TPU cores
IPU available: False, using: 0 IPUs
HPU available: False, using: 0 HPUs

  | Name  | Type              | Params
--------------------------------------------
0 | model | IMDBModel         | 160 K
1 | loss  | BCEWithLogitsLoss | 0
--------------------------------------------
160 K     Trainable params
0         Non-trainable params
160 K     Total params
0.641     Total estimated model params size (MB)
/Users/jonathantaylor/anaconda3/envs/ISLP\_312/lib/python3.12/site-
packages/pytorch\_lightning/loops/fit\_loop.py:298: The number of training batches
(45) is smaller than the logging interval Trainer(log\_every\_n\_steps=50). Set a
lower value for log\_every\_n\_steps if you want to see logs for the training
epoch.
`Trainer.fit` stopped: `max\_epochs=30` reached.
    \end{Verbatim}

    Evaluating the test error yields roughly 86\% accuracy.

    \begin{tcolorbox}[breakable, size=fbox, boxrule=1pt, pad at break*=1mm,colback=cellbackground, colframe=cellborder]
\prompt{In}{incolor}{77}{\boxspacing}
\begin{Verbatim}[commandchars=\\\{\}]
\PY{n}{test\PYZus{}results} \PY{o}{=} \PY{n}{imdb\PYZus{}trainer}\PY{o}{.}\PY{n}{test}\PY{p}{(}\PY{n}{imdb\PYZus{}module}\PY{p}{,} \PY{n}{datamodule}\PY{o}{=}\PY{n}{imdb\PYZus{}dm}\PY{p}{)}
\PY{n}{test\PYZus{}results}
\end{Verbatim}
\end{tcolorbox}

    \begin{Verbatim}[commandchars=\\\{\}]
────────────────────────────────────────────────────────────────────────────────
────────────────────────────────────────
       Test metric             DataLoader 0
────────────────────────────────────────────────────────────────────────────────
────────────────────────────────────────
      test\_accuracy         0.8450000286102295
        test\_loss           1.2167928218841553
────────────────────────────────────────────────────────────────────────────────
────────────────────────────────────────
    \end{Verbatim}

            \begin{tcolorbox}[breakable, size=fbox, boxrule=.5pt, pad at break*=1mm, opacityfill=0]
\prompt{Out}{outcolor}{77}{\boxspacing}
\begin{Verbatim}[commandchars=\\\{\}]
[\{'test\_loss': 1.2167928218841553, 'test\_accuracy': 0.8450000286102295\}]
\end{Verbatim}
\end{tcolorbox}
        
    \hypertarget{comparison-to-lasso}{%
\subsubsection{Comparison to Lasso}\label{comparison-to-lasso}}

We now fit a lasso logistic regression model using
\texttt{LogisticRegression()} from \texttt{sklearn}. Since
\texttt{sklearn} does not recognize the sparse tensors of
\texttt{torch}, we use a sparse matrix that is recognized by
\texttt{sklearn.}

    \begin{tcolorbox}[breakable, size=fbox, boxrule=1pt, pad at break*=1mm,colback=cellbackground, colframe=cellborder]
\prompt{In}{incolor}{78}{\boxspacing}
\begin{Verbatim}[commandchars=\\\{\}]
\PY{p}{(}\PY{p}{(}\PY{n}{X\PYZus{}train}\PY{p}{,} \PY{n}{Y\PYZus{}train}\PY{p}{)}\PY{p}{,}
 \PY{p}{(}\PY{n}{X\PYZus{}valid}\PY{p}{,} \PY{n}{Y\PYZus{}valid}\PY{p}{)}\PY{p}{,}
 \PY{p}{(}\PY{n}{X\PYZus{}test}\PY{p}{,} \PY{n}{Y\PYZus{}test}\PY{p}{)}\PY{p}{)} \PY{o}{=} \PY{n}{load\PYZus{}sparse}\PY{p}{(}\PY{n}{validation}\PY{o}{=}\PY{l+m+mi}{2000}\PY{p}{,}
                                 \PY{n}{random\PYZus{}state}\PY{o}{=}\PY{l+m+mi}{0}\PY{p}{,}
                                 \PY{n}{root}\PY{o}{=}\PY{l+s+s1}{\PYZsq{}}\PY{l+s+s1}{data/IMDB}\PY{l+s+s1}{\PYZsq{}}\PY{p}{)}
\end{Verbatim}
\end{tcolorbox}

    Similar to what we did in Section 10.9.1, we construct a series of 50
values for the lasso reguralization parameter \(\lambda\).

    \begin{tcolorbox}[breakable, size=fbox, boxrule=1pt, pad at break*=1mm,colback=cellbackground, colframe=cellborder]
\prompt{In}{incolor}{79}{\boxspacing}
\begin{Verbatim}[commandchars=\\\{\}]
\PY{n}{lam\PYZus{}max} \PY{o}{=} \PY{n}{np}\PY{o}{.}\PY{n}{abs}\PY{p}{(}\PY{n}{X\PYZus{}train}\PY{o}{.}\PY{n}{T} \PY{o}{*} \PY{p}{(}\PY{n}{Y\PYZus{}train} \PY{o}{\PYZhy{}} \PY{n}{Y\PYZus{}train}\PY{o}{.}\PY{n}{mean}\PY{p}{(}\PY{p}{)}\PY{p}{)}\PY{p}{)}\PY{o}{.}\PY{n}{max}\PY{p}{(}\PY{p}{)}
\PY{n}{lam\PYZus{}val} \PY{o}{=} \PY{n}{lam\PYZus{}max} \PY{o}{*} \PY{n}{np}\PY{o}{.}\PY{n}{exp}\PY{p}{(}\PY{n}{np}\PY{o}{.}\PY{n}{linspace}\PY{p}{(}\PY{n}{np}\PY{o}{.}\PY{n}{log}\PY{p}{(}\PY{l+m+mi}{1}\PY{p}{)}\PY{p}{,}
                                       \PY{n}{np}\PY{o}{.}\PY{n}{log}\PY{p}{(}\PY{l+m+mf}{1e\PYZhy{}4}\PY{p}{)}\PY{p}{,} \PY{l+m+mi}{50}\PY{p}{)}\PY{p}{)}
\end{Verbatim}
\end{tcolorbox}

    With \texttt{LogisticRegression()} the regularization parameter \(C\) is
specified as the inverse of \(\lambda\). There are several solvers for
logistic regression; here we use \texttt{liblinear} which works well
with the sparse input format.

    \begin{tcolorbox}[breakable, size=fbox, boxrule=1pt, pad at break*=1mm,colback=cellbackground, colframe=cellborder]
\prompt{In}{incolor}{80}{\boxspacing}
\begin{Verbatim}[commandchars=\\\{\}]
\PY{n}{logit} \PY{o}{=} \PY{n}{LogisticRegression}\PY{p}{(}\PY{n}{penalty}\PY{o}{=}\PY{l+s+s1}{\PYZsq{}}\PY{l+s+s1}{l1}\PY{l+s+s1}{\PYZsq{}}\PY{p}{,} 
                           \PY{n}{C}\PY{o}{=}\PY{l+m+mi}{1}\PY{o}{/}\PY{n}{lam\PYZus{}max}\PY{p}{,}
                           \PY{n}{solver}\PY{o}{=}\PY{l+s+s1}{\PYZsq{}}\PY{l+s+s1}{liblinear}\PY{l+s+s1}{\PYZsq{}}\PY{p}{,}
                           \PY{n}{warm\PYZus{}start}\PY{o}{=}\PY{k+kc}{True}\PY{p}{,}
                           \PY{n}{fit\PYZus{}intercept}\PY{o}{=}\PY{k+kc}{True}\PY{p}{)}
\end{Verbatim}
\end{tcolorbox}

    The path of 50 values takes approximately 40 seconds to run.

    \begin{tcolorbox}[breakable, size=fbox, boxrule=1pt, pad at break*=1mm,colback=cellbackground, colframe=cellborder]
\prompt{In}{incolor}{81}{\boxspacing}
\begin{Verbatim}[commandchars=\\\{\}]
\PY{n}{coefs} \PY{o}{=} \PY{p}{[}\PY{p}{]}
\PY{n}{intercepts} \PY{o}{=} \PY{p}{[}\PY{p}{]}

\PY{k}{for} \PY{n}{l} \PY{o+ow}{in} \PY{n}{lam\PYZus{}val}\PY{p}{:}
    \PY{n}{logit}\PY{o}{.}\PY{n}{C} \PY{o}{=} \PY{l+m+mi}{1}\PY{o}{/}\PY{n}{l}
    \PY{n}{logit}\PY{o}{.}\PY{n}{fit}\PY{p}{(}\PY{n}{X\PYZus{}train}\PY{p}{,} \PY{n}{Y\PYZus{}train}\PY{p}{)}
    \PY{n}{coefs}\PY{o}{.}\PY{n}{append}\PY{p}{(}\PY{n}{logit}\PY{o}{.}\PY{n}{coef\PYZus{}}\PY{o}{.}\PY{n}{copy}\PY{p}{(}\PY{p}{)}\PY{p}{)}
    \PY{n}{intercepts}\PY{o}{.}\PY{n}{append}\PY{p}{(}\PY{n}{logit}\PY{o}{.}\PY{n}{intercept\PYZus{}}\PY{p}{)}
\end{Verbatim}
\end{tcolorbox}

    The coefficient and intercepts have an extraneous dimension which can be
removed by the \texttt{np.squeeze()} function.

    \begin{tcolorbox}[breakable, size=fbox, boxrule=1pt, pad at break*=1mm,colback=cellbackground, colframe=cellborder]
\prompt{In}{incolor}{82}{\boxspacing}
\begin{Verbatim}[commandchars=\\\{\}]
\PY{n}{coefs} \PY{o}{=} \PY{n}{np}\PY{o}{.}\PY{n}{squeeze}\PY{p}{(}\PY{n}{coefs}\PY{p}{)}
\PY{n}{intercepts} \PY{o}{=} \PY{n}{np}\PY{o}{.}\PY{n}{squeeze}\PY{p}{(}\PY{n}{intercepts}\PY{p}{)}
\end{Verbatim}
\end{tcolorbox}

    We'll now make a plot to compare our neural network results with the
lasso.

    \begin{tcolorbox}[breakable, size=fbox, boxrule=1pt, pad at break*=1mm,colback=cellbackground, colframe=cellborder]
\prompt{In}{incolor}{83}{\boxspacing}
\begin{Verbatim}[commandchars=\\\{\}]
\PY{o}{\PYZpc{}\PYZpc{}capture}
\PY{n}{fig}\PY{p}{,} \PY{n}{axes} \PY{o}{=} \PY{n}{subplots}\PY{p}{(}\PY{l+m+mi}{1}\PY{p}{,} \PY{l+m+mi}{2}\PY{p}{,} \PY{n}{figsize}\PY{o}{=}\PY{p}{(}\PY{l+m+mi}{16}\PY{p}{,} \PY{l+m+mi}{8}\PY{p}{)}\PY{p}{,} \PY{n}{sharey}\PY{o}{=}\PY{k+kc}{True}\PY{p}{)}
\PY{k}{for} \PY{p}{(}\PY{p}{(}\PY{n}{X\PYZus{}}\PY{p}{,} \PY{n}{Y\PYZus{}}\PY{p}{)}\PY{p}{,}
     \PY{n}{data\PYZus{}}\PY{p}{,}
     \PY{n}{color}\PY{p}{)} \PY{o+ow}{in} \PY{n+nb}{zip}\PY{p}{(}\PY{p}{[}\PY{p}{(}\PY{n}{X\PYZus{}train}\PY{p}{,} \PY{n}{Y\PYZus{}train}\PY{p}{)}\PY{p}{,}
                    \PY{p}{(}\PY{n}{X\PYZus{}valid}\PY{p}{,} \PY{n}{Y\PYZus{}valid}\PY{p}{)}\PY{p}{,}
                    \PY{p}{(}\PY{n}{X\PYZus{}test}\PY{p}{,} \PY{n}{Y\PYZus{}test}\PY{p}{)}\PY{p}{]}\PY{p}{,}
                    \PY{p}{[}\PY{l+s+s1}{\PYZsq{}}\PY{l+s+s1}{Training}\PY{l+s+s1}{\PYZsq{}}\PY{p}{,} \PY{l+s+s1}{\PYZsq{}}\PY{l+s+s1}{Validation}\PY{l+s+s1}{\PYZsq{}}\PY{p}{,} \PY{l+s+s1}{\PYZsq{}}\PY{l+s+s1}{Test}\PY{l+s+s1}{\PYZsq{}}\PY{p}{]}\PY{p}{,}
                    \PY{p}{[}\PY{l+s+s1}{\PYZsq{}}\PY{l+s+s1}{black}\PY{l+s+s1}{\PYZsq{}}\PY{p}{,} \PY{l+s+s1}{\PYZsq{}}\PY{l+s+s1}{red}\PY{l+s+s1}{\PYZsq{}}\PY{p}{,} \PY{l+s+s1}{\PYZsq{}}\PY{l+s+s1}{blue}\PY{l+s+s1}{\PYZsq{}}\PY{p}{]}\PY{p}{)}\PY{p}{:}
    \PY{n}{linpred\PYZus{}} \PY{o}{=} \PY{n}{X\PYZus{}} \PY{o}{*} \PY{n}{coefs}\PY{o}{.}\PY{n}{T} \PY{o}{+} \PY{n}{intercepts}\PY{p}{[}\PY{k+kc}{None}\PY{p}{,}\PY{p}{:}\PY{p}{]}
    \PY{n}{label\PYZus{}} \PY{o}{=} \PY{n}{np}\PY{o}{.}\PY{n}{array}\PY{p}{(}\PY{n}{linpred\PYZus{}} \PY{o}{\PYZgt{}} \PY{l+m+mi}{0}\PY{p}{)}
    \PY{n}{accuracy\PYZus{}} \PY{o}{=} \PY{n}{np}\PY{o}{.}\PY{n}{array}\PY{p}{(}\PY{p}{[}\PY{n}{np}\PY{o}{.}\PY{n}{mean}\PY{p}{(}\PY{n}{Y\PYZus{}} \PY{o}{==} \PY{n}{l}\PY{p}{)} \PY{k}{for} \PY{n}{l} \PY{o+ow}{in} \PY{n}{label\PYZus{}}\PY{o}{.}\PY{n}{T}\PY{p}{]}\PY{p}{)}
    \PY{n}{axes}\PY{p}{[}\PY{l+m+mi}{0}\PY{p}{]}\PY{o}{.}\PY{n}{plot}\PY{p}{(}\PY{o}{\PYZhy{}}\PY{n}{np}\PY{o}{.}\PY{n}{log}\PY{p}{(}\PY{n}{lam\PYZus{}val} \PY{o}{/} \PY{n}{X\PYZus{}train}\PY{o}{.}\PY{n}{shape}\PY{p}{[}\PY{l+m+mi}{0}\PY{p}{]}\PY{p}{)}\PY{p}{,}
                 \PY{n}{accuracy\PYZus{}}\PY{p}{,}
                 \PY{l+s+s1}{\PYZsq{}}\PY{l+s+s1}{.\PYZhy{}\PYZhy{}}\PY{l+s+s1}{\PYZsq{}}\PY{p}{,}
                 \PY{n}{color}\PY{o}{=}\PY{n}{color}\PY{p}{,}
                 \PY{n}{markersize}\PY{o}{=}\PY{l+m+mi}{13}\PY{p}{,}
                 \PY{n}{linewidth}\PY{o}{=}\PY{l+m+mi}{2}\PY{p}{,}
                 \PY{n}{label}\PY{o}{=}\PY{n}{data\PYZus{}}\PY{p}{)}
\PY{n}{axes}\PY{p}{[}\PY{l+m+mi}{0}\PY{p}{]}\PY{o}{.}\PY{n}{legend}\PY{p}{(}\PY{p}{)}
\PY{n}{axes}\PY{p}{[}\PY{l+m+mi}{0}\PY{p}{]}\PY{o}{.}\PY{n}{set\PYZus{}xlabel}\PY{p}{(}\PY{l+s+sa}{r}\PY{l+s+s1}{\PYZsq{}}\PY{l+s+s1}{\PYZdl{}\PYZhy{}}\PY{l+s+s1}{\PYZbs{}}\PY{l+s+s1}{log(}\PY{l+s+s1}{\PYZbs{}}\PY{l+s+s1}{lambda)\PYZdl{}}\PY{l+s+s1}{\PYZsq{}}\PY{p}{,} \PY{n}{fontsize}\PY{o}{=}\PY{l+m+mi}{20}\PY{p}{)}
\PY{n}{axes}\PY{p}{[}\PY{l+m+mi}{0}\PY{p}{]}\PY{o}{.}\PY{n}{set\PYZus{}ylabel}\PY{p}{(}\PY{l+s+s1}{\PYZsq{}}\PY{l+s+s1}{Accuracy}\PY{l+s+s1}{\PYZsq{}}\PY{p}{,} \PY{n}{fontsize}\PY{o}{=}\PY{l+m+mi}{20}\PY{p}{)}
\end{Verbatim}
\end{tcolorbox}

    Notice the use of \texttt{\%\%capture}, which suppresses the displaying
of the partially completed figure. This is useful when making a complex
figure, since the steps can be spread across two or more cells. We now
add a plot of the lasso accuracy, and display the composed figure by
simply entering its name at the end of the cell.

    \begin{tcolorbox}[breakable, size=fbox, boxrule=1pt, pad at break*=1mm,colback=cellbackground, colframe=cellborder]
\prompt{In}{incolor}{84}{\boxspacing}
\begin{Verbatim}[commandchars=\\\{\}]
\PY{n}{imdb\PYZus{}results} \PY{o}{=} \PY{n}{pd}\PY{o}{.}\PY{n}{read\PYZus{}csv}\PY{p}{(}\PY{n}{imdb\PYZus{}logger}\PY{o}{.}\PY{n}{experiment}\PY{o}{.}\PY{n}{metrics\PYZus{}file\PYZus{}path}\PY{p}{)}
\PY{n}{summary\PYZus{}plot}\PY{p}{(}\PY{n}{imdb\PYZus{}results}\PY{p}{,}
             \PY{n}{axes}\PY{p}{[}\PY{l+m+mi}{1}\PY{p}{]}\PY{p}{,}
             \PY{n}{col}\PY{o}{=}\PY{l+s+s1}{\PYZsq{}}\PY{l+s+s1}{accuracy}\PY{l+s+s1}{\PYZsq{}}\PY{p}{,}
             \PY{n}{ylabel}\PY{o}{=}\PY{l+s+s1}{\PYZsq{}}\PY{l+s+s1}{Accuracy}\PY{l+s+s1}{\PYZsq{}}\PY{p}{)}
\PY{n}{axes}\PY{p}{[}\PY{l+m+mi}{1}\PY{p}{]}\PY{o}{.}\PY{n}{set\PYZus{}xticks}\PY{p}{(}\PY{n}{np}\PY{o}{.}\PY{n}{linspace}\PY{p}{(}\PY{l+m+mi}{0}\PY{p}{,} \PY{l+m+mi}{30}\PY{p}{,} \PY{l+m+mi}{7}\PY{p}{)}\PY{o}{.}\PY{n}{astype}\PY{p}{(}\PY{n+nb}{int}\PY{p}{)}\PY{p}{)}
\PY{n}{axes}\PY{p}{[}\PY{l+m+mi}{1}\PY{p}{]}\PY{o}{.}\PY{n}{set\PYZus{}ylabel}\PY{p}{(}\PY{l+s+s1}{\PYZsq{}}\PY{l+s+s1}{Accuracy}\PY{l+s+s1}{\PYZsq{}}\PY{p}{,} \PY{n}{fontsize}\PY{o}{=}\PY{l+m+mi}{20}\PY{p}{)}
\PY{n}{axes}\PY{p}{[}\PY{l+m+mi}{1}\PY{p}{]}\PY{o}{.}\PY{n}{set\PYZus{}xlabel}\PY{p}{(}\PY{l+s+s1}{\PYZsq{}}\PY{l+s+s1}{Epoch}\PY{l+s+s1}{\PYZsq{}}\PY{p}{,} \PY{n}{fontsize}\PY{o}{=}\PY{l+m+mi}{20}\PY{p}{)}
\PY{n}{axes}\PY{p}{[}\PY{l+m+mi}{1}\PY{p}{]}\PY{o}{.}\PY{n}{set\PYZus{}ylim}\PY{p}{(}\PY{p}{[}\PY{l+m+mf}{0.5}\PY{p}{,} \PY{l+m+mi}{1}\PY{p}{]}\PY{p}{)}\PY{p}{;}
\PY{n}{axes}\PY{p}{[}\PY{l+m+mi}{1}\PY{p}{]}\PY{o}{.}\PY{n}{axhline}\PY{p}{(}\PY{n}{test\PYZus{}results}\PY{p}{[}\PY{l+m+mi}{0}\PY{p}{]}\PY{p}{[}\PY{l+s+s1}{\PYZsq{}}\PY{l+s+s1}{test\PYZus{}accuracy}\PY{l+s+s1}{\PYZsq{}}\PY{p}{]}\PY{p}{,}
                \PY{n}{color}\PY{o}{=}\PY{l+s+s1}{\PYZsq{}}\PY{l+s+s1}{blue}\PY{l+s+s1}{\PYZsq{}}\PY{p}{,}
                \PY{n}{linestyle}\PY{o}{=}\PY{l+s+s1}{\PYZsq{}}\PY{l+s+s1}{\PYZhy{}\PYZhy{}}\PY{l+s+s1}{\PYZsq{}}\PY{p}{,}
                \PY{n}{linewidth}\PY{o}{=}\PY{l+m+mi}{3}\PY{p}{)}
\PY{n}{fig}
\end{Verbatim}
\end{tcolorbox}
 
            
\prompt{Out}{outcolor}{84}{}
    
    \begin{center}
    \adjustimage{max size={0.9\linewidth}{0.9\paperheight}}{artifacts/latex/Ch10-deeplearning-lab_files/artifacts/latex/Ch10-deeplearning-lab_170_0.png}
    \end{center}
    { \hspace*{\fill} \\}
    

    From the graphs we see that the accuracy of the lasso logistic
regression peaks at about \(0.88\), as it does for the neural network.

Once again, we end with a cleanup.

    \begin{tcolorbox}[breakable, size=fbox, boxrule=1pt, pad at break*=1mm,colback=cellbackground, colframe=cellborder]
\prompt{In}{incolor}{85}{\boxspacing}
\begin{Verbatim}[commandchars=\\\{\}]
\PY{k}{del}\PY{p}{(}\PY{n}{imdb\PYZus{}model}\PY{p}{,}
    \PY{n}{imdb\PYZus{}trainer}\PY{p}{,}
    \PY{n}{imdb\PYZus{}logger}\PY{p}{,}
    \PY{n}{imdb\PYZus{}dm}\PY{p}{,}
    \PY{n}{imdb\PYZus{}train}\PY{p}{,}
    \PY{n}{imdb\PYZus{}test}\PY{p}{)}
\end{Verbatim}
\end{tcolorbox}

    \hypertarget{recurrent-neural-networks}{%
\subsection{Recurrent Neural Networks}\label{recurrent-neural-networks}}

In this lab we fit the models illustrated in Section 10.5.

    \hypertarget{sequential-models-for-document-classification}{%
\subsubsection{Sequential Models for Document
Classification}\label{sequential-models-for-document-classification}}

Here we fit a simple LSTM RNN for sentiment prediction to the
\texttt{IMDb} movie-review data, as discussed in Section 10.5.1. For an
RNN we use the sequence of words in a document, taking their order into
account. We loaded the preprocessed data at the beginning of Section
10.9.5. A script that details the preprocessing can be found in the
\texttt{ISLP} library. Notably, since more than 90\% of the documents
had fewer than 500 words, we set the document length to 500. For longer
documents, we used the last 500 words, and for shorter documents, we
padded the front with blanks.

    \begin{tcolorbox}[breakable, size=fbox, boxrule=1pt, pad at break*=1mm,colback=cellbackground, colframe=cellborder]
\prompt{In}{incolor}{86}{\boxspacing}
\begin{Verbatim}[commandchars=\\\{\}]
\PY{n}{imdb\PYZus{}seq\PYZus{}dm} \PY{o}{=} \PY{n}{SimpleDataModule}\PY{p}{(}\PY{n}{imdb\PYZus{}seq\PYZus{}train}\PY{p}{,}
                               \PY{n}{imdb\PYZus{}seq\PYZus{}test}\PY{p}{,}
                               \PY{n}{validation}\PY{o}{=}\PY{l+m+mi}{2000}\PY{p}{,}
                               \PY{n}{batch\PYZus{}size}\PY{o}{=}\PY{l+m+mi}{300}\PY{p}{,}
                               \PY{n}{num\PYZus{}workers}\PY{o}{=}\PY{n+nb}{min}\PY{p}{(}\PY{l+m+mi}{6}\PY{p}{,} \PY{n}{max\PYZus{}num\PYZus{}workers}\PY{p}{)}
                               \PY{p}{)}
\end{Verbatim}
\end{tcolorbox}

    The first layer of the RNN is an embedding layer of size 32, which will
be learned during training. This layer one-hot encodes each document as
a matrix of dimension \(500 \times 10,003\), and then maps these
\(10,003\) dimensions down to \(32\). \{The extra 3 dimensions
correspond to commonly occurring non-word entries in the reviews.\}
Since each word is represented by an integer, this is effectively
achieved by the creation of an embedding matrix of size
\(10,003\times 32\); each of the 500 integers in the document are then
mapped to the appropriate 32 real numbers by indexing the appropriate
rows of this matrix.

    The second layer is an LSTM with 32 units, and the output layer is a
single logit for the binary classification task. In the last line of the
\texttt{forward()} method below, we take the last 32-dimensional output
of the LSTM and map it to our response.

    \begin{tcolorbox}[breakable, size=fbox, boxrule=1pt, pad at break*=1mm,colback=cellbackground, colframe=cellborder]
\prompt{In}{incolor}{87}{\boxspacing}
\begin{Verbatim}[commandchars=\\\{\}]
\PY{k}{class}\PY{+w}{ }\PY{n+nc}{LSTMModel}\PY{p}{(}\PY{n}{nn}\PY{o}{.}\PY{n}{Module}\PY{p}{)}\PY{p}{:}
    \PY{k}{def}\PY{+w}{ }\PY{n+nf+fm}{\PYZus{}\PYZus{}init\PYZus{}\PYZus{}}\PY{p}{(}\PY{n+nb+bp}{self}\PY{p}{,} \PY{n}{input\PYZus{}size}\PY{p}{)}\PY{p}{:}
        \PY{n+nb}{super}\PY{p}{(}\PY{n}{LSTMModel}\PY{p}{,} \PY{n+nb+bp}{self}\PY{p}{)}\PY{o}{.}\PY{n+nf+fm}{\PYZus{}\PYZus{}init\PYZus{}\PYZus{}}\PY{p}{(}\PY{p}{)}
        \PY{n+nb+bp}{self}\PY{o}{.}\PY{n}{embedding} \PY{o}{=} \PY{n}{nn}\PY{o}{.}\PY{n}{Embedding}\PY{p}{(}\PY{n}{input\PYZus{}size}\PY{p}{,} \PY{l+m+mi}{32}\PY{p}{)}
        \PY{n+nb+bp}{self}\PY{o}{.}\PY{n}{lstm} \PY{o}{=} \PY{n}{nn}\PY{o}{.}\PY{n}{LSTM}\PY{p}{(}\PY{n}{input\PYZus{}size}\PY{o}{=}\PY{l+m+mi}{32}\PY{p}{,}
                            \PY{n}{hidden\PYZus{}size}\PY{o}{=}\PY{l+m+mi}{32}\PY{p}{,}
                            \PY{n}{batch\PYZus{}first}\PY{o}{=}\PY{k+kc}{True}\PY{p}{)}
        \PY{n+nb+bp}{self}\PY{o}{.}\PY{n}{dense} \PY{o}{=} \PY{n}{nn}\PY{o}{.}\PY{n}{Linear}\PY{p}{(}\PY{l+m+mi}{32}\PY{p}{,} \PY{l+m+mi}{1}\PY{p}{)}
    \PY{k}{def}\PY{+w}{ }\PY{n+nf}{forward}\PY{p}{(}\PY{n+nb+bp}{self}\PY{p}{,} \PY{n}{x}\PY{p}{)}\PY{p}{:}
        \PY{n}{val}\PY{p}{,} \PY{p}{(}\PY{n}{h\PYZus{}n}\PY{p}{,} \PY{n}{c\PYZus{}n}\PY{p}{)} \PY{o}{=} \PY{n+nb+bp}{self}\PY{o}{.}\PY{n}{lstm}\PY{p}{(}\PY{n+nb+bp}{self}\PY{o}{.}\PY{n}{embedding}\PY{p}{(}\PY{n}{x}\PY{p}{)}\PY{p}{)}
        \PY{k}{return} \PY{n}{torch}\PY{o}{.}\PY{n}{flatten}\PY{p}{(}\PY{n+nb+bp}{self}\PY{o}{.}\PY{n}{dense}\PY{p}{(}\PY{n}{val}\PY{p}{[}\PY{p}{:}\PY{p}{,}\PY{o}{\PYZhy{}}\PY{l+m+mi}{1}\PY{p}{]}\PY{p}{)}\PY{p}{)}
\end{Verbatim}
\end{tcolorbox}

    We instantiate and take a look at the summary of the model, using the
first 10 documents in the corpus.

    \begin{tcolorbox}[breakable, size=fbox, boxrule=1pt, pad at break*=1mm,colback=cellbackground, colframe=cellborder]
\prompt{In}{incolor}{88}{\boxspacing}
\begin{Verbatim}[commandchars=\\\{\}]
\PY{n}{lstm\PYZus{}model} \PY{o}{=} \PY{n}{LSTMModel}\PY{p}{(}\PY{n}{X\PYZus{}test}\PY{o}{.}\PY{n}{shape}\PY{p}{[}\PY{o}{\PYZhy{}}\PY{l+m+mi}{1}\PY{p}{]}\PY{p}{)}
\PY{n}{summary}\PY{p}{(}\PY{n}{lstm\PYZus{}model}\PY{p}{,}
        \PY{n}{input\PYZus{}data}\PY{o}{=}\PY{n}{imdb\PYZus{}seq\PYZus{}train}\PY{o}{.}\PY{n}{tensors}\PY{p}{[}\PY{l+m+mi}{0}\PY{p}{]}\PY{p}{[}\PY{p}{:}\PY{l+m+mi}{10}\PY{p}{]}\PY{p}{,}
        \PY{n}{col\PYZus{}names}\PY{o}{=}\PY{p}{[}\PY{l+s+s1}{\PYZsq{}}\PY{l+s+s1}{input\PYZus{}size}\PY{l+s+s1}{\PYZsq{}}\PY{p}{,}
                   \PY{l+s+s1}{\PYZsq{}}\PY{l+s+s1}{output\PYZus{}size}\PY{l+s+s1}{\PYZsq{}}\PY{p}{,}
                   \PY{l+s+s1}{\PYZsq{}}\PY{l+s+s1}{num\PYZus{}params}\PY{l+s+s1}{\PYZsq{}}\PY{p}{]}\PY{p}{)}
\end{Verbatim}
\end{tcolorbox}

            \begin{tcolorbox}[breakable, size=fbox, boxrule=.5pt, pad at break*=1mm, opacityfill=0]
\prompt{Out}{outcolor}{88}{\boxspacing}
\begin{Verbatim}[commandchars=\\\{\}]
================================================================================
===================================
Layer (type:depth-idx)                   Input Shape               Output Shape
Param \#
================================================================================
===================================
LSTMModel                                [10, 500]                 [10]
--
├─Embedding: 1-1                         [10, 500]                 [10, 500, 32]
320,096
├─LSTM: 1-2                              [10, 500, 32]             [10, 500, 32]
8,448
├─Linear: 1-3                            [10, 32]                  [10, 1]
33
================================================================================
===================================
Total params: 328,577
Trainable params: 328,577
Non-trainable params: 0
Total mult-adds (Units.MEGABYTES): 45.44
================================================================================
===================================
Input size (MB): 50.00
Forward/backward pass size (MB): 2.56
Params size (MB): 1.31
Estimated Total Size (MB): 53.87
================================================================================
===================================
\end{Verbatim}
\end{tcolorbox}
        
    The 10,003 is suppressed in the summary, but we see it in the parameter
count, since \(10,003\times 32=320,096\).

    \begin{tcolorbox}[breakable, size=fbox, boxrule=1pt, pad at break*=1mm,colback=cellbackground, colframe=cellborder]
\prompt{In}{incolor}{89}{\boxspacing}
\begin{Verbatim}[commandchars=\\\{\}]
\PY{n}{lstm\PYZus{}module} \PY{o}{=} \PY{n}{SimpleModule}\PY{o}{.}\PY{n}{binary\PYZus{}classification}\PY{p}{(}\PY{n}{lstm\PYZus{}model}\PY{p}{)}
\PY{n}{lstm\PYZus{}logger} \PY{o}{=} \PY{n}{CSVLogger}\PY{p}{(}\PY{l+s+s1}{\PYZsq{}}\PY{l+s+s1}{logs}\PY{l+s+s1}{\PYZsq{}}\PY{p}{,} \PY{n}{name}\PY{o}{=}\PY{l+s+s1}{\PYZsq{}}\PY{l+s+s1}{IMDB\PYZus{}LSTM}\PY{l+s+s1}{\PYZsq{}}\PY{p}{)}
\end{Verbatim}
\end{tcolorbox}

    \begin{tcolorbox}[breakable, size=fbox, boxrule=1pt, pad at break*=1mm,colback=cellbackground, colframe=cellborder]
\prompt{In}{incolor}{90}{\boxspacing}
\begin{Verbatim}[commandchars=\\\{\}]
\PY{n}{lstm\PYZus{}trainer} \PY{o}{=} \PY{n}{Trainer}\PY{p}{(}\PY{n}{deterministic}\PY{o}{=}\PY{k+kc}{True}\PY{p}{,}
                       \PY{n}{max\PYZus{}epochs}\PY{o}{=}\PY{l+m+mi}{20}\PY{p}{,}
                       \PY{n}{logger}\PY{o}{=}\PY{n}{lstm\PYZus{}logger}\PY{p}{,}
                       \PY{n}{enable\PYZus{}progress\PYZus{}bar}\PY{o}{=}\PY{k+kc}{False}\PY{p}{,}
                       \PY{n}{callbacks}\PY{o}{=}\PY{p}{[}\PY{n}{ErrorTracker}\PY{p}{(}\PY{p}{)}\PY{p}{]}\PY{p}{)}
\PY{n}{lstm\PYZus{}trainer}\PY{o}{.}\PY{n}{fit}\PY{p}{(}\PY{n}{lstm\PYZus{}module}\PY{p}{,}
                 \PY{n}{datamodule}\PY{o}{=}\PY{n}{imdb\PYZus{}seq\PYZus{}dm}\PY{p}{)}
\end{Verbatim}
\end{tcolorbox}

    \begin{Verbatim}[commandchars=\\\{\}]
GPU available: True (mps), used: True
TPU available: False, using: 0 TPU cores
IPU available: False, using: 0 IPUs
HPU available: False, using: 0 HPUs

  | Name  | Type              | Params
--------------------------------------------
0 | model | LSTMModel         | 328 K
1 | loss  | BCEWithLogitsLoss | 0
--------------------------------------------
328 K     Trainable params
0         Non-trainable params
328 K     Total params
1.314     Total estimated model params size (MB)
`Trainer.fit` stopped: `max\_epochs=20` reached.
    \end{Verbatim}

    The rest is now similar to other networks we have fit. We track the test
performance as the network is fit, and see that it attains 85\%
accuracy.

    \begin{tcolorbox}[breakable, size=fbox, boxrule=1pt, pad at break*=1mm,colback=cellbackground, colframe=cellborder]
\prompt{In}{incolor}{91}{\boxspacing}
\begin{Verbatim}[commandchars=\\\{\}]
\PY{n}{lstm\PYZus{}trainer}\PY{o}{.}\PY{n}{test}\PY{p}{(}\PY{n}{lstm\PYZus{}module}\PY{p}{,} \PY{n}{datamodule}\PY{o}{=}\PY{n}{imdb\PYZus{}seq\PYZus{}dm}\PY{p}{)}
\end{Verbatim}
\end{tcolorbox}

    \begin{Verbatim}[commandchars=\\\{\}]
────────────────────────────────────────────────────────────────────────────────
────────────────────────────────────────
       Test metric             DataLoader 0
────────────────────────────────────────────────────────────────────────────────
────────────────────────────────────────
      test\_accuracy          0.841759979724884
        test\_loss           0.7341755032539368
────────────────────────────────────────────────────────────────────────────────
────────────────────────────────────────
    \end{Verbatim}

            \begin{tcolorbox}[breakable, size=fbox, boxrule=.5pt, pad at break*=1mm, opacityfill=0]
\prompt{Out}{outcolor}{91}{\boxspacing}
\begin{Verbatim}[commandchars=\\\{\}]
[\{'test\_loss': 0.7341755032539368, 'test\_accuracy': 0.841759979724884\}]
\end{Verbatim}
\end{tcolorbox}
        
    We once again show the learning progress, followed by cleanup.

    \begin{tcolorbox}[breakable, size=fbox, boxrule=1pt, pad at break*=1mm,colback=cellbackground, colframe=cellborder]
\prompt{In}{incolor}{92}{\boxspacing}
\begin{Verbatim}[commandchars=\\\{\}]
\PY{n}{lstm\PYZus{}results} \PY{o}{=} \PY{n}{pd}\PY{o}{.}\PY{n}{read\PYZus{}csv}\PY{p}{(}\PY{n}{lstm\PYZus{}logger}\PY{o}{.}\PY{n}{experiment}\PY{o}{.}\PY{n}{metrics\PYZus{}file\PYZus{}path}\PY{p}{)}
\PY{n}{fig}\PY{p}{,} \PY{n}{ax} \PY{o}{=} \PY{n}{subplots}\PY{p}{(}\PY{l+m+mi}{1}\PY{p}{,} \PY{l+m+mi}{1}\PY{p}{,} \PY{n}{figsize}\PY{o}{=}\PY{p}{(}\PY{l+m+mi}{6}\PY{p}{,} \PY{l+m+mi}{6}\PY{p}{)}\PY{p}{)}
\PY{n}{summary\PYZus{}plot}\PY{p}{(}\PY{n}{lstm\PYZus{}results}\PY{p}{,}
             \PY{n}{ax}\PY{p}{,}
             \PY{n}{col}\PY{o}{=}\PY{l+s+s1}{\PYZsq{}}\PY{l+s+s1}{accuracy}\PY{l+s+s1}{\PYZsq{}}\PY{p}{,}
             \PY{n}{ylabel}\PY{o}{=}\PY{l+s+s1}{\PYZsq{}}\PY{l+s+s1}{Accuracy}\PY{l+s+s1}{\PYZsq{}}\PY{p}{)}
\PY{n}{ax}\PY{o}{.}\PY{n}{set\PYZus{}xticks}\PY{p}{(}\PY{n}{np}\PY{o}{.}\PY{n}{linspace}\PY{p}{(}\PY{l+m+mi}{0}\PY{p}{,} \PY{l+m+mi}{20}\PY{p}{,} \PY{l+m+mi}{5}\PY{p}{)}\PY{o}{.}\PY{n}{astype}\PY{p}{(}\PY{n+nb}{int}\PY{p}{)}\PY{p}{)}
\PY{n}{ax}\PY{o}{.}\PY{n}{set\PYZus{}ylabel}\PY{p}{(}\PY{l+s+s1}{\PYZsq{}}\PY{l+s+s1}{Accuracy}\PY{l+s+s1}{\PYZsq{}}\PY{p}{)}
\PY{n}{ax}\PY{o}{.}\PY{n}{set\PYZus{}ylim}\PY{p}{(}\PY{p}{[}\PY{l+m+mf}{0.5}\PY{p}{,} \PY{l+m+mi}{1}\PY{p}{]}\PY{p}{)}
\end{Verbatim}
\end{tcolorbox}

            \begin{tcolorbox}[breakable, size=fbox, boxrule=.5pt, pad at break*=1mm, opacityfill=0]
\prompt{Out}{outcolor}{92}{\boxspacing}
\begin{Verbatim}[commandchars=\\\{\}]
(0.5, 1.0)
\end{Verbatim}
\end{tcolorbox}
        
    \begin{center}
    \adjustimage{max size={0.9\linewidth}{0.9\paperheight}}{artifacts/latex/Ch10-deeplearning-lab_files/artifacts/latex/Ch10-deeplearning-lab_187_1.png}
    \end{center}
    { \hspace*{\fill} \\}
    
    \begin{tcolorbox}[breakable, size=fbox, boxrule=1pt, pad at break*=1mm,colback=cellbackground, colframe=cellborder]
\prompt{In}{incolor}{93}{\boxspacing}
\begin{Verbatim}[commandchars=\\\{\}]
\PY{k}{del}\PY{p}{(}\PY{n}{lstm\PYZus{}model}\PY{p}{,}
    \PY{n}{lstm\PYZus{}trainer}\PY{p}{,}
    \PY{n}{lstm\PYZus{}logger}\PY{p}{,}
    \PY{n}{imdb\PYZus{}seq\PYZus{}dm}\PY{p}{,}
    \PY{n}{imdb\PYZus{}seq\PYZus{}train}\PY{p}{,}
    \PY{n}{imdb\PYZus{}seq\PYZus{}test}\PY{p}{)}
\end{Verbatim}
\end{tcolorbox}

    \hypertarget{time-series-prediction}{%
\subsubsection{Time Series Prediction}\label{time-series-prediction}}

We now show how to fit the models in Section 10.5.2 for time series
prediction. We first load and standardize the data.

    \begin{tcolorbox}[breakable, size=fbox, boxrule=1pt, pad at break*=1mm,colback=cellbackground, colframe=cellborder]
\prompt{In}{incolor}{94}{\boxspacing}
\begin{Verbatim}[commandchars=\\\{\}]
\PY{n}{NYSE} \PY{o}{=} \PY{n}{load\PYZus{}data}\PY{p}{(}\PY{l+s+s1}{\PYZsq{}}\PY{l+s+s1}{NYSE}\PY{l+s+s1}{\PYZsq{}}\PY{p}{)}
\PY{n}{cols} \PY{o}{=} \PY{p}{[}\PY{l+s+s1}{\PYZsq{}}\PY{l+s+s1}{DJ\PYZus{}return}\PY{l+s+s1}{\PYZsq{}}\PY{p}{,} \PY{l+s+s1}{\PYZsq{}}\PY{l+s+s1}{log\PYZus{}volume}\PY{l+s+s1}{\PYZsq{}}\PY{p}{,} \PY{l+s+s1}{\PYZsq{}}\PY{l+s+s1}{log\PYZus{}volatility}\PY{l+s+s1}{\PYZsq{}}\PY{p}{]}
\PY{n}{X} \PY{o}{=} \PY{n}{pd}\PY{o}{.}\PY{n}{DataFrame}\PY{p}{(}\PY{n}{StandardScaler}\PY{p}{(}
                     \PY{n}{with\PYZus{}mean}\PY{o}{=}\PY{k+kc}{True}\PY{p}{,}
                     \PY{n}{with\PYZus{}std}\PY{o}{=}\PY{k+kc}{True}\PY{p}{)}\PY{o}{.}\PY{n}{fit\PYZus{}transform}\PY{p}{(}\PY{n}{NYSE}\PY{p}{[}\PY{n}{cols}\PY{p}{]}\PY{p}{)}\PY{p}{,}
                 \PY{n}{columns}\PY{o}{=}\PY{n}{NYSE}\PY{p}{[}\PY{n}{cols}\PY{p}{]}\PY{o}{.}\PY{n}{columns}\PY{p}{,}
                 \PY{n}{index}\PY{o}{=}\PY{n}{NYSE}\PY{o}{.}\PY{n}{index}\PY{p}{)}
\end{Verbatim}
\end{tcolorbox}

    Next we set up the lagged versions of the data, dropping any rows with
missing values using the \texttt{dropna()} method.

    \begin{tcolorbox}[breakable, size=fbox, boxrule=1pt, pad at break*=1mm,colback=cellbackground, colframe=cellborder]
\prompt{In}{incolor}{95}{\boxspacing}
\begin{Verbatim}[commandchars=\\\{\}]
\PY{k}{for} \PY{n}{lag} \PY{o+ow}{in} \PY{n+nb}{range}\PY{p}{(}\PY{l+m+mi}{1}\PY{p}{,} \PY{l+m+mi}{6}\PY{p}{)}\PY{p}{:}
    \PY{k}{for} \PY{n}{col} \PY{o+ow}{in} \PY{n}{cols}\PY{p}{:}
        \PY{n}{newcol} \PY{o}{=} \PY{n}{np}\PY{o}{.}\PY{n}{zeros}\PY{p}{(}\PY{n}{X}\PY{o}{.}\PY{n}{shape}\PY{p}{[}\PY{l+m+mi}{0}\PY{p}{]}\PY{p}{)} \PY{o}{*} \PY{n}{np}\PY{o}{.}\PY{n}{nan}
        \PY{n}{newcol}\PY{p}{[}\PY{n}{lag}\PY{p}{:}\PY{p}{]} \PY{o}{=} \PY{n}{X}\PY{p}{[}\PY{n}{col}\PY{p}{]}\PY{o}{.}\PY{n}{values}\PY{p}{[}\PY{p}{:}\PY{o}{\PYZhy{}}\PY{n}{lag}\PY{p}{]}
        \PY{n}{X}\PY{o}{.}\PY{n}{insert}\PY{p}{(}\PY{n+nb}{len}\PY{p}{(}\PY{n}{X}\PY{o}{.}\PY{n}{columns}\PY{p}{)}\PY{p}{,} \PY{l+s+s2}{\PYZdq{}}\PY{l+s+si}{\PYZob{}0\PYZcb{}}\PY{l+s+s2}{\PYZus{}}\PY{l+s+si}{\PYZob{}1\PYZcb{}}\PY{l+s+s2}{\PYZdq{}}\PY{o}{.}\PY{n}{format}\PY{p}{(}\PY{n}{col}\PY{p}{,} \PY{n}{lag}\PY{p}{)}\PY{p}{,} \PY{n}{newcol}\PY{p}{)}
\PY{n}{X}\PY{o}{.}\PY{n}{insert}\PY{p}{(}\PY{n+nb}{len}\PY{p}{(}\PY{n}{X}\PY{o}{.}\PY{n}{columns}\PY{p}{)}\PY{p}{,} \PY{l+s+s1}{\PYZsq{}}\PY{l+s+s1}{train}\PY{l+s+s1}{\PYZsq{}}\PY{p}{,} \PY{n}{NYSE}\PY{p}{[}\PY{l+s+s1}{\PYZsq{}}\PY{l+s+s1}{train}\PY{l+s+s1}{\PYZsq{}}\PY{p}{]}\PY{p}{)}
\PY{n}{X} \PY{o}{=} \PY{n}{X}\PY{o}{.}\PY{n}{dropna}\PY{p}{(}\PY{p}{)}
\end{Verbatim}
\end{tcolorbox}

    Finally, we extract the response, training indicator, and drop the
current day's \texttt{DJ\_return} and \texttt{log\_volatility} to
predict only from previous day's data.

    \begin{tcolorbox}[breakable, size=fbox, boxrule=1pt, pad at break*=1mm,colback=cellbackground, colframe=cellborder]
\prompt{In}{incolor}{96}{\boxspacing}
\begin{Verbatim}[commandchars=\\\{\}]
\PY{n}{Y}\PY{p}{,} \PY{n}{train} \PY{o}{=} \PY{n}{X}\PY{p}{[}\PY{l+s+s1}{\PYZsq{}}\PY{l+s+s1}{log\PYZus{}volume}\PY{l+s+s1}{\PYZsq{}}\PY{p}{]}\PY{p}{,} \PY{n}{X}\PY{p}{[}\PY{l+s+s1}{\PYZsq{}}\PY{l+s+s1}{train}\PY{l+s+s1}{\PYZsq{}}\PY{p}{]}
\PY{n}{X} \PY{o}{=} \PY{n}{X}\PY{o}{.}\PY{n}{drop}\PY{p}{(}\PY{n}{columns}\PY{o}{=}\PY{p}{[}\PY{l+s+s1}{\PYZsq{}}\PY{l+s+s1}{train}\PY{l+s+s1}{\PYZsq{}}\PY{p}{]} \PY{o}{+} \PY{n}{cols}\PY{p}{)}
\PY{n}{X}\PY{o}{.}\PY{n}{columns}
\end{Verbatim}
\end{tcolorbox}

            \begin{tcolorbox}[breakable, size=fbox, boxrule=.5pt, pad at break*=1mm, opacityfill=0]
\prompt{Out}{outcolor}{96}{\boxspacing}
\begin{Verbatim}[commandchars=\\\{\}]
Index(['DJ\_return\_1', 'log\_volume\_1', 'log\_volatility\_1', 'DJ\_return\_2',
       'log\_volume\_2', 'log\_volatility\_2', 'DJ\_return\_3', 'log\_volume\_3',
       'log\_volatility\_3', 'DJ\_return\_4', 'log\_volume\_4', 'log\_volatility\_4',
       'DJ\_return\_5', 'log\_volume\_5', 'log\_volatility\_5'],
      dtype='object')
\end{Verbatim}
\end{tcolorbox}
        
    We first fit a simple linear model and compute the \(R^2\) on the test
data using the \texttt{score()} method.

    \begin{tcolorbox}[breakable, size=fbox, boxrule=1pt, pad at break*=1mm,colback=cellbackground, colframe=cellborder]
\prompt{In}{incolor}{97}{\boxspacing}
\begin{Verbatim}[commandchars=\\\{\}]
\PY{n}{M} \PY{o}{=} \PY{n}{LinearRegression}\PY{p}{(}\PY{p}{)}
\PY{n}{M}\PY{o}{.}\PY{n}{fit}\PY{p}{(}\PY{n}{X}\PY{p}{[}\PY{n}{train}\PY{p}{]}\PY{p}{,} \PY{n}{Y}\PY{p}{[}\PY{n}{train}\PY{p}{]}\PY{p}{)}
\PY{n}{M}\PY{o}{.}\PY{n}{score}\PY{p}{(}\PY{n}{X}\PY{p}{[}\PY{o}{\PYZti{}}\PY{n}{train}\PY{p}{]}\PY{p}{,} \PY{n}{Y}\PY{p}{[}\PY{o}{\PYZti{}}\PY{n}{train}\PY{p}{]}\PY{p}{)}
\end{Verbatim}
\end{tcolorbox}

            \begin{tcolorbox}[breakable, size=fbox, boxrule=.5pt, pad at break*=1mm, opacityfill=0]
\prompt{Out}{outcolor}{97}{\boxspacing}
\begin{Verbatim}[commandchars=\\\{\}]
0.4128912938562521
\end{Verbatim}
\end{tcolorbox}
        
    We refit this model, including the factor variable
\texttt{day\_of\_week}. For a categorical series in \texttt{pandas}, we
can form the indicators using the \texttt{get\_dummies()} method.

    \begin{tcolorbox}[breakable, size=fbox, boxrule=1pt, pad at break*=1mm,colback=cellbackground, colframe=cellborder]
\prompt{In}{incolor}{98}{\boxspacing}
\begin{Verbatim}[commandchars=\\\{\}]
\PY{n}{X\PYZus{}day} \PY{o}{=} \PY{n}{pd}\PY{o}{.}\PY{n}{concat}\PY{p}{(}\PY{p}{[}\PY{n}{X}\PY{p}{,} 
                  \PY{n}{pd}\PY{o}{.}\PY{n}{get\PYZus{}dummies}\PY{p}{(}\PY{n}{NYSE}\PY{p}{[}\PY{l+s+s1}{\PYZsq{}}\PY{l+s+s1}{day\PYZus{}of\PYZus{}week}\PY{l+s+s1}{\PYZsq{}}\PY{p}{]}\PY{p}{)}\PY{p}{]}\PY{p}{,}
                  \PY{n}{axis}\PY{o}{=}\PY{l+m+mi}{1}\PY{p}{)}\PY{o}{.}\PY{n}{dropna}\PY{p}{(}\PY{p}{)}
\end{Verbatim}
\end{tcolorbox}

    Note that we do not have to reinstantiate the linear regression model as
its \texttt{fit()} method accepts a design matrix and a response
directly.

    \begin{tcolorbox}[breakable, size=fbox, boxrule=1pt, pad at break*=1mm,colback=cellbackground, colframe=cellborder]
\prompt{In}{incolor}{99}{\boxspacing}
\begin{Verbatim}[commandchars=\\\{\}]
\PY{n}{M}\PY{o}{.}\PY{n}{fit}\PY{p}{(}\PY{n}{X\PYZus{}day}\PY{p}{[}\PY{n}{train}\PY{p}{]}\PY{p}{,} \PY{n}{Y}\PY{p}{[}\PY{n}{train}\PY{p}{]}\PY{p}{)}
\PY{n}{M}\PY{o}{.}\PY{n}{score}\PY{p}{(}\PY{n}{X\PYZus{}day}\PY{p}{[}\PY{o}{\PYZti{}}\PY{n}{train}\PY{p}{]}\PY{p}{,} \PY{n}{Y}\PY{p}{[}\PY{o}{\PYZti{}}\PY{n}{train}\PY{p}{]}\PY{p}{)}
\end{Verbatim}
\end{tcolorbox}

            \begin{tcolorbox}[breakable, size=fbox, boxrule=.5pt, pad at break*=1mm, opacityfill=0]
\prompt{Out}{outcolor}{99}{\boxspacing}
\begin{Verbatim}[commandchars=\\\{\}]
0.4595563133053273
\end{Verbatim}
\end{tcolorbox}
        
    This model achieves an \(R^2\) of about 46\%.

    To fit the RNN, we must reshape the data, as it will expect 5 lagged
versions of each feature as indicated by the \texttt{input\_shape}
argument to the layer \texttt{nn.RNN()} below. We first ensure the
columns of our data frame are such that a reshaped matrix will have the
variables correctly lagged. We use the \texttt{reindex()} method to do
this.

For an input shape \texttt{(5,3)}, each row represents a lagged version
of the three variables. The \texttt{nn.RNN()} layer also expects the
first row of each observation to be earliest in time, so we must reverse
the current order. Hence we loop over \texttt{range(5,0,-1)} below,
which is an example of using a \texttt{slice()} to index iterable
objects. The general notation is \texttt{start:end:step}.

    \begin{tcolorbox}[breakable, size=fbox, boxrule=1pt, pad at break*=1mm,colback=cellbackground, colframe=cellborder]
\prompt{In}{incolor}{100}{\boxspacing}
\begin{Verbatim}[commandchars=\\\{\}]
\PY{n}{ordered\PYZus{}cols} \PY{o}{=} \PY{p}{[}\PY{p}{]}
\PY{k}{for} \PY{n}{lag} \PY{o+ow}{in} \PY{n+nb}{range}\PY{p}{(}\PY{l+m+mi}{5}\PY{p}{,}\PY{l+m+mi}{0}\PY{p}{,}\PY{o}{\PYZhy{}}\PY{l+m+mi}{1}\PY{p}{)}\PY{p}{:}
    \PY{k}{for} \PY{n}{col} \PY{o+ow}{in} \PY{n}{cols}\PY{p}{:}
        \PY{n}{ordered\PYZus{}cols}\PY{o}{.}\PY{n}{append}\PY{p}{(}\PY{l+s+s1}{\PYZsq{}}\PY{l+s+si}{\PYZob{}0\PYZcb{}}\PY{l+s+s1}{\PYZus{}}\PY{l+s+si}{\PYZob{}1\PYZcb{}}\PY{l+s+s1}{\PYZsq{}}\PY{o}{.}\PY{n}{format}\PY{p}{(}\PY{n}{col}\PY{p}{,} \PY{n}{lag}\PY{p}{)}\PY{p}{)}
\PY{n}{X} \PY{o}{=} \PY{n}{X}\PY{o}{.}\PY{n}{reindex}\PY{p}{(}\PY{n}{columns}\PY{o}{=}\PY{n}{ordered\PYZus{}cols}\PY{p}{)}
\PY{n}{X}\PY{o}{.}\PY{n}{columns}
\end{Verbatim}
\end{tcolorbox}

            \begin{tcolorbox}[breakable, size=fbox, boxrule=.5pt, pad at break*=1mm, opacityfill=0]
\prompt{Out}{outcolor}{100}{\boxspacing}
\begin{Verbatim}[commandchars=\\\{\}]
Index(['DJ\_return\_5', 'log\_volume\_5', 'log\_volatility\_5', 'DJ\_return\_4',
       'log\_volume\_4', 'log\_volatility\_4', 'DJ\_return\_3', 'log\_volume\_3',
       'log\_volatility\_3', 'DJ\_return\_2', 'log\_volume\_2', 'log\_volatility\_2',
       'DJ\_return\_1', 'log\_volume\_1', 'log\_volatility\_1'],
      dtype='object')
\end{Verbatim}
\end{tcolorbox}
        
    We now reshape the data.

    \begin{tcolorbox}[breakable, size=fbox, boxrule=1pt, pad at break*=1mm,colback=cellbackground, colframe=cellborder]
\prompt{In}{incolor}{101}{\boxspacing}
\begin{Verbatim}[commandchars=\\\{\}]
\PY{n}{X\PYZus{}rnn} \PY{o}{=} \PY{n}{X}\PY{o}{.}\PY{n}{to\PYZus{}numpy}\PY{p}{(}\PY{p}{)}\PY{o}{.}\PY{n}{reshape}\PY{p}{(}\PY{p}{(}\PY{o}{\PYZhy{}}\PY{l+m+mi}{1}\PY{p}{,}\PY{l+m+mi}{5}\PY{p}{,}\PY{l+m+mi}{3}\PY{p}{)}\PY{p}{)}
\PY{n}{X\PYZus{}rnn}\PY{o}{.}\PY{n}{shape}
\end{Verbatim}
\end{tcolorbox}

            \begin{tcolorbox}[breakable, size=fbox, boxrule=.5pt, pad at break*=1mm, opacityfill=0]
\prompt{Out}{outcolor}{101}{\boxspacing}
\begin{Verbatim}[commandchars=\\\{\}]
(6046, 5, 3)
\end{Verbatim}
\end{tcolorbox}
        
    By specifying the first size as -1, \texttt{numpy.reshape()} deduces its
size based on the remaining arguments.

Now we are ready to proceed with the RNN, which uses 12 hidden units,
and 10\% dropout. After passing through the RNN, we extract the final
time point as \texttt{val{[}:,-1{]}} in \texttt{forward()} below. This
gets passed through a 10\% dropout and then flattened through a linear
layer.

    \begin{tcolorbox}[breakable, size=fbox, boxrule=1pt, pad at break*=1mm,colback=cellbackground, colframe=cellborder]
\prompt{In}{incolor}{102}{\boxspacing}
\begin{Verbatim}[commandchars=\\\{\}]
\PY{k}{class}\PY{+w}{ }\PY{n+nc}{NYSEModel}\PY{p}{(}\PY{n}{nn}\PY{o}{.}\PY{n}{Module}\PY{p}{)}\PY{p}{:}
    \PY{k}{def}\PY{+w}{ }\PY{n+nf+fm}{\PYZus{}\PYZus{}init\PYZus{}\PYZus{}}\PY{p}{(}\PY{n+nb+bp}{self}\PY{p}{)}\PY{p}{:}
        \PY{n+nb}{super}\PY{p}{(}\PY{n}{NYSEModel}\PY{p}{,} \PY{n+nb+bp}{self}\PY{p}{)}\PY{o}{.}\PY{n+nf+fm}{\PYZus{}\PYZus{}init\PYZus{}\PYZus{}}\PY{p}{(}\PY{p}{)}
        \PY{n+nb+bp}{self}\PY{o}{.}\PY{n}{rnn} \PY{o}{=} \PY{n}{nn}\PY{o}{.}\PY{n}{RNN}\PY{p}{(}\PY{l+m+mi}{3}\PY{p}{,}
                          \PY{l+m+mi}{12}\PY{p}{,}
                          \PY{n}{batch\PYZus{}first}\PY{o}{=}\PY{k+kc}{True}\PY{p}{)}
        \PY{n+nb+bp}{self}\PY{o}{.}\PY{n}{dense} \PY{o}{=} \PY{n}{nn}\PY{o}{.}\PY{n}{Linear}\PY{p}{(}\PY{l+m+mi}{12}\PY{p}{,} \PY{l+m+mi}{1}\PY{p}{)}
        \PY{n+nb+bp}{self}\PY{o}{.}\PY{n}{dropout} \PY{o}{=} \PY{n}{nn}\PY{o}{.}\PY{n}{Dropout}\PY{p}{(}\PY{l+m+mf}{0.1}\PY{p}{)}
    \PY{k}{def}\PY{+w}{ }\PY{n+nf}{forward}\PY{p}{(}\PY{n+nb+bp}{self}\PY{p}{,} \PY{n}{x}\PY{p}{)}\PY{p}{:}
        \PY{n}{val}\PY{p}{,} \PY{n}{h\PYZus{}n} \PY{o}{=} \PY{n+nb+bp}{self}\PY{o}{.}\PY{n}{rnn}\PY{p}{(}\PY{n}{x}\PY{p}{)}
        \PY{n}{val} \PY{o}{=} \PY{n+nb+bp}{self}\PY{o}{.}\PY{n}{dense}\PY{p}{(}\PY{n+nb+bp}{self}\PY{o}{.}\PY{n}{dropout}\PY{p}{(}\PY{n}{val}\PY{p}{[}\PY{p}{:}\PY{p}{,}\PY{o}{\PYZhy{}}\PY{l+m+mi}{1}\PY{p}{]}\PY{p}{)}\PY{p}{)}
        \PY{k}{return} \PY{n}{torch}\PY{o}{.}\PY{n}{flatten}\PY{p}{(}\PY{n}{val}\PY{p}{)}
\PY{n}{nyse\PYZus{}model} \PY{o}{=} \PY{n}{NYSEModel}\PY{p}{(}\PY{p}{)}
\end{Verbatim}
\end{tcolorbox}

    We fit the model in a similar fashion to previous networks. We supply
the \texttt{fit} function with test data as validation data, so that
when we monitor its progress and plot the history function we can see
the progress on the test data. Of course we should not use this as a
basis for early stopping, since then the test performance would be
biased.

We form the training dataset similar to our \texttt{Hitters} example.

    \begin{tcolorbox}[breakable, size=fbox, boxrule=1pt, pad at break*=1mm,colback=cellbackground, colframe=cellborder]
\prompt{In}{incolor}{103}{\boxspacing}
\begin{Verbatim}[commandchars=\\\{\}]
\PY{n}{datasets} \PY{o}{=} \PY{p}{[}\PY{p}{]}
\PY{k}{for} \PY{n}{mask} \PY{o+ow}{in} \PY{p}{[}\PY{n}{train}\PY{p}{,} \PY{o}{\PYZti{}}\PY{n}{train}\PY{p}{]}\PY{p}{:}
    \PY{n}{X\PYZus{}rnn\PYZus{}t} \PY{o}{=} \PY{n}{torch}\PY{o}{.}\PY{n}{tensor}\PY{p}{(}\PY{n}{X\PYZus{}rnn}\PY{p}{[}\PY{n}{mask}\PY{p}{]}\PY{o}{.}\PY{n}{astype}\PY{p}{(}\PY{n}{np}\PY{o}{.}\PY{n}{float32}\PY{p}{)}\PY{p}{)}
    \PY{n}{Y\PYZus{}t} \PY{o}{=} \PY{n}{torch}\PY{o}{.}\PY{n}{tensor}\PY{p}{(}\PY{n}{Y}\PY{p}{[}\PY{n}{mask}\PY{p}{]}\PY{o}{.}\PY{n}{astype}\PY{p}{(}\PY{n}{np}\PY{o}{.}\PY{n}{float32}\PY{p}{)}\PY{p}{)}
    \PY{n}{datasets}\PY{o}{.}\PY{n}{append}\PY{p}{(}\PY{n}{TensorDataset}\PY{p}{(}\PY{n}{X\PYZus{}rnn\PYZus{}t}\PY{p}{,} \PY{n}{Y\PYZus{}t}\PY{p}{)}\PY{p}{)}
\PY{n}{nyse\PYZus{}train}\PY{p}{,} \PY{n}{nyse\PYZus{}test} \PY{o}{=} \PY{n}{datasets}
\end{Verbatim}
\end{tcolorbox}

    \begin{Verbatim}[commandchars=\\\{\}]
/var/folders/16/8y65\_zv174qgdp4ktlmpv12h0000gq/T/ipykernel\_8578/3890359531.py:4:
FutureWarning: Series.\_\_getitem\_\_ treating keys as positions is deprecated. In a
future version, integer keys will always be treated as labels (consistent with
DataFrame behavior). To access a value by position, use `ser.iloc[pos]`
  Y\_t = torch.tensor(Y[mask].astype(np.float32))
/var/folders/16/8y65\_zv174qgdp4ktlmpv12h0000gq/T/ipykernel\_8578/3890359531.py:4:
FutureWarning: Series.\_\_getitem\_\_ treating keys as positions is deprecated. In a
future version, integer keys will always be treated as labels (consistent with
DataFrame behavior). To access a value by position, use `ser.iloc[pos]`
  Y\_t = torch.tensor(Y[mask].astype(np.float32))
    \end{Verbatim}

    Following our usual pattern, we inspect the summary.

    \begin{tcolorbox}[breakable, size=fbox, boxrule=1pt, pad at break*=1mm,colback=cellbackground, colframe=cellborder]
\prompt{In}{incolor}{104}{\boxspacing}
\begin{Verbatim}[commandchars=\\\{\}]
\PY{n}{summary}\PY{p}{(}\PY{n}{nyse\PYZus{}model}\PY{p}{,}
        \PY{n}{input\PYZus{}data}\PY{o}{=}\PY{n}{X\PYZus{}rnn\PYZus{}t}\PY{p}{,}
        \PY{n}{col\PYZus{}names}\PY{o}{=}\PY{p}{[}\PY{l+s+s1}{\PYZsq{}}\PY{l+s+s1}{input\PYZus{}size}\PY{l+s+s1}{\PYZsq{}}\PY{p}{,}
                   \PY{l+s+s1}{\PYZsq{}}\PY{l+s+s1}{output\PYZus{}size}\PY{l+s+s1}{\PYZsq{}}\PY{p}{,}
                   \PY{l+s+s1}{\PYZsq{}}\PY{l+s+s1}{num\PYZus{}params}\PY{l+s+s1}{\PYZsq{}}\PY{p}{]}\PY{p}{)}
\end{Verbatim}
\end{tcolorbox}

            \begin{tcolorbox}[breakable, size=fbox, boxrule=.5pt, pad at break*=1mm, opacityfill=0]
\prompt{Out}{outcolor}{104}{\boxspacing}
\begin{Verbatim}[commandchars=\\\{\}]
================================================================================
===================================
Layer (type:depth-idx)                   Input Shape               Output Shape
Param \#
================================================================================
===================================
NYSEModel                                [1770, 5, 3]              [1770]
--
├─RNN: 1-1                               [1770, 5, 3]              [1770, 5, 12]
204
├─Dropout: 1-2                           [1770, 12]                [1770, 12]
--
├─Linear: 1-3                            [1770, 12]                [1770, 1]
13
================================================================================
===================================
Total params: 217
Trainable params: 217
Non-trainable params: 0
Total mult-adds (Units.MEGABYTES): 1.83
================================================================================
===================================
Input size (MB): 0.11
Forward/backward pass size (MB): 0.86
Params size (MB): 0.00
Estimated Total Size (MB): 0.97
================================================================================
===================================
\end{Verbatim}
\end{tcolorbox}
        
    We again put the two datasets into a data module, with a batch size of
64.

    \begin{tcolorbox}[breakable, size=fbox, boxrule=1pt, pad at break*=1mm,colback=cellbackground, colframe=cellborder]
\prompt{In}{incolor}{105}{\boxspacing}
\begin{Verbatim}[commandchars=\\\{\}]
\PY{n}{nyse\PYZus{}dm} \PY{o}{=} \PY{n}{SimpleDataModule}\PY{p}{(}\PY{n}{nyse\PYZus{}train}\PY{p}{,}
                           \PY{n}{nyse\PYZus{}test}\PY{p}{,}
                           \PY{n}{num\PYZus{}workers}\PY{o}{=}\PY{n+nb}{min}\PY{p}{(}\PY{l+m+mi}{4}\PY{p}{,} \PY{n}{max\PYZus{}num\PYZus{}workers}\PY{p}{)}\PY{p}{,}
                           \PY{n}{validation}\PY{o}{=}\PY{n}{nyse\PYZus{}test}\PY{p}{,}
                           \PY{n}{batch\PYZus{}size}\PY{o}{=}\PY{l+m+mi}{64}\PY{p}{)}
\end{Verbatim}
\end{tcolorbox}

    We run some data through our model to be sure the sizes match up
correctly.

    \begin{tcolorbox}[breakable, size=fbox, boxrule=1pt, pad at break*=1mm,colback=cellbackground, colframe=cellborder]
\prompt{In}{incolor}{106}{\boxspacing}
\begin{Verbatim}[commandchars=\\\{\}]
\PY{k}{for} \PY{n}{idx}\PY{p}{,} \PY{p}{(}\PY{n}{x}\PY{p}{,} \PY{n}{y}\PY{p}{)} \PY{o+ow}{in} \PY{n+nb}{enumerate}\PY{p}{(}\PY{n}{nyse\PYZus{}dm}\PY{o}{.}\PY{n}{train\PYZus{}dataloader}\PY{p}{(}\PY{p}{)}\PY{p}{)}\PY{p}{:}
    \PY{n}{out} \PY{o}{=} \PY{n}{nyse\PYZus{}model}\PY{p}{(}\PY{n}{x}\PY{p}{)}
    \PY{n+nb}{print}\PY{p}{(}\PY{n}{y}\PY{o}{.}\PY{n}{size}\PY{p}{(}\PY{p}{)}\PY{p}{,} \PY{n}{out}\PY{o}{.}\PY{n}{size}\PY{p}{(}\PY{p}{)}\PY{p}{)}
    \PY{k}{if} \PY{n}{idx} \PY{o}{\PYZgt{}}\PY{o}{=} \PY{l+m+mi}{2}\PY{p}{:}
        \PY{k}{break}
\end{Verbatim}
\end{tcolorbox}

    \begin{Verbatim}[commandchars=\\\{\}]
torch.Size([64]) torch.Size([64])
torch.Size([64]) torch.Size([64])
torch.Size([64]) torch.Size([64])
    \end{Verbatim}

    We follow our previous example for setting up a trainer for a regression
problem, requesting the \(R^2\) metric to be be computed at each epoch.

    \begin{tcolorbox}[breakable, size=fbox, boxrule=1pt, pad at break*=1mm,colback=cellbackground, colframe=cellborder]
\prompt{In}{incolor}{107}{\boxspacing}
\begin{Verbatim}[commandchars=\\\{\}]
\PY{n}{nyse\PYZus{}optimizer} \PY{o}{=} \PY{n}{RMSprop}\PY{p}{(}\PY{n}{nyse\PYZus{}model}\PY{o}{.}\PY{n}{parameters}\PY{p}{(}\PY{p}{)}\PY{p}{,}
                         \PY{n}{lr}\PY{o}{=}\PY{l+m+mf}{0.001}\PY{p}{)}
\PY{n}{nyse\PYZus{}module} \PY{o}{=} \PY{n}{SimpleModule}\PY{o}{.}\PY{n}{regression}\PY{p}{(}\PY{n}{nyse\PYZus{}model}\PY{p}{,}
                                      \PY{n}{optimizer}\PY{o}{=}\PY{n}{nyse\PYZus{}optimizer}\PY{p}{,}
                                      \PY{n}{metrics}\PY{o}{=}\PY{p}{\PYZob{}}\PY{l+s+s1}{\PYZsq{}}\PY{l+s+s1}{r2}\PY{l+s+s1}{\PYZsq{}}\PY{p}{:}\PY{n}{R2Score}\PY{p}{(}\PY{p}{)}\PY{p}{\PYZcb{}}\PY{p}{)}
\end{Verbatim}
\end{tcolorbox}

    Fitting the model should by now be familiar. The results on the test
data are very similar to the linear AR model.

    \begin{tcolorbox}[breakable, size=fbox, boxrule=1pt, pad at break*=1mm,colback=cellbackground, colframe=cellborder]
\prompt{In}{incolor}{ }{\boxspacing}
\begin{Verbatim}[commandchars=\\\{\}]
\PY{n}{nyse\PYZus{}trainer} \PY{o}{=} \PY{n}{Trainer}\PY{p}{(}\PY{n}{deterministic}\PY{o}{=}\PY{k+kc}{True}\PY{p}{,}
                       \PY{n}{max\PYZus{}epochs}\PY{o}{=}\PY{l+m+mi}{200}\PY{p}{,}
                       \PY{n}{enable\PYZus{}progress\PYZus{}bar}\PY{o}{=}\PY{k+kc}{False}\PY{p}{,}
                       \PY{n}{callbacks}\PY{o}{=}\PY{p}{[}\PY{n}{ErrorTracker}\PY{p}{(}\PY{p}{)}\PY{p}{]}\PY{p}{)}
\PY{n}{nyse\PYZus{}trainer}\PY{o}{.}\PY{n}{fit}\PY{p}{(}\PY{n}{nyse\PYZus{}module}\PY{p}{,}
                 \PY{n}{datamodule}\PY{o}{=}\PY{n}{nyse\PYZus{}dm}\PY{p}{)}
\PY{n}{nyse\PYZus{}trainer}\PY{o}{.}\PY{n}{test}\PY{p}{(}\PY{n}{nyse\PYZus{}module}\PY{p}{,}
                  \PY{n}{datamodule}\PY{o}{=}\PY{n}{nyse\PYZus{}dm}\PY{p}{)}
\end{Verbatim}
\end{tcolorbox}

    \begin{Verbatim}[commandchars=\\\{\}]
GPU available: True (mps), used: True
TPU available: False, using: 0 TPU cores
IPU available: False, using: 0 IPUs
HPU available: False, using: 0 HPUs

  | Name  | Type      | Params
------------------------------------
0 | model | NYSEModel | 217
1 | loss  | MSELoss   | 0
------------------------------------
217       Trainable params
0         Non-trainable params
217       Total params
0.001     Total estimated model params size (MB)
    \end{Verbatim}

    We could also fit a model without the \texttt{nn.RNN()} layer by just
using a \texttt{nn.Flatten()} layer instead. This would be a nonlinear
AR model. If in addition we excluded the hidden layer, this would be
equivalent to our earlier linear AR model.

Instead we will fit a nonlinear AR model using the feature set
\texttt{X\_day} that includes the \texttt{day\_of\_week} indicators. To
do so, we must first create our test and training datasets and a
corresponding data module. This may seem a little burdensome, but is
part of the general pipeline for \texttt{torch}.

    \begin{tcolorbox}[breakable, size=fbox, boxrule=1pt, pad at break*=1mm,colback=cellbackground, colframe=cellborder]
\prompt{In}{incolor}{ }{\boxspacing}
\begin{Verbatim}[commandchars=\\\{\}]
\PY{n}{datasets} \PY{o}{=} \PY{p}{[}\PY{p}{]}
\PY{k}{for} \PY{n}{mask} \PY{o+ow}{in} \PY{p}{[}\PY{n}{train}\PY{p}{,} \PY{o}{\PYZti{}}\PY{n}{train}\PY{p}{]}\PY{p}{:}
    \PY{n}{X\PYZus{}day\PYZus{}t} \PY{o}{=} \PY{n}{torch}\PY{o}{.}\PY{n}{tensor}\PY{p}{(}
                   \PY{n}{np}\PY{o}{.}\PY{n}{asarray}\PY{p}{(}\PY{n}{X\PYZus{}day}\PY{p}{[}\PY{n}{mask}\PY{p}{]}\PY{p}{)}\PY{o}{.}\PY{n}{astype}\PY{p}{(}\PY{n}{np}\PY{o}{.}\PY{n}{float32}\PY{p}{)}\PY{p}{)}
    \PY{n}{Y\PYZus{}t} \PY{o}{=} \PY{n}{torch}\PY{o}{.}\PY{n}{tensor}\PY{p}{(}\PY{n}{np}\PY{o}{.}\PY{n}{asarray}\PY{p}{(}\PY{n}{Y}\PY{p}{[}\PY{n}{mask}\PY{p}{]}\PY{p}{)}\PY{o}{.}\PY{n}{astype}\PY{p}{(}\PY{n}{np}\PY{o}{.}\PY{n}{float32}\PY{p}{)}\PY{p}{)}
    \PY{n}{datasets}\PY{o}{.}\PY{n}{append}\PY{p}{(}\PY{n}{TensorDataset}\PY{p}{(}\PY{n}{X\PYZus{}day\PYZus{}t}\PY{p}{,} \PY{n}{Y\PYZus{}t}\PY{p}{)}\PY{p}{)}
\PY{n}{day\PYZus{}train}\PY{p}{,} \PY{n}{day\PYZus{}test} \PY{o}{=} \PY{n}{datasets}
\end{Verbatim}
\end{tcolorbox}

    Creating a data module follows a familiar pattern.

    \begin{tcolorbox}[breakable, size=fbox, boxrule=1pt, pad at break*=1mm,colback=cellbackground, colframe=cellborder]
\prompt{In}{incolor}{ }{\boxspacing}
\begin{Verbatim}[commandchars=\\\{\}]
\PY{n}{day\PYZus{}dm} \PY{o}{=} \PY{n}{SimpleDataModule}\PY{p}{(}\PY{n}{day\PYZus{}train}\PY{p}{,}
                          \PY{n}{day\PYZus{}test}\PY{p}{,}
                          \PY{n}{num\PYZus{}workers}\PY{o}{=}\PY{n+nb}{min}\PY{p}{(}\PY{l+m+mi}{4}\PY{p}{,} \PY{n}{max\PYZus{}num\PYZus{}workers}\PY{p}{)}\PY{p}{,}
                          \PY{n}{validation}\PY{o}{=}\PY{n}{day\PYZus{}test}\PY{p}{,}
                          \PY{n}{batch\PYZus{}size}\PY{o}{=}\PY{l+m+mi}{64}\PY{p}{)}
\end{Verbatim}
\end{tcolorbox}

    We build a \texttt{NonLinearARModel()} that takes as input the 20
features and a hidden layer with 32 units. The remaining steps are
familiar.

    \begin{tcolorbox}[breakable, size=fbox, boxrule=1pt, pad at break*=1mm,colback=cellbackground, colframe=cellborder]
\prompt{In}{incolor}{ }{\boxspacing}
\begin{Verbatim}[commandchars=\\\{\}]
\PY{k}{class}\PY{+w}{ }\PY{n+nc}{NonLinearARModel}\PY{p}{(}\PY{n}{nn}\PY{o}{.}\PY{n}{Module}\PY{p}{)}\PY{p}{:}
    \PY{k}{def}\PY{+w}{ }\PY{n+nf+fm}{\PYZus{}\PYZus{}init\PYZus{}\PYZus{}}\PY{p}{(}\PY{n+nb+bp}{self}\PY{p}{)}\PY{p}{:}
        \PY{n+nb}{super}\PY{p}{(}\PY{n}{NonLinearARModel}\PY{p}{,} \PY{n+nb+bp}{self}\PY{p}{)}\PY{o}{.}\PY{n+nf+fm}{\PYZus{}\PYZus{}init\PYZus{}\PYZus{}}\PY{p}{(}\PY{p}{)}
        \PY{n+nb+bp}{self}\PY{o}{.}\PY{n}{\PYZus{}forward} \PY{o}{=} \PY{n}{nn}\PY{o}{.}\PY{n}{Sequential}\PY{p}{(}\PY{n}{nn}\PY{o}{.}\PY{n}{Flatten}\PY{p}{(}\PY{p}{)}\PY{p}{,}
                                      \PY{n}{nn}\PY{o}{.}\PY{n}{Linear}\PY{p}{(}\PY{l+m+mi}{20}\PY{p}{,} \PY{l+m+mi}{32}\PY{p}{)}\PY{p}{,}
                                      \PY{n}{nn}\PY{o}{.}\PY{n}{ReLU}\PY{p}{(}\PY{p}{)}\PY{p}{,}
                                      \PY{n}{nn}\PY{o}{.}\PY{n}{Dropout}\PY{p}{(}\PY{l+m+mf}{0.5}\PY{p}{)}\PY{p}{,}
                                      \PY{n}{nn}\PY{o}{.}\PY{n}{Linear}\PY{p}{(}\PY{l+m+mi}{32}\PY{p}{,} \PY{l+m+mi}{1}\PY{p}{)}\PY{p}{)}
    \PY{k}{def}\PY{+w}{ }\PY{n+nf}{forward}\PY{p}{(}\PY{n+nb+bp}{self}\PY{p}{,} \PY{n}{x}\PY{p}{)}\PY{p}{:}
        \PY{k}{return} \PY{n}{torch}\PY{o}{.}\PY{n}{flatten}\PY{p}{(}\PY{n+nb+bp}{self}\PY{o}{.}\PY{n}{\PYZus{}forward}\PY{p}{(}\PY{n}{x}\PY{p}{)}\PY{p}{)}
\end{Verbatim}
\end{tcolorbox}

    \begin{tcolorbox}[breakable, size=fbox, boxrule=1pt, pad at break*=1mm,colback=cellbackground, colframe=cellborder]
\prompt{In}{incolor}{ }{\boxspacing}
\begin{Verbatim}[commandchars=\\\{\}]
\PY{n}{nl\PYZus{}model} \PY{o}{=} \PY{n}{NonLinearARModel}\PY{p}{(}\PY{p}{)}
\PY{n}{nl\PYZus{}optimizer} \PY{o}{=} \PY{n}{RMSprop}\PY{p}{(}\PY{n}{nl\PYZus{}model}\PY{o}{.}\PY{n}{parameters}\PY{p}{(}\PY{p}{)}\PY{p}{,}
                           \PY{n}{lr}\PY{o}{=}\PY{l+m+mf}{0.001}\PY{p}{)}
\PY{n}{nl\PYZus{}module} \PY{o}{=} \PY{n}{SimpleModule}\PY{o}{.}\PY{n}{regression}\PY{p}{(}\PY{n}{nl\PYZus{}model}\PY{p}{,}
                                        \PY{n}{optimizer}\PY{o}{=}\PY{n}{nl\PYZus{}optimizer}\PY{p}{,}
                                        \PY{n}{metrics}\PY{o}{=}\PY{p}{\PYZob{}}\PY{l+s+s1}{\PYZsq{}}\PY{l+s+s1}{r2}\PY{l+s+s1}{\PYZsq{}}\PY{p}{:}\PY{n}{R2Score}\PY{p}{(}\PY{p}{)}\PY{p}{\PYZcb{}}\PY{p}{)}
\end{Verbatim}
\end{tcolorbox}

    We continue with the usual training steps, fit the model, and evaluate
the test error. We see the test \(R^2\) is a slight improvement over the
linear AR model that also includes \texttt{day\_of\_week}.

    \begin{tcolorbox}[breakable, size=fbox, boxrule=1pt, pad at break*=1mm,colback=cellbackground, colframe=cellborder]
\prompt{In}{incolor}{ }{\boxspacing}
\begin{Verbatim}[commandchars=\\\{\}]
\PY{n}{nl\PYZus{}trainer} \PY{o}{=} \PY{n}{Trainer}\PY{p}{(}\PY{n}{deterministic}\PY{o}{=}\PY{k+kc}{True}\PY{p}{,}
                     \PY{n}{max\PYZus{}epochs}\PY{o}{=}\PY{l+m+mi}{20}\PY{p}{,}
                     \PY{n}{enable\PYZus{}progress\PYZus{}bar}\PY{o}{=}\PY{k+kc}{False}\PY{p}{,}
                     \PY{n}{callbacks}\PY{o}{=}\PY{p}{[}\PY{n}{ErrorTracker}\PY{p}{(}\PY{p}{)}\PY{p}{]}\PY{p}{)}
\PY{n}{nl\PYZus{}trainer}\PY{o}{.}\PY{n}{fit}\PY{p}{(}\PY{n}{nl\PYZus{}module}\PY{p}{,} \PY{n}{datamodule}\PY{o}{=}\PY{n}{day\PYZus{}dm}\PY{p}{)}
\PY{n}{nl\PYZus{}trainer}\PY{o}{.}\PY{n}{test}\PY{p}{(}\PY{n}{nl\PYZus{}module}\PY{p}{,} \PY{n}{datamodule}\PY{o}{=}\PY{n}{day\PYZus{}dm}\PY{p}{)} 
\end{Verbatim}
\end{tcolorbox}

    


    % Add a bibliography block to the postdoc
    
    
    
\end{document}
